% Les maladies : la crise sanitaire et la médecine traditionnelle — Allemand Tle A — Cours signé OnBuch+
% Programme officiel MINESEC. Compiler avec : tectonic AL14-les-maladies-la-crise-sanitaire-et-la-medecine-traditionnell.tex
\newcommand{\DOCMATIERE}{Allemand}
\newcommand{\DOCNIVEAU}{Tle A}
\newcommand{\DOCMODULE}{Chapitre 3 : Environnement, santé et bien-être}
\newcommand{\DOCLECON}{AL14}
\newcommand{\DOCTITRE}{Les maladies : la crise sanitaire et la médecine traditionnelle}
% OnBuch+ — préambule « cours signés » ENRICHI (compilé avec tectonic / XeLaTeX)
% Système visuel « Pop » d'OnBuch+ : fond crème (jamais blanc), bordures encre
% épaisses, ombres « dures » décalées, titres Archivo Black en capitales.
% Version MATHÉMATIQUES : graphes (pgfplots/tikz), boîtes pédagogiques
% (définition, propriété, méthode, attention, le savais-tu, exercice résolu,
% exercice, TP, bilan), page de garde et signature OnBuch+ ancrée en bas.
%
% Macros attendues dans main.tex AVANT \input{preamble} :
%   \DOCMATIERE  \DOCNIVEAU  \DOCMODULE  \DOCLECON (numéro)  \DOCTITRE
\documentclass[11pt]{article}

\usepackage{fontspec}
\usepackage[ngerman,french]{babel} % français = langue principale ; l'allemand via \all{..} / textebox
\usepackage[a4paper,margin=2cm,top=3.4cm,bottom=2.8cm,headheight=1.4cm,footskip=1.3cm]{geometry}
\usepackage{amsmath,amssymb,amsfonts}
\usepackage{mathtools} % \xmapsto, \coloneqq... souvent utilisés par les modèles
\usepackage{cancel} % simplifications barrées (bilans de forces)
% \abs{z} (module, valeur absolue) et \norm{u} (norme), utilisés par les modèles
\providecommand{\abs}{}\let\abs\relax\DeclarePairedDelimiter{\abs}{\lvert}{\rvert}
\providecommand{\norm}{}\let\norm\relax\DeclarePairedDelimiter{\norm}{\lVert}{\rVert}
\usepackage[table]{xcolor} % [table] : fournit aussi \rowcolors (alternance de lignes)
\usepackage{pagecolor}
\usepackage{tikz}
\usetikzlibrary{arrows.meta,positioning,calc,shapes.geometric,shapes.misc,shapes.arrows,decorations.pathmorphing,decorations.pathreplacing,decorations.markings,patterns,shadows,mindmap,trees,fit,backgrounds,matrix,intersections,angles,quotes}
\usepackage{pgfplots}
\usepackage[version=4]{mhchem} % équations nucléaires \ce{^{235}_{92}U}
\usepackage[european,siunitx]{circuitikz} % schémas électriques
\pgfplotsset{compat=1.18}
\usepgfplotslibrary{fillbetween,groupplots} % aires entre courbes (calcul intégral)
\usepackage{enumitem}
\usepackage{fancyhdr}
% [most] charge toutes les bibliothèques tcolorbox (listings, minted, raster,
% documentation...) et gonfle inutilement la mémoire TeX sur un document long
% (~300 boîtes) : on ne charge que ce qui est réellement utilisé.
\usepackage[skins,breakable]{tcolorbox}
\usepackage{graphicx}
\usepackage{colortbl}
\usepackage{booktabs}
\usepackage{tabularx}
\usepackage{diagbox} % case d'en-tête diagonale des tableaux à double entrée
% Colonnes extensibles centrée / gauche / droite (C, L, R), souvent utilisées par les modèles
\newcolumntype{C}{>{\centering\arraybackslash}X}
\newcolumntype{L}{>{\raggedright\arraybackslash}X}
\newcolumntype{R}{>{\raggedleft\arraybackslash}X}
\usepackage{array}
\usepackage{multicol}
\usepackage{multirow}
\usepackage{siunitx}
% parse-numbers=false : accepte aussi \SI{2,5 \times 10^{-3}}{...}
\sisetup{locale=FR,output-decimal-marker={,},group-minimum-digits=4,parse-numbers=false}
\DeclareSIUnit\ohm{\ensuremath{\symup{\Omega}}} % U+2126 absent de la police
\DeclareSIUnit\division{div} % réglages d'oscilloscope (ms/div, V/div)
\DeclareSIUnit\tr{tr} % tours (vitesse de rotation, ex. \SI{1500}{\tr\per\minute})
\DeclareSIUnit\molar{M} % concentration molaire (ex. \SI{3}{\molar}, \SI{1}{\micro\molar})
\DeclareSIUnit\an{an} % année (le modèle confond parfois avec \year, un compteur TeX) — ex. \SI{5}{\kilo\meter\per\an}
\DeclareSIUnit\FCFA{FCFA} % franc CFA (ex. \SI{500}{\FCFA\per\kilo\gram})
\DeclareSIUnit\degreeBrix{\textdegree Brix} % teneur en sucre d'un sirop/jus (réfractométrie)
\DeclareSIUnit\month{mois} % le modèle utilise parfois \month, un compteur TeX, comme unité de durée
\usepackage{qrcode}
% \degree : commande fréquemment utilisée par les modèles pour un angle ou une
% température en dehors de \SI{}{} (non fournie par les paquets chargés).
\providecommand{\degree}{^{\circ}}
% \qed (fin de démonstration), utilisé par les modèles sans amsthm
\providecommand{\qed}{\hfill$\square$}
% \barre : double barre des valeurs interdites dans un tableau de variations (tkz-tab)
\providecommand{\barre}{\ensuremath{\|}}
% environnement proof (amsthm non chargé) utilisé par les modèles
\ifdefined\proof\else\newenvironment{proof}[1][Démonstration]{\par\noindent\textit{#1.}\ }{\qed\par}\fi
\usepackage{lastpage}
\usepackage{hyperref}
\hypersetup{hidelinks}

% ============================================================
% Palette « Pop » OnBuch+ — mêmes hex que la classe P (plus_ui.dart)
% ============================================================
\definecolor{poporange}{HTML}{F59321}
\definecolor{popdark}{HTML}{E07A0C}
\definecolor{popcream}{HTML}{FFF6E8}
\definecolor{popink}{HTML}{1C1714}
\definecolor{poppurple}{HTML}{7A5AE0}
\definecolor{popgreen}{HTML}{2E9E6A}
\definecolor{poppink}{HTML}{E0457B}
\definecolor{popblue}{HTML}{2F7DE1}
\definecolor{popgold}{HTML}{C98A12}
\definecolor{popmuted}{HTML}{8A7F76}
\definecolor{popline}{HTML}{EADFD2}
\definecolor{popgray}{HTML}{8A7F76}
\colorlet{popgrey}{popgray}
% Alias tolérants (le modèle nomme parfois les couleurs différemment)
\colorlet{popwhite}{white}
\colorlet{popblack}{popink}
\colorlet{popred}{poppink}
\colorlet{popyellow}{popgold}
% Teintes claires (fonds de tableaux/encadrés) — le modèle nomme parfois
% les couleurs avec un suffixe L (ex. popblueL) sans les avoir définies.
\colorlet{poporangeL}{poporange!20}
\colorlet{popdarkL}{popdark!20}
\colorlet{poppurpleL}{poppurple!20}
\colorlet{popgreenL}{popgreen!20}
\colorlet{poppinkL}{poppink!20}
\colorlet{popblueL}{popblue!20}
\colorlet{popgoldL}{popgold!20}
\colorlet{popredL}{popred!20}
\colorlet{popgrayL}{popgray!20}
% Teintes claires pour fonds de tableaux / zones
\colorlet{popblueL}{popblue!10!white}
\colorlet{poporangeL}{poporange!14!white}
\colorlet{popgreenL}{popgreen!12!white}
\colorlet{poppurpleL}{poppurple!12!white}
\colorlet{poppinkL}{poppink!12!white}
\colorlet{popgoldL}{popgold!14!white}

\pagecolor{popcream}

% ============================================================
% Typographie
% ============================================================
\defaultfontfeatures{Ligatures=TeX}
\setmainfont{PlusJakartaSans-Medium.ttf}[Path=../../../fonts/,BoldFont=PlusJakartaSans-Bold.ttf]
% Maths en sans-serif (Fira Math) avec chiffres et lettres droites de Plus
% Jakarta, pour que les formules se fondent dans le texte.
\usepackage[math-style=ISO,bold-style=ISO]{unicode-math}
\setmathfont{FiraMath-Regular.otf}[Path=../../../fonts/]
\setmathfont{PlusJakartaSans-Medium.ttf}[Path=../../../fonts/,range={up/num,up/{Latin,latin},\mathup}]
\usepackage{icomma} % virgule décimale sans espace en mode math
% Caractères mathématiques Unicode absents des polices de texte (Plus Jakarta,
% Archivo Black) : rendus par leur équivalent mathématique, en texte comme en
% formule — sinon ils s'affichent en carré vide (ex. « Arithmétique dans ℤ »).
\usepackage{newunicodechar}
\newunicodechar{ℝ}{\ensuremath{\mathbb{R}}}
\newunicodechar{ℤ}{\ensuremath{\mathbb{Z}}}
\newunicodechar{ℂ}{\ensuremath{\mathbb{C}}}
\newunicodechar{ℕ}{\ensuremath{\mathbb{N}}}
\newunicodechar{ℚ}{\ensuremath{\mathbb{Q}}}
\newunicodechar{𝔻}{\ensuremath{\mathbb{D}}}
\newunicodechar{α}{\ensuremath{\alpha}}
\newunicodechar{β}{\ensuremath{\beta}}
\newunicodechar{γ}{\ensuremath{\gamma}}
\newunicodechar{δ}{\ensuremath{\delta}}
\newunicodechar{ε}{\ensuremath{\varepsilon}}
\newunicodechar{θ}{\ensuremath{\theta}}
\newunicodechar{λ}{\ensuremath{\lambda}}
\newunicodechar{μ}{\ensuremath{\mu}}
\newunicodechar{σ}{\ensuremath{\sigma}}
\newunicodechar{τ}{\ensuremath{\tau}}
\newunicodechar{φ}{\ensuremath{\varphi}}
\newunicodechar{ψ}{\ensuremath{\psi}}
\newunicodechar{ω}{\ensuremath{\omega}}
\newunicodechar{Σ}{\ensuremath{\Sigma}}
% majuscules produites par \MakeUppercase dans les titres (α-aminés -> Α)
\newunicodechar{Α}{\ensuremath{\alpha}}
\newunicodechar{Β}{\ensuremath{\beta}}
\newunicodechar{Γ}{\ensuremath{\Gamma}}
\newunicodechar{Δ}{\ensuremath{\Delta}}
\newunicodechar{Ε}{\ensuremath{\varepsilon}}
\newunicodechar{Θ}{\ensuremath{\Theta}}
\newunicodechar{Λ}{\ensuremath{\Lambda}}
\newunicodechar{Μ}{\ensuremath{\mu}}
\newunicodechar{Τ}{\ensuremath{\tau}}
\newunicodechar{Φ}{\ensuremath{\Phi}}
\newunicodechar{Ψ}{\ensuremath{\Psi}}
\newunicodechar{Ω}{\ensuremath{\Omega}}
\newunicodechar{↦}{\ensuremath{\mapsto}}
\newunicodechar{∈}{\ensuremath{\in}}
\newunicodechar{∉}{\ensuremath{\notin}}
\newunicodechar{∧}{\ensuremath{\wedge}}
\newunicodechar{⇔}{\ensuremath{\Leftrightarrow}}
\newunicodechar{⟺}{\ensuremath{\Longleftrightarrow}}
\newunicodechar{⇒}{\ensuremath{\Rightarrow}}
\newunicodechar{⟹}{\ensuremath{\Longrightarrow}}
\newunicodechar{∘}{\ensuremath{\circ}}
\newunicodechar{∪}{\ensuremath{\cup}}
\newunicodechar{∩}{\ensuremath{\cap}}
\newunicodechar{≡}{\ensuremath{\equiv}}
\newunicodechar{⊂}{\ensuremath{\subset}}
\newunicodechar{⊥}{\ensuremath{\perp}}
\newunicodechar{∃}{\ensuremath{\exists}}
\newunicodechar{∀}{\ensuremath{\forall}}
\newunicodechar{∛}{\ensuremath{\sqrt[3]{\,}}}
\newunicodechar{∜}{\ensuremath{\sqrt[4]{\,}}}
\newunicodechar{⃗}{\ensuremath{{}^{\to}}}
\newunicodechar{ⁿ}{\ifmmode^{n}\else\textsuperscript{\ensuremath{n}}\fi}
\newunicodechar{ˣ}{\ifmmode^{x}\else\textsuperscript{\ensuremath{x}}\fi}
\newunicodechar{ᵇ}{\ifmmode^{b}\else\textsuperscript{\ensuremath{b}}\fi}
\newunicodechar{ʳ}{\ifmmode^{r}\else\textsuperscript{\ensuremath{r}}\fi}
\newunicodechar{ᵘ}{\ifmmode^{u}\else\textsuperscript{\ensuremath{u}}\fi}
\newunicodechar{ʸ}{\ifmmode^{y}\else\textsuperscript{\ensuremath{y}}\fi}
\newunicodechar{ᵃ}{\ifmmode^{a}\else\textsuperscript{\ensuremath{a}}\fi}
\newunicodechar{ᵉ}{\ifmmode^{e}\else\textsuperscript{\ensuremath{e}}\fi}
\newunicodechar{ᵏ}{\ifmmode^{k}\else\textsuperscript{\ensuremath{k}}\fi}
\newunicodechar{ᵖ}{\ifmmode^{p}\else\textsuperscript{\ensuremath{p}}\fi}
\newunicodechar{ᴺ}{\ifmmode^{N}\else\textsuperscript{\ensuremath{N}}\fi}
\newunicodechar{ᶜ}{\ifmmode^{c}\else\textsuperscript{\ensuremath{c}}\fi}
\newunicodechar{ʰ}{\ifmmode^{h}\else\textsuperscript{\ensuremath{h}}\fi}
\newunicodechar{ᵠ}{\ifmmode^{q}\else\textsuperscript{\ensuremath{q}}\fi}
\newunicodechar{ₙ}{\ifmmode_{n}\else\textsubscript{\ensuremath{n}}\fi}
\newunicodechar{ₐ}{\ifmmode_{a}\else\textsubscript{\ensuremath{a}}\fi}
\newunicodechar{ᵢ}{\ifmmode_{i}\else\textsubscript{\ensuremath{i}}\fi}
\newunicodechar{ᵣ}{\ifmmode_{r}\else\textsubscript{\ensuremath{r}}\fi}
\newunicodechar{ₖ}{\ifmmode_{k}\else\textsubscript{\ensuremath{k}}\fi}
\newunicodechar{ᵦ}{\ifmmode_{\beta}\else\textsubscript{\ensuremath{\beta}}\fi}
\newunicodechar{ₓ}{\ifmmode_{x}\else\textsubscript{\ensuremath{x}}\fi}
\newfontfamily\popdisplay{ArchivoBlack-Regular.ttf}[Path=../../../fonts/]
\newfontfamily\poplogo{SpaceGrotesk-Bold.ttf}[Path=../../../fonts/]
\newfontfamily\popsemi{PlusJakartaSans-SemiBold.ttf}[Path=../../../fonts/]
\newfontfamily\popxbold{PlusJakartaSans-ExtraBold.ttf}[Path=../../../fonts/]
\newfontfamily\popxboldls{PlusJakartaSans-ExtraBold.ttf}[Path=../../../fonts/,LetterSpace=3.0]

\setlist[enumerate]{leftmargin=1.7em,itemsep=5pt,topsep=4pt}
\setlist[itemize]{leftmargin=1.7em,itemsep=4pt,topsep=4pt,label={\color{poporange}\textbullet}}
\setlength{\parindent}{0pt}
\setlength{\parskip}{5pt}
\renewcommand{\arraystretch}{1.25}

% pgfplots : style Pop par défaut
\pgfplotsset{
  every axis/.append style={axis line style={popink,line width=1pt},tick style={popink},
    grid style={popline,line width=0.5pt},label style={font=\small},tick label style={font=\footnotesize}},
  popcurve/.style={poporange,line width=1.8pt,no markers,smooth},
}
% Même style disponible dans un \draw TikZ simple (hors pgfplots)
\tikzset{popcurve/.style={poporange,line width=1.8pt}}
\tikzset{popgrid/.style={popline,very thin}}
\colorlet{popcurve}{poporange} % le nom du style est parfois utilisé comme couleur

% ============================================================
% Pill / squiggle
% ============================================================
\newcommand{\poppill}[3]{%
  \tikz[baseline=-0.55ex]\node[draw=popink,line width=1.2pt,rounded corners=2.6pt,%
    fill=#2,inner xsep=6.5pt,inner ysep=3pt]{%
    \popxboldls\fontsize{7}{7}\selectfont\color{#3}\MakeUppercase{#1}};%
}
\newcommand{\popsquiggle}[1]{%
  \tikz[baseline=-0.4ex]{
    \draw[#1,line width=2.6pt,line cap=round]
      plot [smooth] coordinates {(0,0.09) (0.34,0.01) (0.68,0.21) (1.02,0.01) (1.36,0.10)};
  }%
}
% Mot-clé surligné (marqueur orange sous le texte)
\newcommand{\cle}[1]{{\popxbold\color{popdark}#1}}

% ============================================================
% En-tête / pied
% ============================================================
\pagestyle{fancy}
\fancyhf{}
\renewcommand{\headrulewidth}{0pt}
\renewcommand{\footrulewidth}{0pt}
\lhead{\raisebox{-0.15ex}{\poplogo\fontsize{13}{13}\selectfont\color{popink}OnBuch\color{poporange}+}}
\rhead{\poppill{\DOCMATIERE\ \textbullet\ \DOCNIVEAU\ \textbullet\ Leçon \DOCLECON}{white}{popink}}
\lfoot{\popxbold\fontsize{7.6}{7.6}\selectfont\color{popmuted}\MakeUppercase{Cours signé OnBuch+}\ \textbullet\ {\popsemi\DOCTITRE}}
\rfoot{\tikz[baseline=-0.6ex]\node[rounded corners=8.5pt,fill=popink,minimum height=17pt,minimum width=17pt,inner xsep=5pt,inner sep=0]{\popxbold\fontsize{8}{8}\selectfont\color{popcream}\thepage\,/\,\pageref*{LastPage}};}

% ============================================================
% Titres
% ============================================================
\newcommand{\chaptitre}[2]{%
  \begin{tcolorbox}[enhanced,colback=poporange,colframe=popink,boxrule=2.6pt,arc=11pt,
      drop shadow=popink,left=15pt,right=15pt,top=12pt,bottom=11pt,before skip=3pt,after skip=18pt]
    \poppill{#2}{popink}{popcream}\par\vspace{9pt}
    {\raggedright\hyphenpenalty=10000\exhyphenpenalty=10000\popdisplay\fontsize{22}{24}\selectfont\color{popcream}\MakeUppercase{#1}\par}
  \end{tcolorbox}
}
% Section numérotée : pastille orange avec le numéro + titre Archivo
\newcounter{popsec}
\newcommand{\coursec}[1]{%
  \stepcounter{popsec}\setcounter{popsub}{0}%
  \par\vspace{16pt}\noindent
  \begin{minipage}{\linewidth}
  \tikz[baseline=-0.7ex]\node[draw=popink,line width=1.6pt,fill=poporange,rounded corners=4pt,minimum size=22pt,inner sep=0]{\popdisplay\fontsize{12}{12}\selectfont\color{popcream}\thepopsec};%
  \hspace{8pt}{\popdisplay\fontsize{15}{17}\selectfont\color{popink}\MakeUppercase{#1}}\par
  \vspace{4pt}\noindent{\color{poporange}\rule{36pt}{3.6pt}}
  \end{minipage}\par\nopagebreak\vspace{8pt}
}
\newcounter{popsub}
\newcommand{\courssub}[1]{%
  \stepcounter{popsub}%
  \par\vspace{9pt}\noindent{\popxbold\fontsize{11.5}{13}\selectfont\color{popink}{\color{poporange}\thepopsec.\thepopsub}\hspace{6pt}#1}\par\nopagebreak\vspace{4pt}
}

% ============================================================
% Boîtes pédagogiques — même recette : fond blanc, bordure épaisse colorée,
% ombre dure encre, badge plein. Titre optionnel [#1].
% ============================================================
\tcbset{popbase/.style={breakable,enhanced,colback=white,boxrule=2.3pt,arc=9pt,
  shadow={3pt}{-3pt}{0pt}{popink},left=14pt,right=14pt,top=11pt,bottom=11pt,before skip=11pt,after skip=15pt,
  fontupper=\popsemi\fontsize{10.6}{14.5}\selectfont\color{popink},
  fontlower=\popsemi\fontsize{10.6}{14.5}\selectfont\color{popink}}}
% Coche dessinée (le glyphe ✓ n'existe pas dans les polices du document)
\newcommand{\popcheck}[1]{\tikz[baseline=-0.5ex]\draw[#1,line width=1.6pt,line cap=round,line join=round](-0.13,0.0)--(-0.04,-0.09)--(0.13,0.1);}
\providecommand{\coche}{\popcheck{popgreen}}
\providecommand{\resultat}[1]{\par\smallskip\noindent\fcolorbox{popgreen}{popgreenL}{\parbox{\dimexpr\linewidth-2\fboxsep-2\fboxrule\relax}{\textbf{Résultat.} #1}}\par\smallskip}
\newcommand{\popboxhead}[3]{\poppill{#1}{#2}{white}\ifx\relax#3\relax\else\hspace{6pt}{\popxbold\fontsize{10.6}{12}\selectfont\color{popink}#3}\fi\par\vspace{7pt}}

\newtcolorbox{aretenir}[1][]{popbase,colframe=popgreen,before upper={\popboxhead{À retenir}{popgreen}{#1}}}
% Texte en allemand (document de lecture, dialogue, extrait) : cadre
% d'étude. Un titre optionnel (source, auteur) s'affiche dans l'en-tête.
\newtcolorbox{textebox}[1][]{popbase,colframe=popblue,before upper={\popboxhead{Text}{popblue}{#1}\begin{otherlanguage}{ngerman}},after upper={\end{otherlanguage}}}
% Marques ✓ / ✗ (forme correcte / fautive) souvent utilisées par les modèles
\providecommand{\cross}{\textcolor{poppink}{\ensuremath{\times}}}
\providecommand{\xmark}{\cross}\providecommand{\crossmark}{\cross}\providecommand{\cmark}{\ensuremath{\checkmark}}
% Ligne de réponse / justification à compléter (inventée par les modèles)
\providecommand{\justification}[1]{\par\noindent\textit{Justification~:} \dotfill\par}
% \textipa (package tipa non chargé) : la police ne couvre pas l'API -> simple passage
\providecommand{\textipa}[1]{#1}
% Soulignement (ulem non chargé) : \uline, \ul
\providecommand{\uline}[1]{\underline{#1}}\providecommand{\ul}[1]{\underline{#1}}
% \alert (beamer) inventée par les modèles : texte mis en évidence
\providecommand{\alert}[1]{\textbf{\textcolor{poppink}{#1}}}
% variantes ASCII d'ordinaux écrites en macro
\providecommand{\iere}{\textsuperscript{re}}\providecommand{\ieme}{\textsuperscript{e}}\providecommand{\ere}{\textsuperscript{re}}\providecommand{\eme}{\textsuperscript{e}}\providecommand{\er}{\textsuperscript{er}}\providecommand{\ie}{\textsuperscript{e}}
% Ordinaux français écrits en macro par les modèles : 1\ière, 3\ième
\newcommand{\ière}{\textsuperscript{re}}\newcommand{\ième}{\textsuperscript{e}}
% \exemple (inventée par les modèles) : étiquette « Exemple : »
\providecommand{\exemple}{\textbf{Exemple~:}\ }
% \square utilisable aussi hors mode math (cases à cocher)
\AtBeginDocument{\let\squareMath\square\renewcommand{\square}{\ensuremath{\squareMath}}}
% Mot ou phrase en allemand à l'intérieur d'un paragraphe français
\newcommand{\all}[1]{\ifmmode\text{\textit{#1}}\else\begingroup\selectlanguage{ngerman}\itshape #1\endgroup\fi}
\let\esp\all % alias toléré
\newtcolorbox{exemplebox}[1][]{popbase,colframe=poporange,before upper={\popboxhead{Exemple}{poporange}{#1}}}
\newtcolorbox{definition}[1][]{popbase,colframe=popblue,before upper={\popboxhead{Définition}{popblue}{#1}}}
\newtcolorbox{propriete}[1][]{popbase,colframe=poppurple,before upper={\popboxhead{Propriété}{poppurple}{#1}}}
\newtcolorbox{methode}[1][]{popbase,colframe=popgold,before upper={\popboxhead{Méthode}{popgold}{#1}}}
\newtcolorbox{attention}[1][]{popbase,colframe=poppink,before upper={\popboxhead{Attention}{poppink}{#1}}}
\newtcolorbox{experience}[1][]{popbase,colframe=popblue,colback=popblueL,before upper={\popboxhead{Expérience}{popblue}{#1}}}
% « Le savais-tu ? » : carte violette pleine (contexte, culture, Cameroun)
\newtcolorbox{savaistu}[1][]{popbase,colback=poppurple,colframe=popink,
  fontupper=\popsemi\fontsize{10.4}{14.2}\selectfont\color{white},
  before upper={\poppill{Le savais-tu\,?}{white}{poppurple}\ifx\relax#1\relax\else\hspace{6pt}{\popxbold\color{white}#1}\fi\par\vspace{7pt}}}
% Exercice résolu : énoncé en haut, \tcblower puis la solution sur fond crème
\newcounter{popexo}\newcounter{popexr}
\newtcolorbox{exoresolu}[1][]{popbase,colframe=poporange,colbacklower=popcream,
  segmentation style={popink,line width=1.2pt,dashed},
  before upper={\stepcounter{popexr}\popboxhead{Exercice résolu \thepopexr}{poporange}{#1}},
  before lower={\poppill{Solution}{popink}{popcream}\par\vspace{6pt}}}
% Exercice (non résolu en place) — la correction est dans la partie Corrigés
\newtcolorbox{exercice}[1][]{popbase,colframe=popink,
  before upper={\stepcounter{popexo}\popboxhead{Exercice \thepopexo}{popink}{#1}}}
% Corrigé
\newtcolorbox{corrige}[1][]{popbase,colframe=popgreen,colback=popgreenL,
  before upper={\popboxhead{Corrigé}{popgreen}{#1}}}
% Objectifs / prérequis / bilan
\newtcolorbox{objectifs}{popbase,colframe=popink,colback=poporangeL,
  before upper={\popboxhead{Objectifs de la leçon}{popink}{}}}
\newtcolorbox{prerequis}{popbase,colframe=popink,colback=popblueL,
  before upper={\popboxhead{Prérequis}{popblue}{}}}
\newtcolorbox{bilan}{popbase,colframe=popink,colback=white,boxrule=2.8pt,
  before upper={\popboxhead{Fiche bilan}{poporange}{L'essentiel à connaître pour l'examen}}}
% Figure encadrée avec légende
\newtcolorbox{popfigure}[1][]{enhanced,breakable=false,colback=white,colframe=popline,boxrule=1.4pt,arc=8pt,
  left=10pt,right=10pt,top=10pt,bottom=8pt,before skip=10pt,after skip=12pt,halign=center,
  lower separated=false,halign lower=center,
  fontlower=\popsemi\fontsize{9}{11.5}\selectfont\color{popmuted},
  #1}
\newcommand{\legende}[1]{\par\vspace{6pt}{\popsemi\fontsize{9}{11.5}\selectfont\color{popmuted}#1}}

% ============================================================
% Signature OnBuch+
% ============================================================
\newcommand{\poplogoinline}[1]{{\poplogo\fontsize{#1}{#1}\selectfont\color{popink}OnBuch\color{poporange}+}}
\newcommand{\signatureonbuch}{%
  \begin{tcolorbox}[enhanced,colback=popink,colframe=popink,boxrule=2.2pt,arc=10pt,
      drop shadow=poporange,left=14pt,right=14pt,top=10pt,bottom=10pt,before skip=0pt,after skip=0pt]
    \tikz[baseline=-0.7ex]\node[circle,fill=popgreen,draw=popcream,line width=1.3pt,minimum size=22pt,inner sep=0]{\popcheck{white}};%
    \hspace{9pt}\begin{minipage}[c]{0.62\linewidth}
      {\popxbold\fontsize{12}{13}\selectfont\color{popcream}Signé\ \textbullet\ {\poplogo OnBuch\color{poporange}+}}\par\vspace{2pt}
      {\raggedright\popsemi\fontsize{8}{10.5}\selectfont\color{popline}\MakeUppercase{Équipe pédagogique OnBuch+}\\\MakeUppercase{Conforme au programme officiel MINESEC}\par}
    \end{minipage}\hfill
    {\poplogo\fontsize{22}{22}\selectfont\color{popcream}OnBuch\color{poporange}+}
  \end{tcolorbox}%
}

% ============================================================
% Page de garde
% #1 titre · #2 matière · #3 niveau · #4 module · #5 n° leçon · #6 durée ·
% #7 description / objectifs (texte ou itemize)
% ============================================================
\newcommand{\pagegarde}[7]{%
  \thispagestyle{empty}
  \newgeometry{a4paper,margin=1.7cm,top=1.6cm,bottom=1.4cm}%
  \begin{tikzpicture}[remember picture,overlay]
    \fill[poporange!18] ([xshift=-3.2cm,yshift=-3.6cm]current page.north east) circle (4.6cm);
    \fill[poppurple!14] ([xshift=2.4cm,yshift=7.6cm]current page.south west) circle (3.8cm);
  \end{tikzpicture}%
  {\poplogo\fontsize{30}{30}\selectfont\color{popink}OnBuch\color{poporange}+}\hfill
  \raisebox{0.6ex}{\poppill{Cours signé \textbullet\ Premium}{popink}{popcream}}\par
  \vspace{1pt}\popsquiggle{poppurple}\par
  \vspace{8pt}
  \begin{tcolorbox}[enhanced,colback=poporange,colframe=popink,boxrule=2.8pt,arc=13pt,
      drop shadow=popink,left=18pt,right=18pt,top=12pt,bottom=13pt,after skip=10pt]
    \poppill{#2 \textbullet\ #3}{popink}{popcream}\hspace{5pt}\poppill{#4}{white}{popink}\par\vspace{8pt}
    {\popdisplay\fontsize{14}{14}\selectfont\color{popink}LEÇON #5}\par\vspace{4pt}
    {\raggedright\hyphenpenalty=10000\exhyphenpenalty=10000\popdisplay\fontsize{25}{27}\selectfont\color{popcream}\MakeUppercase{#1}\par}
    \vspace{8pt}
    {\popxbold\fontsize{9.5}{9.5}\selectfont\color{popink}\MakeUppercase{Durée conseillée : #6}}
  \end{tcolorbox}
  \begin{tcolorbox}[enhanced,colback=white,colframe=popink,boxrule=2.2pt,arc=10pt,
      drop shadow=popink,left=16pt,right=16pt,top=10pt,bottom=10pt,after skip=8pt]
    \poppill{Dans cette leçon}{poppurple}{white}\par\vspace{5pt}
    \popsemi\fontsize{9.3}{12.6}\selectfont\color{popink}#7
  \end{tcolorbox}
  \noindent\begin{minipage}[t]{0.487\linewidth}
    \begin{tcolorbox}[enhanced,colback=popgreen,colframe=popink,boxrule=2.2pt,arc=10pt,
        drop shadow=popink,left=11pt,right=11pt,top=10pt,bottom=10pt,height=3.1cm,valign=center]
      \tcbox[colback=white,colframe=popink,boxrule=1.3pt,arc=5pt,boxsep=0pt,nobeforeafter,
        left=3pt,right=3pt,top=3pt,bottom=3pt,valign=center]{\qrcode[height=1.9cm]{https://play.google.com/store/apps/details?id=com.luvvix.onbuch&hl=fr}}%
      \hspace{8pt}\begin{minipage}[c]{0.5\linewidth}
        {\popdisplay\fontsize{11}{12.5}\selectfont\color{white}\MakeUppercase{Télécharge OnBuch}}\par\vspace{4pt}
        {\raggedright\popsemi\fontsize{8.4}{11}\selectfont\color{white}Gratuit \textbullet\ Google Play\\Tuteur IA, annales, concours, résultats\par}
      \end{minipage}
    \end{tcolorbox}
  \end{minipage}\hfill
  \begin{minipage}[t]{0.487\linewidth}
    \begin{tcolorbox}[enhanced,colback=poppurple,colframe=popink,boxrule=2.2pt,arc=10pt,
        drop shadow=popink,left=11pt,right=11pt,top=10pt,bottom=10pt,height=3.1cm,valign=center]
      \tikz[baseline=-0.7ex]\node[circle,fill=white,draw=popink,line width=1.3pt,minimum size=1.6cm,inner sep=0]{\popdisplay\fontsize{24}{24}\selectfont\color{poppurple}+};%
      \hspace{8pt}\begin{minipage}[c]{0.6\linewidth}
        {\popdisplay\fontsize{11}{12.5}\selectfont\color{white}\MakeUppercase{Passe à OnBuch+}}\par\vspace{4pt}
        {\raggedright\popsemi\fontsize{8.4}{11}\selectfont\color{white}Tous les cours signés, Tuteur IA illimité, fiches PDF — onglet OnBuch+ de l'app\par}
      \end{minipage}
    \end{tcolorbox}
  \end{minipage}
  \vspace{8pt}
  \popcontenu
  \vfill
  \signatureonbuch
  \restoregeometry
  \newpage
}

% Rangée « Ce cours contient » (page de garde)
\newcommand{\poptile}[3]{% #1 couleur #2 gros texte #3 légende
  \begin{tcolorbox}[enhanced,colback=#1,colframe=popink,boxrule=2pt,arc=9pt,shadow={2.5pt}{-2.5pt}{0pt}{popink},
      left=6pt,right=6pt,top=9pt,bottom=9pt,halign=center,valign=center,height=1.75cm,width=\linewidth,nobeforeafter]
    {\popdisplay\fontsize{12.5}{13}\selectfont\color{white}\MakeUppercase{#2}}\par\vspace{3pt}
    {\centering\popsemi\fontsize{7.8}{9.5}\selectfont\color{white}#3\par}
  \end{tcolorbox}}
\newcommand{\popcontenu}{%
  \par\noindent{\popxboldls\fontsize{7.5}{7.5}\selectfont\color{popmuted}CE COURS SIGNÉ CONTIENT}\par\vspace{7pt}
  \noindent\begin{minipage}[t]{0.235\linewidth}\poptile{popblue}{Cours}{complet, illustré, pas à pas}\end{minipage}\hfill
  \begin{minipage}[t]{0.235\linewidth}\poptile{poporange}{Méthodes}{et exercices résolus}\end{minipage}\hfill
  \begin{minipage}[t]{0.235\linewidth}\poptile{poppink}{Exercices}{type examen + corrigés}\end{minipage}\hfill
  \begin{minipage}[t]{0.235\linewidth}\poptile{popgreen}{Fiche bilan}{l'essentiel à retenir}\end{minipage}\par}

% Sommaire
\newtcolorbox{sommaire}{enhanced,colback=white,colframe=popink,boxrule=2.2pt,arc=10pt,
  drop shadow=popink,left=18pt,right=18pt,top=13pt,bottom=13pt,before skip=6pt,after skip=20pt,
  before upper={\poppill{Sommaire}{poppurple}{white}\par\vspace{9pt}\popsemi\fontsize{10.6}{14.5}\selectfont\color{popink}}}

% Signature de fin de cours — poussée tout en bas de la dernière page
\newcommand{\findecours}{%
  \par\vfill\nopagebreak
  \begin{center}\popsquiggle{poporange}\end{center}\vspace{4pt}
  \signatureonbuch
}
\AtBeginDocument{\def\llbracket{\mathopen{[\mkern-3.5mu[}}\def\rrbracket{\mathclose{]\mkern-3.5mu]}}} % intervalles d'entiers (glyphe ⟦ absent de la police)
% Glyphes absents de Fira Math : redéfinis à partir de glyphes disponibles
\AtBeginDocument{%
  \def\setminus{\mathbin{\backslash}}\let\smallsetminus\setminus
  \def\vdots{\mathord{\vbox{\baselineskip3.2pt\lineskiplimit0pt\kern4pt\hbox{.}\hbox{.}\hbox{.}}}}%
  \def\ddots{\mathinner{\mkern1mu\raise6.5pt\hbox{.}\mkern2mu\raise3.6pt\hbox{.}\mkern2mu\raise0.7pt\hbox{.}\mkern1mu}}%
  \def\triangle{\mathord{\scalebox{0.9}{$\symup{\Delta}$}}}%
  \def\nabla{\mathord{\raisebox{\depth}{\scalebox{1}[-1]{$\symup{\Delta}$}}}}%
  \def\checkmark{\popcheck{popdark}}%
  \def\star{\mathbin{\ast}}\def\bigstar{\mathord{\ast}}%
  \def\diamond{\mathbin{\scalebox{0.6}{$\Diamond$}}}%
  \def\bigcup{\mathop{\mathchoice{\raisebox{-0.3em}{\scalebox{1.7}{$\cup$}}}{\scalebox{1.25}{$\cup$}}{\cup}{\cup}}\displaylimits}%
  \def\bigcap{\mathop{\mathchoice{\raisebox{-0.3em}{\scalebox{1.7}{$\cap$}}}{\scalebox{1.25}{$\cap$}}{\cap}{\cap}}\displaylimits}%
  \def\lesseqqgtr{\mathrel{\lessgtr}}\def\bigoplus{\mathop{\oplus}\displaylimits}%
}
\providecommand{\ssi}{si et seulement si} % abréviation fréquente
\AtBeginDocument{\providecommand{\wideparen}{\overparen}} % arc de cercle

%% FIGFIX-BEGIN v5 — correctifs globaux des figures (docs/onbuch-generation/tools/figfix)
\usepackage{adjustbox}\usepackage{etoolbox}\tcbuselibrary{raster}
% toute figure TikZ/pgfplots plus large que la ligne est réduite à la largeur disponible
\BeforeBeginEnvironment{tikzpicture}{\begin{adjustbox}{max width=\linewidth}}
\AfterEndEnvironment{tikzpicture}{\end{adjustbox}}
% légendes pgfplots lisibles par-dessus les courbes
\pgfplotsset{every axis/.append style={legend style={fill=white,fill opacity=0.92,text opacity=1,draw=black!15,inner sep=3pt}}}
% étiquettes d'arêtes (syntaxe edge["..."]) avec fond blanc
\tikzset{every edge quotes/.append style={fill=white,inner sep=1.2pt}}
%% FIGFIX-END
\begin{document}

\pagegarde{Les maladies : la crise sanitaire et la médecine traditionnelle}{Allemand}{Tle A}{Chapitre 3 : Environnement, santé et bien-être}{AL14}{2 h}{Cette leçon propose l'étude d'un texte sur les maladies, la crise sanitaire et la médecine traditionnelle. Elle permet d'acquérir le vocabulaire thématique, de maîtriser les verbes pronominaux et le conseil au subjonctif II, et de produire un court texte argumenté.
\vspace{4pt}
\begin{multicols}{2}\small
\begin{itemize}
\item À la fin de cette leçon, je sais comprendre globalement et en détail un texte allemand sur les maladies et la médecine traditionnelle.
\item Je sais employer correctement le vocabulaire du thème : die Krankheit, der Arzt, das Heilmittel, die Heilpflanze, die Impfung, die Pandemie.
\item Je sais reconnaître et conjuguer les verbes pronominaux au présent, au passé composé et au prétérit.
\item Je sais distinguer les verbes pronominaux réfléchis, réciproques et idiomatiques.
\item Je sais formuler un conseil avec le subjonctif II (sollen, müssten, könnten, wäre besser).
\item Je sais répondre à des questions de compréhension en phrases complètes.
\end{itemize}\end{multicols}}

\begin{sommaire}
\begin{multicols}{2}
\begin{enumerate}
\item Texte support : lecture et première compréhension
\item Lexique thématique : la santé et la médecine
\item Grammaire 1 : les verbes pronominaux
\item Grammaire 2 : le conseil avec le subjonctif II
\item Méthode du commentaire de texte et commentaire modèle
\item Production guidée et réinvestissement
\item Activité d'intégration
\item Méthodes et exercices résolus
\item Exercices
\item Corrigés des exercices
\item Fiche bilan
\end{enumerate}
\end{multicols}
\end{sommaire}

\begin{objectifs}
\begin{itemize}
\item À la fin de cette leçon, je sais comprendre globalement et en détail un texte allemand sur les maladies et la médecine traditionnelle.
\item Je sais employer correctement le vocabulaire du thème : die Krankheit, der Arzt, das Heilmittel, die Heilpflanze, die Impfung, die Pandemie.
\item Je sais reconnaître et conjuguer les verbes pronominaux au présent, au passé composé et au prétérit.
\item Je sais distinguer les verbes pronominaux réfléchis, réciproques et idiomatiques.
\item Je sais formuler un conseil avec le subjonctif II (sollen, müssten, könnten, wäre besser).
\item Je sais répondre à des questions de compréhension en phrases complètes.
\item Je sais résumer un texte et rédiger un court paragraphe sur la santé au Cameroun.
\item Je sais commenter un texte selon la méthode en trois parties (introduction, développement, conclusion).
\end{itemize}
\end{objectifs}

\begin{prerequis}
\begin{itemize}
\item Le présent et le passé composé des verbes faibles et forts (Première)
\item Les pronoms personnels à l'accusatif et au datif (mich/mir, dich/dir, sich)
\item Les verbes modaux au présent et au prétérit (können, müssen, sollen, dürfen, wollen, mögen)
\item La formation de base du subjonctif II (hätte, wäre, würde + infinitif)
\item Le vocabulaire général du corps et de la santé vu en Seconde/Première
\end{itemize}
\end{prerequis}

\newpage

\coursec{Texte support : lecture et première compréhension}

\begin{experience}[Accroche : vaccin ou tisane ?]
Pendant la crise du Covid-19, les Camerounais ont débattu entre vaccins et tisanes traditionnelles à base de citronnelle ou d'artemisia. Sur les marchés de Yaoundé et de Douala, on vendait des bottes de \all{Zitronengras} (citronnelle) et les grands-mères préparaient leurs infusions. En Allemagne aussi, la pandémie a bouleversé le quotidien : masques, \all{Impfung} (vaccination), confinement. Mais en même temps, beaucoup d'Allemands se tournent vers la médecine douce et les plantes médicinales (\all{Heilpflanzen}). \textbf{Problématique :} comment parler en allemand des maladies, des crises sanitaires et de la médecine traditionnelle ? C'est l'objet de cette leçon. Notre point de départ sera un texte de presse original, \all{Zwischen Impfung und Heilpflanze} (« Entre vaccin et plante médicinale »).
\end{experience}

\courssub{Le texte : Zwischen Impfung und Heilpflanze}

\begin{methode}[Comment aborder un texte allemand ?]
Avant de lire, adoptez la bonne démarche, en cinq étapes :
\begin{enumerate}
\item Lire le \textbf{titre} et deviner le thème : \all{Zwischen Impfung und Heilpflanze} annonce une opposition ou une complémentarité entre deux types de médecine.
\item Faire une \textbf{première lecture globale}, sans dictionnaire, pour saisir l'idée générale.
\item \textbf{Souligner les mots transparents} (mots proches du français) : ils donnent des repères gratuits.
\item Faire une \textbf{deuxième lecture détaillée} avec le glossaire.
\item Répondre aux questions \textbf{en phrases complètes}, en reprenant les mots du texte.
\end{enumerate}
\end{methode}

\textbf{Lecture silencieuse.} Lisez le texte une première fois en silence, sans vous arrêter sur les mots inconnus. Puis lisez-le à voix haute, en respectant la mélodie de la phrase allemande : le verbe conjugué occupe toujours la deuxième place.

\begin{textebox}[Zwischen Impfung und Heilpflanze]
Als die Pandemie Kamerun erreichte, fürchteten sich viele Menschen. Die Krankenhäuser waren voll, und die Ärzte arbeiteten Tag und Nacht. Die Regierung empfahl die Impfung, aber viele Bürger misstrauten dem neuen Heilmittel. In den Dörfern bereiteten die traditionellen Heiler Tees aus Heilpflanzen wie Zitronengras und Artemisia zu. „Unsere Großeltern haben sich immer mit Pflanzen geheilt“, sagt Mama Ngo, eine Heilerin aus Bafoussam. Dr. Mbarga, Arzt in Yaoundé, erklärt dagegen: „Die Impfung schützt vor schweren Krankheiten. Man sollte sich impfen lassen.“ Heute versuchen viele Kameruner, beide Medizinen zu verbinden: Sie gehen zum Arzt, aber sie trinken auch den Tee der Großmutter. Diese Krise hat gezeigt, dass moderne und traditionelle Medizin zusammenarbeiten können.
\end{textebox}

\textbf{Traduction guidée.} « Lorsque la pandémie a atteint le Cameroun, beaucoup de gens ont eu peur. Les hôpitaux étaient pleins et les médecins travaillaient jour et nuit. Le gouvernement a recommandé la vaccination, mais beaucoup de citoyens se méfiaient du nouveau remède. Dans les villages, les guérisseurs traditionnels préparaient des tisanes à base de plantes médicinales comme la citronnelle et l'artemisia. "Nos grands-parents se sont toujours soignés avec des plantes", dit Mama Ngo, une guérisseuse de Bafoussam. Le Dr Mbarga, médecin à Yaoundé, explique au contraire : "La vaccination protège des maladies graves. On devrait se faire vacciner." Aujourd'hui, beaucoup de Camerounais essaient de combiner les deux médecines : ils vont chez le médecin, mais ils boivent aussi la tisane de la grand-mère. Cette crise a montré que la médecine moderne et la médecine traditionnelle peuvent travailler ensemble. »

\courssub{Les mots transparents : des amis précieux}

\begin{definition}[Die Pandemie]
\all{Die Pandemie} (pluriel : \all{die Pandemien}) : \emph{eine Krankheit, die sich über viele Länder oder die ganze Welt ausbreitet} — une maladie qui se répand sur de nombreux pays ou sur le monde entier : une pandémie.
\end{definition}

Un \cle{mot transparent} est un mot allemand dont la forme ressemble au français : on peut le comprendre sans l'avoir appris. Le texte en contient plusieurs :

\begin{center}
\begin{tabularx}{\linewidth}{|X|X|X|}
\hline
\rowcolor{popblueL} Deutsch & Français & Remarque \\
\hline
die Pandemie (-n) & la pandémie & presque identique \\
\hline
das Virus (Viren) & le virus & attention : neutre en allemand \\
\hline
die Impfung (-en) & la vaccination & de \all{impfen} : vacciner \\
\hline
der Patient (-en) & le patient & masculin faible : \all{den Patienten} \\
\hline
die Medizin (-en) & la médecine & aussi : le médicament \\
\hline
das Zitronengras & la citronnelle & mot composé : \all{Zitrone} + \all{Gras} \\
\hline
\end{tabularx}
\end{center}

\begin{attention}[Faux amis et pièges]
\begin{itemize}
\item \all{Die Medizin} signifie « la médecine » (la science) mais aussi « le médicament » ; pour « un médicament » précis, on dit plutôt \all{das Medikament} ou \all{das Heilmittel}.
\item \all{das Virus} est neutre, alors que « le virus » est masculin en français : \all{das Virus, die Viren}.
\item \all{der Heiler} n'est pas « le healer » : c'est « le guérisseur ».
\item N'oubliez pas la majuscule à tous les noms : \all{die Impfung}, \all{die Krise}, \all{der Arzt}.
\end{itemize}
\end{attention}

\courssub{Compréhension globale}

Après la première lecture, répondez mentalement aux quatre questions fondamentales :

\begin{itemize}
\item \textbf{De quoi parle le texte ?} (\all{Worum geht es?}) Il parle de la crise sanitaire au Cameroun et de la coexistence entre médecine moderne et médecine traditionnelle.
\item \textbf{Qui ?} (\all{Wer?}) Les Camerounais, les médecins (Dr. Mbarga), les guérisseurs (Mama Ngo), le gouvernement.
\item \textbf{Où ?} (\all{Wo?}) Au Cameroun : dans les hôpitaux de Yaoundé, dans les villages, à Bafoussam.
\item \textbf{Quel problème ?} (\all{Welches Problem?}) Pendant la pandémie, la population hésite entre la vaccination recommandée par l'État et les remèdes traditionnels à base de plantes.
\end{itemize}

\begin{exemplebox}[Deux phrases clés du texte, traduites]
\all{Die Regierung empfahl die Impfung.} = « Le gouvernement a recommandé la vaccination. »\par
\all{Sie trinken den Tee der Großmutter.} = « Ils boivent la tisane de la grand-mère. »
\end{exemplebox}

\courssub{Compréhension détaillée : les deux arguments}

Le texte est construit comme un débat équilibré. Repérons les arguments de chaque camp :

\begin{center}
\begin{tabularx}{\linewidth}{|X|X|}
\hline
\rowcolor{popgreenL} Pour la médecine moderne & Pour la médecine traditionnelle \\
\hline
\all{Die Impfung schützt vor schweren Krankheiten.} (La vaccination protège des maladies graves.) & \all{Unsere Großeltern haben sich immer mit Pflanzen geheilt.} (Nos grands-parents se sont toujours soignés avec des plantes.) \\
\hline
\all{Man sollte sich impfen lassen.} (On devrait se faire vacciner.) & \all{Die Heiler bereiten Tees aus Heilpflanzen zu.} (Les guérisseurs préparent des tisanes de plantes médicinales.) \\
\hline
\all{Die Ärzte arbeiteten Tag und Nacht.} (Les médecins travaillaient jour et nuit.) & \all{Viele Bürger misstrauten dem neuen Heilmittel.} (Beaucoup de citoyens se méfiaient du nouveau remède.) \\
\hline
\end{tabularx}
\end{center}

\textbf{La conclusion du texte} est une synthèse : \all{Moderne und traditionelle Medizin können zusammenarbeiten.} « La médecine moderne et la médecine traditionnelle peuvent travailler ensemble. » Le verbe \all{zusammenarbeiten} (collaborer) est un verbe à particule séparable : \all{Sie arbeiten zusammen.}

\courssub{Glossaire des mots nouveaux}

Apprenez chaque nom avec son \textbf{article} et son \textbf{pluriel} : c'est indispensable en allemand.

\begin{center}
\begin{tabularx}{\linewidth}{|X|X|}
\hline
\rowcolor{popblueL} Deutsch & Français \\
\hline
die Krankheit (-en) & la maladie \\
\hline
das Krankenhaus ("-häuser) & l'hôpital \\
\hline
der Arzt ("-Ärzte) / die Ärztin (-nen) & le médecin / la femme médecin \\
\hline
die Impfung (-en) & la vaccination \\
\hline
das Heilmittel (-) & le remède \\
\hline
die Heilpflanze (-n) & la plante médicinale \\
\hline
der Heiler (-) / die Heilerin (-nen) & le guérisseur / la guérisseuse \\
\hline
der Bürger (-) & le citoyen \\
\hline
die Regierung (-en) & le gouvernement \\
\hline
der Tee (-s) & le thé, la tisane \\
\hline
die Krise (-n) & la crise \\
\hline
sich fürchten & avoir peur (verbe pronominal) \\
\hline
misstrauen (+ Dat.) & se méfier de \\
\hline
empfehlen (empfiehlt, empfahl, empfohlen) & recommander \\
\hline
schützen (vor + Dat.) & protéger (de) \\
\hline
sich impfen lassen & se faire vacciner \\
\hline
verbinden (verbindet, verband, verbunden) & combiner, relier \\
\hline
zubereiten (séparable) & préparer \\
\hline
\end{tabularx}
\end{center}

\begin{attention}[Prononciation et genre]
\begin{itemize}
\item \all{der Arzt} se prononce « artst » ; son pluriel \all{die Ärzte} prend l'Umlaut : « èr-ts-té ».
\item \all{die Krankheit} : le suffixe \cle{-heit} est toujours féminin, comme \all{die Gesundheit} (la santé), \all{die Freiheit}.
\item \all{das Krankenhaus} : le pluriel \all{die Krankenhäuser} prend l'Umlaut, comme \all{das Haus, die Häuser}.
\item \all{sich fürchten} : « zi-kh fü-rkh-tèn » ; le \all{ch} après \all{ü} est doux, jamais « k ».
\end{itemize}
\end{attention}

\courssub{Carte mentale du vocabulaire}

\begin{popfigure}
\begin{tikzpicture}[mindmap, grow cyclic, every node/.style={concept, align=center, font=\small},
root concept/.append style={concept color=popblue, text=white, font=\bfseries},
level 1/.append style={level distance=3.4cm, sibling angle=90},
level 1 concept/.append style={concept color=poporange!80, text=black}]
\node[root concept]{Gesundheit}
  child { node[align=center]{moderne Medizin\\ der Arzt\\ die Impfung\\ das Krankenhaus} }
  child { node[align=center]{traditionelle Medizin\\ der Heiler\\ die Heilpflanze\\ der Tee} }
  child { node[align=center]{Krankheiten\\ das Virus\\ der Patient\\ das Heilmittel} }
  child { node[align=center]{Krise\\ die Pandemie\\ die Regierung\\ die Angst} };
\end{tikzpicture}
\legende{Carte mentale : le champ lexical de la santé autour de \all{die Gesundheit}.}
\end{popfigure}

\courssub{Questions de compréhension}

Répondez par des \textbf{phrases complètes}, en reprenant les mots de la question. Rappel : en allemand, la réponse commence par le sujet (ou un complément), et le verbe conjugué reste en deuxième position.

\begin{exoresolu}[W-Fragen et Ja/Nein-Fragen]
\textbf{Énoncé.} Répondez aux questions suivantes :
\begin{enumerate}
\item \all{Wer arbeitete Tag und Nacht?}
\item \all{Was empfahl die Regierung?}
\item \all{Woher kommt Mama Ngo?}
\item \all{Misstrauten alle Bürger der Impfung?}
\item \all{Was trinken viele Kameruner heute?}
\end{enumerate}
\tcblower
\textbf{Solution détaillée.}
\begin{enumerate}
\item \all{Die Ärzte arbeiteten Tag und Nacht.} (Les médecins travaillaient jour et nuit.) — On reprend le verbe au Präteritum, comme dans la question.
\item \all{Die Regierung empfahl die Impfung.} (Le gouvernement a recommandé la vaccination.) — \all{empfehlen} est un verbe fort : \all{empfahl} au Präteritum.
\item \all{Mama Ngo kommt aus Bafoussam.} (Mama Ngo vient de Bafoussam.) — \all{Woher} demande l'origine : \all{kommen aus} + nom de ville sans article.
\item \all{Nein, nicht alle Bürger misstrauten der Impfung, aber viele misstrauten dem neuen Heilmittel.} (Non, pas tous, mais beaucoup se méfiaient du nouveau remède.) — Attention : \all{misstrauen} se construit avec le \textbf{datif} : \all{dem Heilmittel}, \all{der Impfung}.
\item \all{Sie trinken den Tee der Großmutter.} (Ils boivent la tisane de la grand-mère.)
\end{enumerate}
\end{exoresolu}

\begin{attention}[Erreurs fréquentes des francophones]
\begin{itemize}
\item Ne traduisez pas « se faire vacciner » mot à mot : il faut \all{sich impfen lassen} (\all{lassen} + infinitif), jamais \all{sich machen impfen}.
\item \all{Woher} (d'où) n'est pas \all{wo} (où) : \all{Woher kommt sie? — Aus Bafoussam.}
\item Dans la réponse, le verbe reste en deuxième position : \all{Die Ärzte arbeiteten...} et non « Arbeiteten die Ärzte » (ordre réservé aux questions sans mot interrogatif).
\item \all{viele Menschen} = « beaucoup de gens » ; \all{Menschen} prend ici la terminaison \all{-en} (déclinaison faible au pluriel, c'est normal).
\end{itemize}
\end{attention}

\courssub{Reformulation : die Zusammenfassung}

Résumer un texte (\all{zusammenfassen}) consiste à en redire l'essentiel avec ses propres mots, en trois ou quatre phrases, sans recopier les citations. Modèle de résumé :

\begin{exemplebox}[Modèle de Zusammenfassung]
\all{Der Text handelt von der Gesundheitskrise in Kamerun. Viele Menschen hatten Angst, und die Krankenhäuser waren voll. Der Staat wollte die Impfung, aber ein Teil der Bevölkerung vertraute lieber den Heilpflanzen der traditionellen Heiler. Am Ende zeigt der Text, dass beide Medizinen zusammenarbeiten können.}\par
« Le texte traite de la crise sanitaire au Cameroun. Beaucoup de gens avaient peur et les hôpitaux étaient pleins. L'État voulait la vaccination, mais une partie de la population faisait plutôt confiance aux plantes médicinales des guérisseurs traditionnels. Enfin, le texte montre que les deux médecines peuvent collaborer. »
\end{exemplebox}

\begin{methode}[Réussir sa Zusammenfassung]
\begin{enumerate}
\item Relevez les \textbf{informations essentielles} : qui, où, quel problème, quelle conclusion.
\item \textbf{Paraphrasez} : remplacez les mots du texte par des synonymes (\all{fürchten sich} $\rightarrow$ \all{Angst haben} ; \all{empfehlen} $\rightarrow$ \all{wollen}).
\item \textbf{Supprimez} les citations directes et les exemples secondaires (Zitronengras, Artemisia).
\item Rédigez au \textbf{présent} ou au \textbf{Präteritum}, en 3--4 phrases, avec des connecteurs : \all{zuerst} (d'abord), \all{aber} (mais), \all{am Ende} (enfin).
\end{enumerate}
\end{methode}

\begin{exercice}[À vous de résumer]
Rédigez en allemand votre propre résumé du texte en quatre phrases maximum, sans recopier les citations de Mama Ngo et du Dr Mbarga. Utilisez au moins trois des mots suivants : \all{die Krise, die Impfung, die Heilpflanzen, verbinden, zusammenarbeiten}.
\end{exercice}

\begin{corrige}[Exercice : la Zusammenfassung]
Proposition de corrigé : \all{Der Text beschreibt die Gesundheitskrise in Kamerun. Die Regierung empfahl die Impfung, aber viele Menschen bevorzugten die Heilpflanzen der Heiler. Heute versuchen die Kameruner, beide Medizinen zu verbinden. Die Krise zeigt: Moderne und traditionelle Medizin können zusammenarbeiten.} — Toute formulation correcte qui reprend les quatre idées (crise, vaccination, plantes, coopération) avec un allemand correct est acceptée.
\end{corrige}

\begin{savaistu}[La phytothérapie en Allemagne]
En Allemagne, la phytothérapie (\all{die Phytotherapie}) est reconnue officiellement : les médecins peuvent prescrire des plantes médicinales, et certaines sont même remboursées par l'assurance maladie (\all{die Krankenkasse}). La camomille (\all{die Kamille}) et la valériane (\all{der Baldrian}) font partie des remèdes les plus vendus en pharmacie. Comme au Cameroun, la médecine moderne et les plantes ne s'opposent donc pas toujours : elles se complètent souvent.
\end{savaistu}

\coursec{Lexique thématique : la santé et la médecine}

Dans la section précédente, nous avons lu et exploité le texte support sur la crise sanitaire et la médecine traditionnelle au Cameroun. Vous y avez rencontré des mots comme \all{die Krankheit}, \all{der Heiler}, \all{die Heilpflanze} ou \all{die Impfung}. Pour bien comprendre un tel texte — et surtout pour pouvoir en parler et écrire à son sujet — il faut maintenant organiser ce vocabulaire de manière systématique. C'est l'objet de cette section : construire, champ lexical par champ lexical, un réseau de mots solide, avec pour chaque nom son \cle{article} et son \cle{pluriel}, car en allemand, apprendre un nom sans son article, c'est ne l'apprendre qu'à moitié.

\courssub{Champ 1 : les maladies}

\begin{definition}[Les mots des maladies]
\begin{itemize}
\item \all{die Krankheit, -en} : la maladie
\item \all{die Erkältung, -en} : le rhume, le refroidissement
\item \all{das Fieber, -} : la fièvre
\item \all{die Grippe, -n} : la grippe
\item \all{die Epidemie, -n} : l'épidémie
\item \all{die Pandemie, -n} : la pandémie
\item \all{das Virus, Viren} : le virus
\item \all{die Ansteckung, -en} : la contagion, la contamination
\end{itemize}
\end{definition}

Observez la différence entre \all{die Epidemie} et \all{die Pandemie} : une \emph{épidémie} touche une région ou un pays (une épidémie de choléra à Douala), une \emph{pandémie} touche le monde entier. Notez aussi le pluriel irrégulier de \all{das Virus} : \all{Viren} (prononcez « vi-ren »), hérité du latin.

\begin{exemplebox}[Exemples traduits]
\all{Die Grippe ist eine ansteckende Krankheit.} « La grippe est une maladie contagieuse. »

\all{Viele Schüler in Yaoundé haben sich eine Erkältung geholt.} « Beaucoup d'élèves à Yaoundé ont attrapé un rhume. »

\all{Die Impfung schützt vor einer Ansteckung.} « La vaccination protège d'une contamination. »
\end{exemplebox}

\begin{aretenir}[Expressions figées à mémoriser]
\begin{itemize}
\item \all{sich eine Erkältung holen} : attraper un rhume (littéralement « se chercher un refroidissement »)
\item \all{Fieber haben} : avoir de la fièvre
\item \all{vor einer Krankheit schützen} : protéger d'une maladie
\item \all{an einer Krankheit leiden} : souffrir d'une maladie
\end{itemize}
\end{aretenir}

\courssub{Champ 2 : les personnes}

\begin{definition}[Les acteurs de la santé]
\begin{itemize}
\item \all{der Arzt, Ärzte} : le médecin (homme)
\item \all{die Ärztin, -nen} : la médecin (femme)
\item \all{der Patient, -en} : le patient
\item \all{die Patientin, -nen} : la patiente
\item \all{der Heiler, -} : le guérisseur, le tradipraticien
\item \all{die Heilerin, -nen} : la guérisseuse
\item \all{der Krankenpfleger, -} : l'infirmier
\item \all{die Krankenpflegerin, -nen} : l'infirmière
\end{itemize}
\end{definition}

Remarquez la régularité de la formation du féminin : on ajoute \all{-in} au nom masculin, et le pluriel devient \all{-nen} : \all{der Arzt} $\rightarrow$ \all{die Ärztin, -nen}. Attention, l'Umlaut apparaît souvent au féminin : \all{Arzt} $\rightarrow$ \all{Ärztin}.

\begin{exemplebox}[Exemples traduits]
\all{Der Patient leidet an Fieber.} « Le patient souffre de fièvre. »

\all{Die Ärztin behandelt die Krankheit.} « La médecin traite la maladie. »

\all{Der Heiler im Dorf kennt viele Heilpflanzen.} « Le tradipraticien du village connaît beaucoup de plantes médicinales. »
\end{exemplebox}

\courssub{Champ 3 : les remèdes et les soins}

\begin{definition}[Heilpflanze]
\all{die Heilpflanze, -n} : eine Pflanze, die gegen Krankheiten hilft — une plante médicinale, c'est-à-dire une plante qui aide à soigner les maladies. Exemple : \all{Artemisia und Zitronengras sind bekannte Heilpflanzen in Kamerun.} « L'artemisia et la citronnelle sont des plantes médicinales connues au Cameroun. »
\end{definition}

\begin{definition}[Les mots des remèdes]
\begin{itemize}
\item \all{das Heilmittel, -} : le remède
\item \all{die Medizin, -en} : la médecine (la science)
\item \all{das Medikament, -e} : le médicament
\item \all{die Heilpflanze, -n} : la plante médicinale
\item \all{der Tee, Tees} : le thé, la tisane
\item \all{die Impfung, -en} : la vaccination
\item \all{die Spritze, -n} : la piqûre, l'injection
\item \all{die Behandlung, -en} : le traitement
\end{itemize}
\end{definition}

\begin{attention}[Faux ami : « die Medizin » n'est pas « un médicament »]
Erreur très fréquente chez les francophones : \all{die Medizin} signifie \emph{la médecine} au sens de la science, de la discipline : \all{Er studiert Medizin.} « Il fait des études de médecine. » Pour dire « un médicament », employez \all{das Medikament} ou \all{das Heilmittel} : \all{Die Ärztin verschreibt ein Medikament.} « La médecin prescrit un médicament. »
\end{attention}

\courssub{Champ 4 : les lieux}

\begin{definition}[Où se soigne-t-on ?]
\begin{itemize}
\item \all{das Krankenhaus, Krankenhäuser} : l'hôpital
\item \all{die Klinik, -en} : la clinique
\item \all{die Praxis, Praxen} : le cabinet (médical)
\item \all{die Apotheke, -n} : la pharmacie
\end{itemize}
\end{definition}

Distinction utile : \all{das Krankenhaus} est le grand hôpital public (comme l'hôpital central de Yaoundé), \all{die Klinik} est souvent un établissement privé ou spécialisé, et \all{die Praxis} (pluriel \all{Praxen}) est le cabinet du médecin de ville. Notez le pluriel avec Umlaut : \all{das Krankenhaus} $\rightarrow$ \all{die Krankenhäuser}.

\courssub{Champ 5 : les verbes utiles}

\begin{definition}[Verbes du champ de la santé]
\begin{itemize}
\item \all{sich erkälten} : attraper un rhume, s'enrhumer
\item \all{sich fühlen} : se sentir (\all{Ich fühle mich krank.} « Je me sens malade. »)
\item \all{sich anstecken} : se contaminer, attraper (une maladie)
\item \all{heilen} : guérir
\item \all{behandeln} : traiter, soigner
\item \all{impfen} : vacciner
\item \all{leiden an (+ Dativ)} : souffrir de (\all{Er leidet an Fieber.})
\item \all{gesund werden} : guérir, retrouver la santé (devenir sain)
\end{itemize}
\end{definition}

Trois de ces verbes sont \cle{pronominaux} (\all{sich erkälten}, \all{sich fühlen}, \all{sich anstecken}) : ils seront étudiés en détail dans la section de grammaire. Retenez dès maintenant que le pronom réfléchi \all{sich} fait partie du verbe et se conjugue avec lui : \all{ich erkälte mich, du erkältest dich, er erkältet sich}.

\begin{exemplebox}[Exemples traduits]
\all{Nach der Behandlung wird der Patient gesund.} « Après le traitement, le patient guérit. »

\all{Die Krankenpflegerin impft die Kinder gegen die Krankheit.} « L'infirmière vaccine les enfants contre la maladie. »

\all{Ich habe mich bei meinem Bruder angesteckt.} « Je me suis contaminé chez mon frère / mon frère m'a transmis la maladie. »
\end{exemplebox}

\courssub{Tableau récapitulatif du lexique}

\begin{center}
\begin{tabularx}{\linewidth}{|X|X|X|}\hline
\rowcolor{popblueL} Deutsch & Français & Remarque \\ \hline
die Krankheit, -en & la maladie & féminin en \all{-heit} \\ \hline
die Erkältung, -en & le rhume & \all{sich eine Erkältung holen} \\ \hline
das Fieber, - & la fièvre & \all{Fieber haben} \\ \hline
das Virus, Viren & le virus & pluriel irrégulier \\ \hline
der Arzt, Ärzte / die Ärztin, -nen & le/la médecin & féminin en \all{-in} \\ \hline
der Heiler, - / die Heilerin, -nen & le/la tradipraticien(ne) & famille de \all{heilen} \\ \hline
das Heilmittel, - & le remède & neutre \\ \hline
das Medikament, -e & le médicament & ne pas confondre avec \all{die Medizin} \\ \hline
die Heilpflanze, -n & la plante médicinale & mot composé \\ \hline
die Impfung, -en & la vaccination & féminin en \all{-ung} \\ \hline
die Spritze, -n & la piqûre & \\ \hline
die Behandlung, -en & le traitement & féminin en \all{-ung} \\ \hline
das Krankenhaus, Krankenhäuser & l'hôpital & pluriel avec Umlaut \\ \hline
die Apotheke, -n & la pharmacie & \\ \hline
die Praxis, Praxen & le cabinet médical & \\ \hline
\end{tabularx}
\end{center}

\courssub{Une règle d'or pour les articles}

\begin{propriete}[Les suffixes féminins]
Les noms formés avec les suffixes \cle{-ung}, \cle{-heit}, \cle{-keit} et \cle{-ion} sont \textbf{toujours féminins} et prennent l'article \all{die} :
\begin{itemize}
\item \all{die Impfung, die Erkältung, die Behandlung, die Ansteckung} (suffixe \all{-ung})
\item \all{die Krankheit} (suffixe \all{-heit})
\item \all{die Tradition} (suffixe \all{-ion})
\end{itemize}
Cette règle est sans exception : elle vous donne gratuitement le genre de centaines de noms.
\end{propriete}

\courssub{Les familles de mots}

Apprendre le vocabulaire par familles est la méthode la plus efficace : un seul effort de mémorisation vous donne cinq ou six mots.

\begin{popfigure}
\begin{center}
\begin{tikzpicture}[mindmap, grow cyclic, every node/.style={concept, align=center, font=\small}, root concept/.append style={concept color=popblue, font=\large\bfseries}, level 1/.append style={level distance=3.4cm, sibling angle=60}, concept connection/.append style={color=popmuted}]
\node[root concept, align=center]{heil-\\ (sain, entier)}
child[concept color=popgreen] { node[align=center]{ \all{heilen}\\ guérir } }
child[concept color=poporange] { node[align=center]{ \all{der Heiler}\\ die Heilerin\\ le guérisseur } }
child[concept color=poppurple] { node[align=center]{ \all{das Heilmittel}\\ le remède } }
child[concept color=poppink] { node[align=center]{ \all{die Heilpflanze}\\ la plante\\ médicinale } }
child[concept color=popgold] { node[align=center]{ \all{die Heilung}\\ la guérison } }
child[concept color=popblue!60] { node[align=center]{ \all{der Heiland}\\ le Sauveur } };
\end{tikzpicture}
\end{center}
\legende{Carte mentale de la famille de mots autour de la racine « heil- »}
\end{popfigure}

De la même façon, observez les deux autres familles du thème :

\begin{center}
\begin{tabularx}{\linewidth}{|X|X|X|}\hline
\rowcolor{popblueL} Racine & Dérivés & Français \\ \hline
\all{krank} (malade) & \all{die Krankheit, das Krankenhaus, der Krankenpfleger} & la maladie, l'hôpital, l'infirmier \\ \hline
\all{impfen} (vacciner) & \all{die Impfung, der Impfstoff} & la vaccination, le vaccin \\ \hline
\all{heilen} (guérir) & \all{der Heiler, das Heilmittel, die Heilpflanze} & le guérisseur, le remède, la plante médicinale \\ \hline
\end{tabularx}
\end{center}

\begin{savaistu}[La racine « Heil- »]
Le mot \all{Heil-} vient de la même racine que l'adjectif \all{heil}, qui signifie « sain, entier, indemne ». De là découlent \all{heilen} (guérir, c'est-à-dire « rendre entier »), \all{das Heilmittel} (le remède, littéralement « le moyen de rendre sain ») et même \all{der Heiland} (le Sauveur, dans un contexte religieux). La langue allemande construit ainsi des familles de sens très transparentes : guérir, c'est « rendre entier ».
\end{savaistu}

\courssub{Texte d'application : vocabulaire en contexte}

\begin{textebox}[Beim Heiler im Dorf]
In einem Dorf in der Nähe von Bafoussam lebt der bekannte Heiler Herr Nkeng. Viele Patienten kommen zu ihm, weil sie an Fieber oder an einer Erkältung leiden. Herr Nkeng kennt viele Heilpflanzen: Artemisia gegen das Fieber, Zitronengras als Tee gegen die Erkältung. „Die Natur gibt uns viele Heilmittel“, sagt er oft.

Aber Herr Nkeng ist vorsichtig. Wenn ein Patient sehr krank ist, schickt er ihn ins Krankenhaus in die Stadt. „Manche Krankheiten kann ich nicht heilen“, erklärt er. „Ein Virus braucht manchmal die moderne Medizin und eine Impfung.“

Die Ärztin im Krankenhaus, Frau Dr. Mbarga, respektiert seine Arbeit. Sie behandelt die Patienten mit Medikamenten und Spritzen, aber sie sagt: „Die traditionelle und die moderne Medizin müssen zusammenarbeiten.“ So werden viele Menschen im Dorf gesund.
\end{textebox}

\textbf{Glossaire} : \all{der Heiler} le tradipraticien ; \all{die Heilpflanze} la plante médicinale ; \all{vorsichtig} prudent ; \all{schicken} envoyer ; \all{respektieren} respecter ; \all{zusammenarbeiten} collaborer ; \all{gesund werden} guérir.

\textbf{Questions de compréhension} :
\begin{enumerate}
\item \all{Warum kommen viele Patienten zu Herrn Nkeng ?}
\item \all{Welche Heilpflanzen kennt der Heiler ?}
\item \all{Wann schickt Herr Nkeng einen Patienten ins Krankenhaus ?}
\item \all{Was sagt Frau Dr. Mbarga über die traditionelle Medizin ?}
\end{enumerate}

\courssub{Exercice résolu}

\begin{exoresolu}[Complétez avec le bon article et le bon mot]
Complétez : \all{1. \_\_\_ Patient leidet an Fieber. 2. Die Ärztin verschreibt ein \_\_\_ (médicament). 3. \_\_\_ Impfung schützt vor der Krankheit. 4. Er holt sich eine \_\_\_ (rhume). 5. Artemisia ist eine bekannte \_\_\_.}
\tcblower
\textbf{Solution détaillée} :
\begin{enumerate}
\item \all{Der Patient leidet an Fieber.} — nominatif masculin singulier, sujet du verbe \all{leiden}.
\item \all{Die Ärztin verschreibt ein Medikament.} — attention au faux ami : « un médicament » se dit \all{das Medikament} (ou \all{das Heilmittel}), pas \all{die Medizin}.
\item \all{Die Impfung schützt vor der Krankheit.} — les noms en \all{-ung} sont toujours féminins.
\item \all{Er holt sich eine Erkältung.} — expression figée \all{sich eine Erkältung holen} : « attraper un rhume ».
\item \all{Artemisia ist eine bekannte Heilpflanze.} — mot composé : \all{Heil- + Pflanze}, féminin comme \all{die Pflanze}.
\end{enumerate}
\end{exoresolu}

\begin{attention}[Pièges à éviter]
\begin{itemize}
\item Ne dites jamais \all{der Krankheit} : les noms en \all{-heit} sont féminins (\all{die Krankheit}).
\item « Souffrir de » se traduit par \all{leiden an (+ Dativ)} : \all{Er leidet an Fieber}, et non par une construction avec \all{von}.
\item N'oubliez pas la majuscule à tous les noms : \all{das Fieber, die Spritze, das Krankenhaus}.
\item Le pluriel de \all{das Virus} est \all{Viren}, et celui de \all{das Krankenhaus} est \all{Krankenhäuser} avec Umlaut.
\end{itemize}
\end{attention}

\begin{exercice}[À vous de jouer]
\begin{enumerate}
\item Donnez l'article et le pluriel : \all{Arzt, Heilpflanze, Krankenhaus, Apotheke, Medikament, Behandlung}.
\item Traduisez en allemand : « Le tradipraticien soigne le patient avec des plantes médicinales. » / « La vaccination protège les enfants contre la maladie. »
\item Formez le féminin : \all{der Arzt, der Heiler, der Patient, der Krankenpfleger}.
\end{enumerate}
\end{exercice}

\begin{corrige}[Exercice « À vous de jouer »]
\begin{enumerate}
\item \all{der Arzt, Ärzte} ; \all{die Heilpflanze, -n} ; \all{das Krankenhaus, Krankenhäuser} ; \all{die Apotheke, -n} ; \all{das Medikament, -e} ; \all{die Behandlung, -en}.
\item \all{Der Heiler behandelt den Patienten mit Heilpflanzen.} / \all{Die Impfung schützt die Kinder vor der Krankheit.}
\item \all{die Ärztin, die Heilerin, die Patientin, die Krankenpflegerin}.
\end{enumerate}
\end{corrige}

Ce lexique étant désormais structuré, nous pouvons passer à la grammaire : nous verrons dans la section suivante comment conjuguer correctement les verbes pronominaux comme \all{sich erkälten} et \all{sich anstecken}, que vous venez de rencontrer.

\coursec{Grammaire 1 : les verbes pronominaux}

Dans la section précédente, nous avons constitué le vocabulaire de la santé : \all{die Krankheit}, \all{der Arzt}, \all{das Heilmittel}, \all{die Heilpflanze}, \all{die Impfung}, \all{die Pandemie}. Ce lexique prend vie dans des phrases. Or, en relisant notre texte support sur la crise sanitaire et la médecine traditionnelle, deux phrases retiennent l'attention par leur construction particulière :

\begin{center}
\all{Viele Menschen fürchteten sich vor der Pandemie.}\\
« Beaucoup de gens avaient peur de la pandémie. »\\[0.3em]
\all{Unsere Großeltern haben sich mit Heilpflanzen geheilt.}\\
« Nos grands-parents se sont soignés avec des plantes médicinales. »
\end{center}

Que signifie ce petit mot \all{sich} ? Pourquoi ne peut-on pas le supprimer ? C'est toute la question des \cle{verbes pronominaux} (\all{reflexive Verben}), que nous allons étudier maintenant.

\courssub{Observation : à quoi sert le pronom réfléchi ?}

\begin{exemplebox}[Observer avant d'apprendre]
Comparons deux phrases :

\all{Der Arzt wäscht die Hände des Patienten.} — « Le médecin lave les mains du patient. »

\all{Der Arzt wäscht sich.} — « Le médecin se lave (lui-même). »

Dans la première phrase, l'action part du sujet (\all{der Arzt}) et va vers un autre objet (\all{die Hände des Patienten}). Dans la seconde, l'action \emph{revient} sur le sujet : le médecin lave... le médecin. Le pronom \all{sich} marque ce retour de l'action sur son auteur. C'est exactement le rôle du « se » français.
\end{exemplebox}

\begin{definition}[Le verbe pronominal]
Un \cle{verbe pronominal} (\all{ein reflexives Verb}) est un verbe qui se conjugue avec un \cle{pronom réfléchi} (\all{das Reflexivpronomen}) qui renvoie au sujet. À l'infinitif et dans les dictionnaires, on le note avec \all{sich} : \all{sich waschen} (se laver), \all{sich freuen} (se réjouir), \all{sich erkälten} (s'enrhumer).
\end{definition}

\courssub{Les pronoms réfléchis : accusatif et datif}

Le pronom réfléchi varie selon la personne, comme en français (me, te, se, nous, vous, se). Mais l'allemand ajoute une difficulté : il faut choisir entre \cle{accusatif} et \cle{datif}.

\begin{center}
\begin{tabular}{|l|l|l|l|}
\hline
\rowcolor{popblueL} \textbf{Personne} & \textbf{Français} & \textbf{Accusatif} & \textbf{Datif} \\
\hline
1re pers. sg. (ich) & me & mich & mir \\
\hline
2e pers. sg. (du) & te & dich & dir \\
\hline
3e pers. sg. (er/sie/es) & se & sich & sich \\
\hline
1re pers. pl. (wir) & nous & uns & uns \\
\hline
2e pers. pl. (ihr) & vous & euch & euch \\
\hline
3e pers. pl. + Sie & se & sich & sich \\
\hline
\end{tabular}
\end{center}

\begin{aretenir}[Deux choses à retenir immédiatement]
\begin{itemize}
\item Seules les 1re et 2e personnes du singulier distinguent les deux cas : \all{mich/mir}, \all{dich/dir}. Toutes les autres formes sont identiques à l'accusatif et au datif.
\item \all{sich} s'écrit toujours en minuscules à l'intérieur de la phrase ; il ne prend la majuscule que s'il ouvre réellement la phrase, ce qui est rare.
\end{itemize}
\end{aretenir}

\courssub{Accusatif ou datif ? La règle du complément}

C'est LA grande question des verbes pronominaux. La règle est simple et logique.

\begin{propriete}[Choix du cas du pronom réfléchi]
\begin{itemize}
\item Si le verbe n'a \textbf{pas d'autre complément d'objet direct}, le pronom réfléchi est à l'\textbf{accusatif} : \all{Ich wasche mich.} (Je me lave.)
\item Si le verbe a \textbf{déjà un complément d'objet direct} (un COD), le pronom réfléchi passe au \textbf{datif} : \all{Ich wasche mir die Hände.} (Je me lave les mains.) Ici, \all{die Hände} occupe la place de l'accusatif ; le pronom indique alors « pour moi, à moi ».
\end{itemize}
\end{propriete}

\begin{methode}[Choisir entre mich et mir en trois étapes]
\begin{enumerate}
\item Conjugue le verbe avec son sujet : \all{Ich wasche...}
\item Pose la question « quoi ? » (ou « qui ? ») après le verbe. Y a-t-il déjà une réponse dans la phrase ? Ici : oui, \all{die Hände}.
\item Si oui, le pronom est au \textbf{datif} : \all{Ich wasche mir die Hände.} Si non, il est à l'\textbf{accusatif} : \all{Ich wasche mich.}
\end{enumerate}
\end{methode}

\begin{exemplebox}[Accusatif et datif en contexte médical]
\all{Ich fühle mich krank.} — « Je me sens malade. » (pas d'autre COD → accusatif)

\all{Du putzt dir die Zähne.} — « Tu te brosses les dents. » (\all{die Zähne} = COD → datif)

\all{Er bricht sich das Bein.} — « Il se casse la jambe. » (\all{das Bein} = COD → datif)

\all{Wir treffen uns vor dem Krankenhaus.} — « Nous nous retrouvons devant l'hôpital. » (accusatif)

\all{Sie kauft sich ein neues Medikament.} — « Elle s'achète un nouveau médicament. » (\all{ein neues Medikament} = COD → datif)
\end{exemplebox}

\begin{popfigure}
\begin{tikzpicture}[
  node distance=0.9cm,
  startstop/.style={rectangle, rounded corners, draw=popblue, fill=popblueL, align=center, inner sep=6pt},
  decision/.style={diamond, draw=poporange, fill=poporangeL, align=center, inner sep=2pt, aspect=2},
  result/.style={rectangle, rounded corners, draw=popgreen, fill=popgreenL, align=center, inner sep=6pt},
  arr/.style={-{Stealth}, thick, popdark}
]
\node[startstop, align=center] (start){Verbe pronominal\\ \all{sich waschen, sich brechen...}};
\node[decision, below=of start, align=center] (q){Le verbe a-t-il\\ déjà un COD ?};
\node[result, below left=1.2cm and -0.5cm of q, align=center] (oui){OUI $\rightarrow$ pronom au \textbf{datif}\\ \all{Ich wasche \textbf{mir} die Hände.}};
\node[result, below right=1.2cm and -0.5cm of q, align=center] (non){NON $\rightarrow$ pronom à l'\textbf{accusatif}\\ \all{Ich wasche \textbf{mich}.}};
\draw[arr] (start) -- (q);
\draw[arr] (q) -| node[near start, above left]{ja} (oui);
\draw[arr] (q) -| node[near start, above right]{nein} (non);
\end{tikzpicture}
\legende{Organigramme : choisir le cas du pronom réfléchi}
\end{popfigure}

\courssub{Les trois catégories de verbes pronominaux}

\begin{definition}[Trois emplois du pronom réfléchi]
\begin{enumerate}
\item \textbf{Verbes réfléchis} (sens propre) : l'action revient sur le sujet. Le pronom peut être remplacé par un autre complément. \all{Ich wasche mich.} / \all{Ich wasche das Kind.} (Je lave l'enfant.)
\item \textbf{Verbes réciproques} : l'action s'échange entre plusieurs personnes (sens de « l'un l'autre »). \all{Wir treffen uns.} (Nous nous rencontrons.) \all{Sie helfen sich.} (Ils s'entraident.)
\item \textbf{Verbes idiomatiques} (vrais verbes pronominaux) : le pronom est \textbf{obligatoire} et ne peut ni être supprimé ni être remplacé. \all{sich freuen} (se réjouir), \all{sich erkälten} (s'enrhumer), \all{sich interessieren für} (s'intéresser à), \all{sich erinnern an} (se souvenir de), \all{sich beeilen} (se dépêcher), \all{sich fürchten vor} (avoir peur de).
\end{enumerate}
\end{definition}

\begin{exemplebox}[Verbes idiomatiques en contexte]
\all{Sie interessiert sich für Heilpflanzen.} — « Elle s'intéresse aux plantes médicinales. »

\all{Wir freuen uns über die gute Nachricht.} — « Nous nous réjouissons de la bonne nouvelle. »

\all{Der Patient erholt sich schnell.} — « Le patient se rétablit vite. »

\all{Erinnerst du dich an den alten Heiler?} — « Te souviens-tu du vieux guérisseur ? »
\end{exemplebox}

\begin{attention}[Pièges pour les francophones]
\begin{itemize}
\item \textbf{Ne jamais oublier le pronom} avec les verbes idiomatiques : on dit \all{sich freuen}, jamais \all{freuen} seul. \all{Ich freue mich.} et non \all{Ich freue.}
\item \textbf{Le français et l'allemand ne coïncident pas toujours.} Certains verbes pronominaux en français ne le sont pas en allemand : « se lever » = \all{aufstehen} (\all{Ich stehe um sechs Uhr auf.}) ; « se taire » = \all{schweigen}. À l'inverse, \all{sich beeilen} est pronominal alors que le français dit « dépêche-toi » sans pronom visible : \all{Beeil dich!}
\item \textbf{Prépositions à mémoriser avec le verbe} : \all{sich interessieren für} (+ Akk.), \all{sich erinnern an} (+ Akk.), \all{sich freuen über} (+ Akk., joie devant un fait présent ou passé), \all{sich freuen auf} (+ Akk., joie devant un événement à venir), \all{sich fürchten vor} (+ Dat.).
\end{itemize}
\end{attention}

\begin{center}
\begin{tabularx}{\linewidth}{|X|X|X|}
\hline
\rowcolor{popblueL} \textbf{Français} & \textbf{Deutsch} & \textbf{Remarque} \\
\hline
se souvenir de & sich erinnern an (+ Akk.) & pronominal dans les deux langues \\
\hline
s'intéresser à & sich interessieren für (+ Akk.) & préposition \all{für} \\
\hline
se lever & aufstehen & pas pronominal en allemand ! \\
\hline
se dépêcher & sich beeilen & pronominal en allemand seulement \\
\hline
s'enrhumer & sich erkälten & idiomatique, pronom obligatoire \\
\hline
se rétablir & sich erholen & idiomatique \\
\hline
se sentir (bien/mal) & sich fühlen & \all{Ich fühle mich krank.} \\
\hline
\end{tabularx}
\end{center}

\courssub{Conjugaison : sich erkälten au présent et au Perfekt}

\begin{propriete}[Formation et auxiliaire]
Le verbe pronominal se conjugue comme un verbe ordinaire ; le pronom réfléchi s'accorde simplement à la même personne que le sujet. Au \cle{Perfekt}, \textbf{tous} les verbes pronominaux se conjuguent avec l'auxiliaire \all{haben}, sans exception : \all{Er hat sich erkältet.} (Il s'est enrhumé.) C'est une différence majeure avec le français, qui emploie « être » (il s'\emph{est} enrhumé).
\end{propriete}

\begin{center}
\begin{tabular}{|l|l|l|}
\hline
\rowcolor{popblueL} \textbf{Personne} & \textbf{Präsens} & \textbf{Perfekt} \\
\hline
ich & erkälte \textbf{mich} & habe \textbf{mich} erkältet \\
\hline
du & erkältest \textbf{dich} & hast \textbf{dich} erkältet \\
\hline
er/sie/es & erkältet \textbf{sich} & hat \textbf{sich} erkältet \\
\hline
wir & erkälten \textbf{uns} & haben \textbf{uns} erkältet \\
\hline
ihr & erkältet \textbf{euch} & habt \textbf{euch} erkältet \\
\hline
sie/Sie & erkälten \textbf{sich} & haben \textbf{sich} erkältet \\
\hline
\end{tabular}
\end{center}

\begin{popfigure}
\begin{tikzpicture}[node distance=0.55cm and 0.4cm]
\node[draw=popblue, fill=popblueL, rounded corners, align=center, inner sep=6pt] (titre) {\textbf{\all{sich erkälten}} — s'enrhumer};
\node[draw=popgreen, fill=popgreenL, rounded corners, align=center, inner sep=5pt, below left=0.8cm and -1.5cm of titre] (pres) {\textbf{Präsens}\\ \all{ich erkälte mich}\\ \all{du erkältest dich}\\ \all{er erkältet sich}\\ \all{wir erkälten uns}\\ \all{ihr erkältet euch}\\ \all{sie erkälten sich}};
\node[draw=poporange, fill=poporangeL, rounded corners, align=center, inner sep=5pt, below right=0.8cm and -1.5cm of titre] (perf) {\textbf{Perfekt (haben !)}\\ \all{ich habe mich erkältet}\\ \all{du hast dich erkältet}\\ \all{er hat sich erkältet}\\ \all{wir haben uns erkältet}\\ \all{ihr habt euch erkältet}\\ \all{sie haben sich erkältet}};
\draw[-{Stealth}, thick, popdark] (titre) -- (pres);
\draw[-{Stealth}, thick, popdark] (titre) -- (perf);
\end{tikzpicture}
\legende{Carte de conjugaison : sich erkälten au présent et au Perfekt}
\end{popfigure}

\courssub{La place du pronom dans la phrase}

\begin{propriete}[Position du pronom réfléchi]
\begin{itemize}
\item En phrase déclarative, le pronom se place \textbf{juste après le verbe conjugué} : \all{Ich erkälte mich oft im Regen.} (Je m'enrhume souvent sous la pluie.)
\item En cas d'inversion (le sujet après le verbe), le pronom se place \textbf{avant un sujet nominal} : \all{Im Winter erkältet sich mein Bruder oft.} (En hiver, mon frère s'enrhume souvent.) — mais \textbf{après un pronom personnel sujet} : \all{Im Winter erkältet er sich oft.}
\item À l'impératif : \all{Ruh dich aus!} (Repose-toi !) ; \all{Beeilt euch!} (Dépêchez-vous !) ; \all{Erinnern Sie sich an den Arzt!} (Souvenez-vous du médecin !)
\item Au Perfekt, le pronom reste près de l'auxiliaire, le participe va en fin de phrase : \all{Sie hat sich für Heilpflanzen interessiert.}
\item Après un verbe modal, le pronom précède l'infinitif renvoyé en fin de phrase : \all{Du musst dich ausruhen.} (Tu dois te reposer.)
\end{itemize}
\end{propriete}

\courssub{Un mini-dialogue à l'infirmerie du lycée}

\begin{textebox}[Beim Schulkrankenpfleger — Chez l'infirmier scolaire]
„Guten Morgen, Chantal! Wie fühlst du dich heute?“

„Guten Morgen, Herr Krankenpfleger. Ich fühle mich nicht gut. Ich habe mich gestern auf dem Sportplatz erkältet. Mein Kopf tut weh, und ich habe die ganze Nacht gehustet.“

„Hast du dich denn nicht warm angezogen? Es regnet viel in dieser Jahreszeit in Yaoundé.“

„Doch, aber ich habe mich beeilt und meine Jacke vergessen. Meine Mutter ärgert sich immer, wenn ich mich so unvorsichtig verhalte.“

„Dann solltest du dich jetzt ausruhen und viel Tee trinken. Meine Großmutter bereitet sich immer einen Tee aus Ingwer und Zitrone zu, wenn sie sich erkältet hat. Diese Heilpflanzen helfen wirklich!“

„Meine Großeltern vertrauen auch der traditionellen Medizin. Sie kennen sich gut mit Heilpflanzen aus.“

„Das ist gut. Aber wenn das Fieber steigt, musst du dich von einem Arzt untersuchen lassen. Und vergiss nicht: Erhol dich gut, bevor du in die Klasse zurückkehrst!“

„Danke, Herr Krankenpfleger. Ich werde mich an Ihren Rat halten.“
\end{textebox}

\textbf{Glossaire :} \all{der Schulkrankenpfleger, -} (l'infirmier scolaire) ; \all{husten} (tousser) ; \all{sich warm anziehen} (s'habiller chaudement) ; \all{die Jahreszeit, -en} (la saison) ; \all{sich ärgern} (se fâcher) ; \all{sich verhalten} (se comporter) ; \all{sich ausruhen} (se reposer) ; \all{der Ingwer} (le gingembre) ; \all{sich auskennen mit} (s'y connaître en) ; \all{das Fieber, -} (la fièvre) ; \all{sich untersuchen lassen} (se faire examiner) ; \all{der Rat} (le conseil) ; \all{sich an den Rat halten} (suivre le conseil).

\textbf{Questions de compréhension :} 1. Warum fühlt sich Chantal nicht gut? 2. Was empfiehlt der Krankenpfleger? 3. Welche Rolle spielen die Großeltern und die traditionelle Medizin im Dialog?

\courssub{Exercice résolu et entraînement}

\begin{exoresolu}[Conjuguer et choisir le bon pronom]
Énoncé — Complète avec le pronom réfléchi correct (accusatif ou datif) :

1. Ich wasche \ldots{} die Hände vor dem Essen. 2. Wir freuen \ldots{} über die Impfung. 3. Du musst \ldots{} beeilen, der Arzt wartet! 4. Die Kinder haben \ldots{} im Regen erkältet. 5. Interessierst du \ldots{} für traditionelle Medizin?

\tcblower
\textbf{Solution détaillée :}

1. \all{Ich wasche \textbf{mir} die Hände.} — Il y a déjà un COD (\all{die Hände}) : datif.

2. \all{Wir freuen \textbf{uns}.} — Verbe idiomatique, pas d'autre COD : accusatif.

3. \all{Du musst \textbf{dich} beeilen.} — \all{sich beeilen} est idiomatique ; après un modal, le pronom précède l'infinitif renvoyé en fin de phrase : accusatif.

4. \all{Die Kinder haben \textbf{sich} erkältet.} — Perfekt avec \all{haben}, pronom accusatif, 3e personne du pluriel.

5. \all{Interessierst du \textbf{dich} für traditionelle Medizin?} — \all{sich interessieren für} : pronom obligatoire, accusatif.
\end{exoresolu}

\begin{exercice}[À ton tour]
1. Traduis en allemand : a) Je me sens malade. b) Nous nous sommes reposés après la crise. c) Elle se lave les mains avec du savon. d) Souvenez-vous du médecin traditionnel !

2. Conjugue \all{sich erholen} (se rétablir) au présent et au Perfekt à toutes les personnes.

3. Trouve l'erreur et corrige : \all{Ich habe mich die Zähne geputzt.} / \all{Er freut über die Nachricht.}
\end{exercice}

\begin{corrige}[Exercice « À ton tour »]
1. a) \all{Ich fühle mich krank.} b) \all{Wir haben uns nach der Krise ausgeruht.} c) \all{Sie wäscht sich die Hände mit Seife.} d) \all{Erinnern Sie sich an den traditionellen Heiler!}

2. Présent : \all{ich erhole mich, du erholst dich, er erholt sich, wir erholen uns, ihr erholt euch, sie erholen sich.} Perfekt : \all{ich habe mich erholt, du hast dich erholt, er hat sich erholt, wir haben uns erholt, ihr habt euch erholt, sie haben sich erholt.}

3. \all{Ich habe \textbf{mir} die Zähne geputzt} (COD présent → datif) ; \all{Er freut \textbf{sich} über die Nachricht} (verbe idiomatique : le pronom est obligatoire).
\end{corrige}

\begin{aretenir}[L'essentiel des verbes pronominaux]
\begin{itemize}
\item Verbe conjugué + pronom réfléchi à la même personne que le sujet : \all{ich... mich/mir}, \all{du... dich/dir}, \all{er... sich}.
\item Déjà un COD dans la phrase ? → pronom au datif : \all{Ich wasche mir die Hände.}
\item Au Perfekt : toujours \all{haben} : \all{Er hat sich erkältet.}
\item Les verbes idiomatiques (\all{sich freuen, sich erinnern an, sich beeilen}) perdent leur sens sans le pronom : ne l'oublie jamais.
\end{itemize}
\end{aretenir}

Ces verbes pronominaux nous seront immédiatement utiles dans la section suivante : pour donner un conseil à un malade — \all{Du solltest dich ausruhen!} —, il nous faudra combiner le pronom réfléchi et le subjonctif II, que nous allons maintenant étudier.

\coursec{Grammaire 2 : le conseil avec le subjonctif II}

Dans la section précédente, nous avons étudié les \cle{verbes pronominaux} (\all{sich ausruhen}, \all{sich erkälten}, \all{sich impfen lassen}). Or, dans notre texte support, ces verbes apparaissaient très souvent accompagnés d'un verbe modal à une forme particulière : \all{Man sollte sich impfen lassen.} Cette forme, \all{sollte}, n'est pas un simple passé : c'est le \cle{subjonctif II} (\all{der Konjunktiv II}), le mode du conseil, de la recommandation et de l'hypothèse. C'est précisément ce que nous allons étudier maintenant : comment, en allemand, donner un conseil poli et nuancé — une compétence exigée au Baccalauréat, notamment dans les productions écrites sur la santé.

\courssub{Observation : que remarque-t-on dans le texte ?}

Reprenons deux phrases tirées de notre texte sur la crise sanitaire et la médecine traditionnelle :

\begin{exemplebox}[Deux phrases du texte support]
\all{„Man sollte sich impfen lassen.“} « On devrait se faire vacciner. »

\all{„Dann solltest du dich ausruhen!“} « Alors tu devrais te reposer ! »
\end{exemplebox}

Observons attentivement :

\begin{itemize}
\item Le verbe \all{sollen} apparaît sous la forme \all{sollte} / \all{solltest} : ce n'est ni le présent (\all{soll}, \all{sollst}) ni un simple prétérit narratif. C'est la forme du \cle{subjonctif II}.
\item Le sens n'est pas « tu devais » (passé), mais « tu \textbf{devrais} » : le locuteur donne un \cle{conseil}, une \cle{recommandation}, il n'affirme pas un fait.
\item L'infinitif (\all{impfen lassen}, \all{ausruhen}) se trouve \textbf{en fin de phrase}, comme toujours avec un verbe modal.
\item Dans la deuxième phrase, le verbe pronominal conserve son pronom réfléchi \all{dich}, placé juste après le sujet.
\end{itemize}

\begin{definition}[Le subjonctif II (Konjunktiv II)]
Le \cle{subjonctif II} est le mode allemand de l'\textbf{irréel}, du \textbf{souhait}, de l'\textbf{hypothèse} et du \textbf{conseil poli}. Il correspond le plus souvent au \textbf{conditionnel français} : \all{du solltest} = « tu devrais », \all{ich würde gehen} = « j'irais ». Dans cette leçon, nous nous limitons à son emploi pour \textbf{conseiller} et \textbf{recommander}.
\end{definition}

\courssub{Formation du subjonctif II des verbes modaux}

Bonne nouvelle : pour les verbes modaux, le subjonctif II est très régulier. On part du \textbf{prétérit} du modal, et on ajoute un \textbf{tréma} (Umlaut) sur la voyelle radicale \textbf{a, o, u} quand c'est possible.

\begin{propriete}[Formation du Konjunktiv II des modaux]
\textbf{Règle :} Konjunktiv II = prétérit du modal + tréma sur a, o, u.

\begin{itemize}
\item \all{können} → prétérit \all{konnte} → KII \all{könnte} (tréma sur o)
\item \all{müssen} → prétérit \all{musste} → KII \all{müsste} (tréma sur u)
\item \all{dürfen} → prétérit \all{durfte} → KII \all{dürfte} (tréma sur u)
\item \all{sollen} → prétérit \all{sollte} → KII \all{sollte} (\textbf{pas de tréma} : la forme KII est identique au prétérit)
\item \all{wollen} → prétérit \all{wollte} → KII \all{wollte} (\textbf{pas de tréma} non plus)
\end{itemize}

\textbf{Exception à retenir :} \all{sollen} et \all{wollen} ne prennent \textbf{jamais} de tréma. Seuls le contexte et le sens permettent de distinguer \all{sollte} « devrait » (conseil) de \all{sollte} « devait » (passé).
\end{propriete}

\begin{popfigure}
\begin{center}
\begin{tikzpicture}[
  box/.style={draw=popblue, fill=popblueL, rounded corners=2pt, align=center, font=\small, minimum height=0.9cm},
  kii/.style={draw=poporange, fill=poporangeL, rounded corners=2pt, align=center, font=\small\bfseries, minimum height=0.9cm},
  fleche/.style={-{Stealth[length=2.5mm]}, thick, popdark}]
\node[box] (s1) at (0,0) {sollen};
\node[kii, right=1.6cm of s1] (s2) {sollte};
\draw[fleche] (s1) -- (s2);
\node[box] (m1) at (0,-1.2) {müssen};
\node[kii, right=1.6cm of m1] (m2) {müsste};
\draw[fleche] (m1) -- (m2);
\node[box] (k1) at (0,-2.4) {können};
\node[kii, right=1.6cm of k1] (k2) {könnte};
\draw[fleche] (k1) -- (k2);
\node[box] (d1) at (0,-3.6) {dürfen};
\node[kii, right=1.6cm of d1] (d2) {dürfte};
\draw[fleche] (d1) -- (d2);
\node[box] (w1) at (0,-4.8) {wollen};
\node[kii, right=1.6cm of w1] (w2) {wollte};
\draw[fleche] (w1) -- (w2);
\node[align=center, font=\footnotesize\itshape, popmuted] at (3.2,-5.8) {prétérit + tréma (sauf sollen / wollen)};
\end{tikzpicture}
\end{center}
\legende{Les verbes modaux au subjonctif II : du présent à la forme de conseil.}
\end{popfigure}

Voici la conjugaison complète des quatre modaux les plus utiles pour le conseil, au subjonctif II :

\begin{center}
\begin{tabular}{|l|l|l|l|l|}
\hline
\rowcolor{popblueL} Personne & sollen & müssen & können & dürfen \\
\hline
ich & sollte & müsste & könnte & dürfte \\
\hline
du & solltest & müsstest & könntest & dürftest \\
\hline
er/sie/es/man & sollte & müsste & könnte & dürfte \\
\hline
wir & sollten & müssten & könnten & dürften \\
\hline
ihr & solltet & müsstet & könntet & dürftet \\
\hline
sie/Sie & sollten & müssten & könnten & dürften \\
\hline
\end{tabular}
\end{center}

\begin{aretenir}[L'essentiel sur la formation]
\begin{itemize}
\item Les terminaisons sont celles du prétérit des modaux : \all{-e, -est, -e, -en, -et, -en}.
\item Tréma obligatoire pour \all{können}, \all{müssen}, \all{dürfen} ; jamais pour \all{sollen} et \all{wollen}.
\item Pour le conseil, les trois modaux rois sont : \all{sollte} (devrait), \all{könnte} (pourrait), \all{müsste} (devrait vraiment / il faudrait).
\end{itemize}
\end{aretenir}

\courssub{Structure de la phrase de conseil}

La structure est exactement celle de la phrase avec modal que vous connaissez depuis la classe de Seconde : le modal conjugué occupe la \textbf{deuxième position}, et l'infinitif se place \textbf{en toute fin de phrase}.

\begin{popfigure}
\begin{center}
\begin{tikzpicture}[
  bloc/.style={draw=popdark, rounded corners=3pt, align=center, font=\small, minimum height=1cm, inner sep=5pt},
  fleche/.style={-{Stealth[length=2.5mm]}, thick, poporange}]
\node[bloc, fill=popblueL, align=center] (suj) at (0,0){1. Sujet\\Du};
\node[bloc, fill=poporangeL, right=0.5cm of suj, align=center] (mod){2. Modal KII\\solltest};
\node[bloc, fill=poppurpleL, right=0.5cm of mod, align=center] (mil){3. Compléments\\mehr Tee};
\node[bloc, fill=popgreenL, right=0.5cm of mil, align=center] (inf){4. Infinitif\\trinken.};
\draw[fleche] (suj) -- (mod);
\draw[fleche] (mod) -- (mil);
\draw[fleche] (mil) -- (inf);
\end{tikzpicture}
\end{center}
\legende{Structure de la phrase de conseil : sujet + modal au KII + compléments + infinitif final.}
\end{popfigure}

\begin{exemplebox}[Exemples traduits]
\all{Du solltest zum Arzt gehen.} « Tu devrais aller chez le médecin. »

\all{Man könnte Heilpflanzen und Medikamente verbinden.} « On pourrait combiner plantes médicinales et médicaments. »

\all{Wir müssten mehr über traditionelle Heilmittel lernen.} « Nous devrions apprendre davantage sur les remèdes traditionnels. »

\all{Ihr könntet euch besser vor den Krankheiten schützen.} « Vous pourriez mieux vous protéger contre les maladies. »

\all{Es wäre besser, wenn du dich impfen ließest.} « Il vaudrait mieux que tu te fasses vacciner. »
\end{exemplebox}

\begin{propriete}[Cas particulier : le conseil avec un verbe pronominal]
Quand le conseil porte sur un verbe pronominal, le \textbf{pronom réfléchi} reste à sa place habituelle (après le sujet ou après le verbe conjugué), et l'infinitif pronominal va en fin de phrase :

\all{Du solltest \textbf{dich} ausruhen.} « Tu devrais te reposer. »

\all{Man sollte \textbf{sich} impfen lassen.} « On devrait se faire vacciner. »

\all{Ihr müsstet \textbf{euch} mehr bewegen.} « Vous devriez bouger davantage. »

Le pronom s'accorde toujours avec le sujet : \all{ich — mich}, \all{du — dich}, \all{er/sie/es/man — sich}, \all{wir — uns}, \all{ihr — euch}, \all{sie/Sie — sich}.
\end{propriete}

\courssub{Les nuances du conseil : sollte, könnte, müsste}

Les trois modaux ne disent pas la même chose. Choisir le bon modal, c'est doser la force de son conseil — exactement comme en français entre « tu devrais », « tu pourrais » et « il faudrait vraiment que ».

\begin{center}
\begin{tabularx}{\linewidth}{|l|X|X|X|}
\hline
\rowcolor{popblueL} Modal KII & Français & Nuance & Exemple \\
\hline
\all{sollte} & devrait & conseil ferme, direct & \all{Du solltest zum Arzt gehen.} \\
\hline
\all{könnte} & pourrait & suggestion ouverte, polie & \all{Du könntest mehr Tee trinken.} \\
\hline
\all{müsste} & il faudrait (vraiment) & obligation morale, insistance & \all{Du müsstest dich endlich impfen lassen!} \\
\hline
\all{dürfte} & il est permis / probable & hypothèse prudente & \all{Das dürfte kein Problem sein.} \\
\hline
\end{tabularx}
\end{center}

\courssub{Autres tournures de conseil à connaître}

Pour enrichir vos productions écrites au Baccalauréat, voici trois tournures très appréciées des correcteurs :

\begin{exemplebox}[Tournures de conseil de niveau Terminale]
\textbf{1.} \all{Es wäre besser, wenn + KII} = « Il vaudrait mieux que... »

\all{Es wäre besser, wenn du mehr schliefest.} « Il vaudrait mieux que tu dormes davantage. »

\textbf{2.} \all{An deiner Stelle würde ich...} = « À ta place, je... »

\all{An deiner Stelle würde ich einen Arzt besuchen.} « À ta place, je consulterais un médecin. »

\textbf{3.} \all{Ich würde dir raten, ... zu + infinitif} = « Je te conseillerais de... »

\all{Ich würde dir raten, dich impfen zu lassen.} « Je te conseillerais de te faire vacciner. »
\end{exemplebox}

\begin{center}
\begin{tabularx}{\linewidth}{|X|X|X|}
\hline
\rowcolor{popblueL} Français & Deutsch & Remarque \\
\hline
Tu devrais te reposer. & \all{Du solltest dich ausruhen.} & modal seul au KII, sans \all{würde} \\
\hline
Il vaudrait mieux que tu consultes. & \all{Es wäre besser, wenn du zum Arzt gingest.} & \all{wenn} + KII en fin de proposition \\
\hline
À ta place, je me ferais vacciner. & \all{An deiner Stelle würde ich mich impfen lassen.} & \all{würde} + infinitif \\
\hline
On pourrait combiner les deux médecines. & \all{Man könnte beide Medizinen verbinden.} & suggestion ouverte \\
\hline
\end{tabularx}
\end{center}

\courssub{Un dialogue pour réinvestir : chez le médecin à Yaoundé}

\begin{textebox}[Beim Arzt in Jaunde — dialogue]
\textbf{Frau Mbarga:} Guten Tag, Herr Doktor! Ich fühle mich seit drei Tagen krank. Ich habe Fieber und Kopfschmerzen.

\textbf{Arzt:} Guten Tag! Sie sollten sich zuerst untersuchen lassen. Haben Sie sich impfen lassen?

\textbf{Frau Mbarga:} Nein, leider nicht. Meine Großmutter sagt, man könnte alles mit Heilpflanzen heilen. Sie gibt mir immer einen Tee aus Rinde und Blättern.

\textbf{Arzt:} Traditionelle Heilmittel können manchmal helfen, das ist wahr. Aber bei einer schweren Krankheit sollte man vorsichtig sein. Es wäre besser, wenn Sie die traditionelle und die moderne Medizin verbinden würden.

\textbf{Frau Mbarga:} Was sollte ich also tun?

\textbf{Arzt:} Sie sollten sich ausruhen, viel Wasser trinken und diese Tabletten nehmen. An Ihrer Stelle würde ich auch einen Bluttest machen. Und Sie müssten sich wirklich impfen lassen — das ist der beste Schutz!

\textbf{Frau Mbarga:} Danke, Herr Doktor. Ich würde meiner Großmutter auch raten, öfter ins Krankenhaus zu kommen.

\textbf{Arzt:} Gute Idee! Man könnte in Ihrem Viertel eine Informationswoche über Impfungen organisieren. Gute Besserung!

\emph{Glossaire :} das Fieber (–) : la fièvre ; der Kopfschmerz (–en) : le mal de tête ; sich untersuchen lassen : se faire examiner ; die Heilpflanze (–n) : la plante médicinale ; die Rinde (–n) : l'écorce ; das Blatt (¨er) : la feuille ; heilen : guérir ; vorsichtig : prudent ; die Tablette (–n) : le comprimé ; der Bluttest (–s) : la prise de sang ; der Schutz (–) : la protection ; das Viertel (–) : le quartier ; Gute Besserung! : Bon rétablissement !
\end{textebox}

\textbf{Questions de compréhension :}

\begin{enumerate}
\item Quels symptômes Madame Mbarga présente-t-elle ?
\item Que pense la grand-mère de la médecine traditionnelle ? Relevez la phrase avec un modal au KII.
\item Quels conseils le médecin donne-t-il ? Relevez trois tournures de conseil différentes.
\item Quelle proposition finale le médecin fait-il pour le quartier ?
\end{enumerate}

\courssub{Exercice résolu et pièges à éviter}

\begin{exoresolu}[Conseils à la pharmacie]
\textbf{Énoncé.} Traduisez en allemand, en utilisant le modal indiqué entre parenthèses :

1. Tu devrais boire plus d'eau. (sollen)

2. On pourrait aussi utiliser des plantes médicinales. (können)

3. Il faudrait vraiment que vous vous fassiez vacciner. (müssen)

4. Il vaudrait mieux que tu ailles chez le médecin. (\all{es wäre besser, wenn})
\tcblower
\textbf{Solution détaillée.}

1. \all{Du solltest mehr Wasser trinken.} — Modal \all{sollen} au KII, 2\textsuperscript{e} personne : \all{solltest}. Infinitif \all{trinken} en fin de phrase. « Plus d'eau » = \all{mehr Wasser} (sans article).

2. \all{Man könnte auch Heilpflanzen benutzen.} — « On » se traduit par \all{man} ; tréma obligatoire : \all{könnte}. « Plantes médicinales » = \all{die Heilpflanzen} (pluriel).

3. \all{Sie müssten sich wirklich impfen lassen.} — Verbe pronominal \all{sich impfen lassen} : le pronom \all{sich} suit le sujet, l'infinitif double \all{impfen lassen} ferme la phrase. \all{müssten} avec tréma.

4. \all{Es wäre besser, wenn du zum Arzt gingest.} (ou : \all{... wenn du zum Arzt gehen würdest.}) — Après \all{wenn}, le verbe conjugué se place en fin de subordonnée. Le KII de \all{gehen} est \all{gingest} (forme soutenue) ; la tournure avec \all{würde} + infinitif est acceptée et plus fréquente à l'oral.
\end{exoresolu}

\begin{attention}[L'erreur classique du francophone]
Ne traduisez \textbf{jamais} « tu devrais » par \all{du würdest sollen} ! C'est un calque du conditionnel français (« tu \textbf{devr}\textbf{ais} » perçu comme « tu ferais devoir »). En allemand, \textbf{le modal seul au subjonctif II suffit} :

« Tu devrais aller chez le médecin. » = \all{Du solltest zum Arzt gehen.} (et non \all{Du würdest zum Arzt gehen sollen.})

Autres pièges fréquents :

\begin{itemize}
\item Oublier le tréma : écrire \all{du konntest} au lieu de \all{du könntest} change le sens (passé réel au lieu du conseil).
\item Placer l'infinitif trop tôt : \all{Du solltest gehen zum Arzt} est fautif ; l'infinitif ferme la phrase : \all{Du solltest zum Arzt gehen.}
\item Oublier le pronom réfléchi : \all{Du solltest impfen lassen} est incorrect ; il faut \all{Du solltest \textbf{dich} impfen lassen.}
\item Confondre \all{sollte} (devrait, conseil) et \all{wollte} (voulait / voudrait) : \all{Ich wollte} au KII signifie « je voudrais », pas « je devrais ».
\end{itemize}
\end{attention}

\begin{methode}[Donner un conseil au Baccalauréat en quatre étapes]
\begin{enumerate}
\item \textbf{Identifier le problème} dans le sujet ou le texte : \all{Das Problem: Er hat Fieber und geht nicht zum Arzt.}
\item \textbf{Choisir le modal} selon la nuance : \all{sollte} pour un conseil direct, \all{könnte} pour une suggestion polie, \all{müsste} pour insister.
\item \textbf{Construire la phrase} : sujet + modal au KII + compléments + infinitif en fin de phrase : \all{Du solltest sofort zum Arzt gehen.}
\item \textbf{Ajouter une justification} avec \all{weil} (verbe en fin de subordonnée) : \all{Du solltest sofort zum Arzt gehen, weil Fieber gefährlich sein kann.} « Tu devrais aller tout de suite chez le médecin, parce que la fièvre peut être dangereuse. »
\end{enumerate}
\end{methode}

\begin{savaistu}[« Sich impfen lassen » : la construction avec lassen]
Dans \all{Du solltest dich impfen lassen}, le verbe \all{lassen} a un sens particulier : il exprime le fait de \textbf{faire faire quelque chose par quelqu'un d'autre} — c'est le « faire » causatif français. \all{Ich lasse mich impfen.} = « Je me fais vacciner » (c'est le médecin qui vaccine). Autres exemples utiles : \all{Ich lasse mir die Haare schneiden.} « Je me fais couper les cheveux » ; \all{Er lässt sein Fahrrad reparieren.} « Il fait réparer son vélo ». Cette construction, combinée au subjonctif II (\all{Du solltest dich untersuchen lassen} — « Tu devrais te faire examiner »), est très valorisée dans les copies de Terminale.
\end{savaistu}

\begin{exercice}[À vous de conseiller]
1. Conjuguez le modal au subjonctif II : a) \all{Ihr} (müssen) \all{mehr Sport machen.} b) \all{Man} (können) \all{Heilpflanzen im Garten pflanzen.} c) \all{Wir} (sollen) \all{uns gesund ernähren.}

2. Traduisez : a) Tu devrais te reposer, parce que tu es malade. b) On pourrait organiser une semaine de la santé au lycée. c) À ta place, je me ferais vacciner.

3. Production : votre ami(e) refuse la vaccination et ne croit qu'aux remèdes traditionnels. Écrivez-lui trois conseils nuancés (un avec \all{sollte}, un avec \all{könnte}, un avec \all{es wäre besser, wenn}).
\end{exercice}

\begin{corrige}[Exercice : À vous de conseiller]
1. a) \all{Ihr müsstet mehr Sport machen.} b) \all{Man könnte Heilpflanzen im Garten pflanzen.} c) \all{Wir sollten uns gesund ernähren.}

2. a) \all{Du solltest dich ausruhen, weil du krank bist.} b) \all{Man könnte eine Gesundheitswoche am Gymnasium organisieren.} c) \all{An deiner Stelle würde ich mich impfen lassen.}

3. Exemple de production : \all{Du solltest dich informieren, weil Impfungen viele Leben retten. Du könntest auch mit einem Arzt über deine Angst sprechen. Es wäre besser, wenn du die traditionelle Medizin und die moderne Medizin verbinden würdest.}
\end{corrige}

\coursec{Méthode du commentaire de texte et commentaire modèle}

Après avoir maîtrisé le \cle{subjonctif II} pour formuler des conseils (\all{Man sollte...}, \all{Ich würde Ihnen raten...}), nous abordons maintenant l'exercice phare de l'écrit au Baccalauréat : le \cle{commentaire de texte}. Cet exercice mobilise toutes vos compétences : compréhension fine, analyse lexicale et grammaticale, structuration logique et expression personnelle argumentée. Le thème de notre leçon — la crise sanitaire et la médecine traditionnelle — offre un terrain idéal pour confronter des arguments, nuancer son propos et utiliser le \all{Konjunktiv II} dans la conclusion pour émettre des hypothèses ou des recommandations.

\courssub{La structure en trois parties du commentaire}

Un commentaire de texte réussi en allemand suit une architecture rigoureuse en \cle{trois parties} : \all{Einleitung} (introduction), \all{Hauptteil} (développement) et \all{Schluss} (conclusion). Chaque partie a une fonction précise et répond à des codes formels stricts.

\begin{definition}[Les trois piliers du commentaire]
\begin{itemize}
    \item \textbf{Einleitung (Introduction) :} Présenter le document (titre, auteur, source, date, type de texte), annoncer le \cle{thème} général, formuler la \cle{problématique} (souvent une question ouverte) et annoncer le \cle{plan} (souvent en deux ou trois arguments).
    \item \textbf{Hauptteil (Développement) :} Analyser le texte en suivant le plan annoncé. Chaque paragraphe traite un argument : \textbf{Idée directrice} (Thesen Satz) $\rightarrow$ \textbf{Preuve textuelle} (citation courte entre guillemets allemands) $\rightarrow$ \textbf{Analyse/Explication} (lien avec le thème, vocabulaire, grammaire, contexte).
    \item \textbf{Schluss (Conclusion) :} Faire le \cle{bilan} (réponse synthétique à la problématique), ouvrir sur une \cle{réflexion personnelle} (avis argumenté, perspectives) et éventuellement lancer une \cle{ouverture} vers un autre sujet ou l'actualité.
\end{itemize}
\end{definition}

\begin{center}
\begin{tabularx}{\linewidth}{|X|X|X|}
\hline
\rowcolor{popblueL} \textbf{Partie} & \textbf{Contenu clé (Deutsch)} & \textbf{Expressions-types (Redemittel)} \\
\hline
Einleitung & Textvorstellung, Thema, Problemfrage, Gliederung & \all{Der Text ... handelt von ... ; Es geht um die Frage, ob ... ; Zunächst ... , dann ... , schließlich ...} \\
\hline
Hauptteil & Argument 1, 2, 3 : These – Beleg – Analyse & \all{Der Autor argumentiert, dass ... ; Dies wird deutlich durch das Zitat „...“ ; Das zeigt, dass ... ; Ein weiteres Argument ist ...} \\
\hline
Schluss & Zusammenfassung, eigene Meinung, Ausblick & \all{Zusammenfassend lässt sich sagen ... ; Meiner Meinung nach ... ; Man könnte folgern, dass ... ; Es bleibt zu hoffen, dass ...} \\
\hline
\end{tabularx}
\end{center}

\begin{popfigure}
\begin{tikzpicture}[
    node distance=1.2cm and 2cm,
    every node/.style={font=\small, align=center},
    box/.style={draw=popblue, thick, rounded corners, fill=popblueL, minimum width=3.2cm, minimum height=1.8cm},
    arrow/.style={-Stealth, thick, popdark}
]
\node[box, align=center] (einl){\textbf{Einleitung}\\\textit{(Présentation)}\\Thème + Problématique + Plan};
\node[box, below=of einl, align=center] (haupt){\textbf{Hauptteil}\\\textit{(Analyse)}\\Arg. 1 : These – Beleg – Analyse\\Arg. 2 : These – Beleg – Analyse\\(Arg. 3)};
\node[box, below=of haupt, align=center] (schl){\textbf{Schluss}\\\textit{(Bilan + Opinion)}\\Réponse + Avis personnel + Ouverture};

\draw[arrow] (einl) -- (haupt);
\draw[arrow] (haupt) -- (schl);

% Side annotations
\node[left=0.5cm of einl, text width=2.5cm, align=right, popmuted] {10\% du texte};
\node[left=0.5cm of haupt, text width=2.5cm, align=right, popmuted] {70\% du texte};
\node[left=0.5cm of schl, text width=2.5cm, align=right, popmuted] {20\% du texte};
\end{tikzpicture}
\legende{Organigramme de la structure du commentaire : proportions indicatives et enchaînement logique.}
\end{popfigure}

\courssub{La problématique et l'introduction (Einleitung)}

La problématique est le \cle{moteur} de votre commentaire. Elle ne doit pas être une question fermée (oui/non) mais une question d'analyse ou d'évaluation. Pour notre thème, une problématique type est :

\begin{center}
\cle{Kann die traditionelle Medizin die moderne Medizin ergänzen?}
\end{center}
(La médecine traditionnelle peut-elle compléter la médecine moderne ?)

\begin{methode}[Rédiger l'Einleitung en 4 étapes]
\begin{enumerate}
    \item \textbf{Accroche contextuelle :} Situer le thème (crise sanitaire, Cameroun, Allemagne, etc.).
    \item \textbf{Présentation du texte :} Titre (entre guillemets), auteur, source, date, type (article, interview, reportage).
    \item \textbf{Problématique :} La question directrice (souvent introduite par \all{Es geht um die Frage, ob...} ou \all{Zentrale Frage ist...}).
    \item \textbf{Annonce du plan :} \all{Der Text gliedert sich in drei Teile: Zunächst ... , dann ... , schließlich ...}
\end{enumerate}
\end{methode}

\begin{textebox}[Modèle d'Einleitung — Texte support fictif : „Zwischen Impfung und Heilpflanze“]
\all{Der Zeitungsartikel „Zwischen Impfung und Heilpflanze“, veröffentlicht am 12. Oktober 2023 in der „Cameroon Tribune“, behandelt die Gesundheitskrise in Kamerun und die Frage, ob traditionelle und moderne Medizin zusammenarbeiten können. Der Autor schildert die Spannungen zwischen dem Vertrauen in Heilpflanzen und der schulmedizinischen Impfkampagne. Zentrale Problemfrage ist: Kann die traditionelle Medizin die moderne Medizin ergänzen? Der Text gliedert sich in drei Teile: Zunächst wird die Position der Heilerin Mama Ngo vorgestellt, dann die Argumentation des Arztes Dr. Mbarga dargestellt, schließlich werden Möglichkeiten der Kooperation skizziert.}
\end{textebox}

\begin{exemplebox}[Glossaire de l'introduction]
\begin{itemize}
    \item \all{der Zeitungsartikel, die -n} : l'article de presse
    \item \all{veröffentlichen} : publier (Partizip II : \all{veröffentlicht})
    \item \all{behandeln} : traiter (un sujet) — \all{handeln von} (+ Datif) : traiter de
    \item \all{die Gesundheitskrise, -n} : la crise sanitaire
    \item \all{die Heilerin, -nen} : la guérisseuse / la tradipraticienne
    \item \all{die Impfkampagne, -n} : la campagne de vaccination
    \item \all{die Spannung, -en} : la tension
    \item \all{skizzieren} : esquisser, décrire en gros traits
\end{itemize}
\end{exemplebox}

\courssub{Le développement : analyse et argumentation (Hauptteil)}

Le \all{Hauptteil} est le cœur du commentaire. Il ne s'agit \textbf{pas} de résumer le texte paragraphe par paragraphe, mais de le \cle{décortiquer} en suivant votre plan. Chaque argument suit le schéma : \textbf{These – Beleg – Analyse}.

\begin{propriete}[Les connecteurs logiques pour structurer le Hauptteil]
Utilisez des \cle{Konjunktionen} et \cle{Adverbien} pour lier vos arguments. Attention à l'ordre des mots (verbe en 2e position après adverbe de liaison).
\end{propriete}

\begin{center}
\begin{tabularx}{\linewidth}{|l|X|X|}
\hline
\rowcolor{popgreenL} \textbf{Fonction} & \textbf{Allemand (Deutsch)} & \textbf{Français} \\
\hline
Énumérer (Plan) & \all{zunächst / erstens / zuerst ; dann / zweitens / als Nächstes ; schließlich / drittens / zum Schluss} & Premièrement ; Deuxièmement ; Finalement \\
\hline
Ajouter / Renforcer & \all{außerdem ; darüber hinaus ; zudem ; auch ; nicht nur ... sondern auch} & De plus ; En outre ; Aussi ; Non seulement ... mais aussi \\
\hline
Contraster / Opposer & \all{aber ; jedoch ; andererseits ; im Gegensatz dazu ; während ; obwohl} & Mais ; Cependant ; D'un autre côté ; En revanche ; Alors que ; Bien que \\
\hline
Illustrer / Citer & \all{beispielsweise ; zum Beispiel ; dies zeigt sich daran, dass ... ; ein Beleg dafür ist ...} & Par exemple ; Cela se voit par le fait que... ; Une preuve en est... \\
\hline
Conclure un paragraphe & \all{daher ; deshalb ; deswegen ; folglich ; man kann also sagen, dass ...} & C'est pourquoi ; Par conséquent ; On peut donc dire que... \\
\hline
\end{tabularx}
\end{center}

\begin{exoresolu}[Analyser un argument avec citation et Konjunktiv II]
\textbf{Énoncé :} À partir de la phrase du texte : \all{„Die Impfung ist der einzige Weg, die Pandemie zu stoppen,“ betont Dr. Mbarga.}, rédigez un paragraphe d'analyse (These – Beleg – Analyse) en utilisant le subjonctif II pour nuancer.

\tcblower

\textbf{Solution détaillée :}
\all{Ein zentrales Argument des Textes ist die Notwendigkeit der Impfung. Dr. Mbarga betont: „Die Impfung ist der einzige Weg, die Pandemie zu stoppen.“ Hierbei verwendet er \textbf{den Superlativ} „der einzige Weg“, um die Alternativlosigkeit der Schulmedizin zu unterstreichen. \textbf{Analyse :} Der Autor stellt die wissenschaftliche Evidenz über die traditionelle Erfahrung. \textbf{Konjunktiv II für Nuance :} \textit{Man könnte jedoch einwenden, dass} viele Patienten \textit{auf Pflanzen vertrauen würden}, \textit{wenn} die Impfstoffe nicht verfügbar \textit{wären}. (On pourrait cependant objecter que beaucoup de patients feraient confiance aux plantes si les vaccins n'étaient pas disponibles.)}
\end{exoresolu}

\courssub{La conclusion : bilan et opinion personnelle (Schluss)}

Le \all{Schluss} est la \textbf{seule} partie où votre opinion personnelle (\cle{Standpunkt}) est attendue et notée. Elle se structure en trois temps : \textbf{Bilan} (réponse à la problématique) $\rightarrow$ \textbf{Meinung} (avis justifié) $\rightarrow$ \textbf{Ausblick} (ouverture).

\begin{attention}[Piège majeur : L'ordre des mots après « Meiner Meinung nach »]
C'est l'erreur \textbf{n°1} des francophones ! En allemand, \all{Meiner Meinung nach} est un \cle{complément de phrase} placé en \textbf{position 1}. Le verbe conjugué \textbf{doit} suivre immédiatement (position 2). \textbf{Jamais de virgule, jamais d'inversion du sujet et du verbe.}
\begin{itemize}
    \item \textbf{Richtig :} \all{Meiner Meinung nach \textbf{sollten} beide Medizinen zusammenarbeiten.} (Verbe position 2)
    \item \textbf{Falsch :} \all{Meiner Meinung nach, \textbf{sollten} beide...} (Virgule interdite)
    \item \textbf{Falsch :} \all{Meiner Meinung nach \textbf{both} Medizinen \textbf{sollten}...} (Sujet avant verbe = position 3 = erreur)
\end{itemize}
\emph{Règle d'or :} \cle{[Complément] + [Verbe] + [Sujet] + ...}
\end{attention}

\begin{exemplebox}[Expressions clés pour le Schluss (avec traduction)]
\begin{itemize}
    \item \all{Zusammenfassend lässt sich sagen, dass die Krise beide Systeme herausfordert.} = En résumé, on peut dire que la crise met les deux systèmes au défi.
    \item \all{Meiner Meinung nach \textbf{müssen} traditionelle Heiler und Ärzte \textbf{zusammenarbeiten}.} = À mon avis, les guérisseurs traditionnels et les médecins \textbf{doivent} travailler ensemble. (\textit{Verbe position 2 !})
    \item \all{Ich bin der Ansicht, dass eine Integration der Heilpflanzen in das Gesundheitssystem \textbf{sinnvoll wäre}.} = Je suis d'avis qu'une intégration des plantes médicinales dans le système de santé \textbf{serait judicieuse}. (Konjunktiv II : \all{wäre})
    \item \all{Es bleibt zu hoffen, dass die Regierung klare Regeln für die Kooperation schafft.} = Il reste à espérer que le gouvernement crée des règles claires pour la coopération.
    \item \all{Abschließend stellt sich die Frage, ob andere afrikanische Länder diesem Beispiel folgen werden.} = En guise de conclusion, se pose la question de savoir si d'autres pays africains suivront cet exemple.
\end{itemize}
\end{exemplebox}

\begin{exercice}[Entraînement : Rédiger votre Schluss]
Rédigez en allemand une conclusion (environ 60-80 mots) pour le commentaire sur le texte \all{„Zwischen Impfung und Heilpflanze“}.
\textbf{Consignes :}
\begin{enumerate}
    \item Faites le bilan : répondez à la question \all{Kann die traditionelle Medizin die moderne Medizin ergänzen?}
    \item Donnez votre avis personnel justifié (utilisez \all{Meiner Meinung nach} + verbe position 2, et au moins un \all{Konjunktiv II} : \all{wäre / könnte / sollte}).
    \item Terminez par une ouverture (avenir, recherche, politique de santé).
\end{enumerate}
\end{exercice}

\begin{corrige}[Exercice : Proposition de Schluss]
\all{Zusammenfassend lässt sich sagen, dass die traditionelle Medizin die moderne Medizin durchaus ergänzen kann, besonders in ländlichen Regionen Kameruns, wo Ärzte fehlen. Meiner Meinung nach \textbf{sollte} der Staat eine offizielle Anerkennung der Heiler \textbf{schaffen}, damit Patienten nicht zwischen zwei Welten wählen \textbf{müssen}. Eine enge Zusammenarbeit \textbf{wäre} der beste Weg, das Vertrauen der Bevölkerung zu gewinnen. \textbf{Abschließend} bleibt die Frage, ob die wissenschaftliche Forschung die Wirksamkeit der einheimischen Heilpflanzen systematisch prüfen wird.}
\end{corrige}

\courssub{Commentaire modèle guidé (extrait complet)}

Voici un extrait cohérent réunissant \all{Einleitung}, la structure du \all{Hauptteil} et le \all{Schluss} pour le texte fictif étudié. Observez l'enchaînement des temps (Präteritum pour le récit du texte, Präsens pour l'analyse, Konjunktiv II pour l'opinion).

\begin{textebox}[Commentaire modèle : „Zwischen Impfung und Heilpflanze“]
\all{\textbf{Einleitung:} Der Zeitungsartikel „Zwischen Impfung und Heilpflanze“ (Cameroon Tribune, 12.10.2023) behandelt die Gesundheitskrise in Kamerun und die Frage, ob traditionelle und moderne Medizin zusammenarbeiten können. Der Autor stellt zwei Positionen gegenüber: die Heilerin Mama Ngo vertraut auf die Kraft einheimischer Pflanzen, während Dr. Mbarga die Impfung als einzige wissenschaftliche Lösung verteidigt. Problemfrage: Kann die traditionelle Medizin die moderne Medizin ergänzen? Der Text gliedert sich in drei Teile: Zunächst die Sicht der Tradition, dann das Argument der Schulmedizin, schließlich Perspektiven der Kooperation.}

\all{\textbf{Hauptteil:} \textit{Argument 1 (Tradition):} Der erste Teil porträtiert Mama Ngo. Das Zitat „Unsere Vorfahren kannten die Medizin des Waldes“ zeigt das tiefe Vertrauen in das überlieferte Wissen. Der Autor nutzt lexikalisch das \textbf{Wortfeld} „Natur, Vertrauen, Erfahrung“ (Heilpflanze, Wissen, Generation). \textit{Argument 2 (Moderne):} Im Kontrast dazu steht Dr. Mbarga. Er argumentiert mit „Evidenz“, „Studie“, „Sicherheit“. Die Aussage „Ohne Impfung keine Herdenimmunität“ verdeutlicht den wissenschaftlichen Anspruch. \textit{Argument 3 (Synthese):} Der dritte Teil skizziert Kooperationsmodelle. \textbf{Der Ausdruck} „Brücken bauen“ (Metapher) signalisiert den Wunsch nach Integration.}

\all{\textbf{Schluss:} Zusammenfassend zeigt der Text, dass beide Systeme Legitimität besitzen, aber unterschiedliche Logiken folgen. Meiner Meinung nach \textbf{muss} der Staat rechtliche Rahmen \textbf{schaffen}, damit Kooperation nicht nur ein Schlagwort bleibt. Eine institutionalisierte Zusammenarbeit \textbf{könnte} die Gesundheitsversorgung in Kamerun nachhaltig verbessern. Es bleibt zu hoffen, dass die nächste Generation von Ärzten auch traditionelles Wissen im Studium kennenlernt.}
\end{textebox}

\begin{methode}[Check-list « Spécial Bac » pour réussir le commentaire]
\begin{itemize}
    \item \textbf{Ne recopiez JAMAIS} de longs passages du texte. Citez \textbf{court} (max 1 ligne) entre guillemets allemands \all{„ ... “}.
    \item \textbf{Séparez visuellement} les trois parties (saut de ligne, alinéa).
    \item \textbf{Analysez, ne résumez pas :} Dites \all{WARUM} l'auteur choisit ce mot, cette structure, ce temps verbal (lien avec le thème).
    \item \textbf{Opinion UNIQUEMENT dans le Schluss.} Pas de \all{Ich finde...} dans l'Einleitung ou le Hauptteil.
    \item \textbf{Relisez l'ordre des mots :} Verbe position 2 (surtout après \all{Meiner Meinung nach, Zunächst, Allerdings, Daher}).
    \item \textbf{Soignez l'orthographe des noms} (majuscules !) et les Umlaute (\all{ä, ö, ü, ß}).
\end{itemize}
\end{methode}

\begin{savaistu}[Contexte culturel : Médecine traditionnelle au Cameroun et en Allemagne]
Au Cameroun, l'OMS estime que 80\% de la population recourt aux médecines traditionnelles (tradipraticiens, plantes comme \all{Prunus africana} ou \all{Moringa}). Le Ministère de la Santé publique a créé un \all{Département de Médecine Traditionnelle} pour encadrer la pratique. En Allemagne, la \all{Naturheilkunde} (médecine naturelle) et l'\all{Homöopathie} sont très populaires et souvent remboursées par les caisses d'assurance maladie (\all{Krankenkassen}). Le terme \all{Integrative Medizin} désigne aujourd'hui la combinaison voulue des deux approches — exactement le sujet de notre problématique !
\end{savaistu}

\coursec{Production guidée et réinvestissement}

Dans la section précédente, nous avons appris à commenter méthodiquement le texte sur la crise sanitaire et la médecine traditionnelle : repérage du thème, analyse des idées, appréciation personnelle. Il est temps maintenant de \cle{réinvestir} tout ce que nous avons acquis — le vocabulaire de la santé, les verbes pronominaux, le conseil au subjonctif II — dans des productions concrètes : réponses écrites, résumé, paragraphe guidé et débat oral. C'est l'étape où le savoir devient un \cle{savoir-faire}, exactement ce qui sera exigé de vous à l'épreuve du Baccalauréat.

\courssub{Réponses écrites aux questions de compréhension}

Répondre à une question de compréhension, c'est rédiger une \cle{phrase complète} en allemand, en reprenant les termes de la question, sans jamais copier mot à mot de longs passages du texte. Le verbe conjugué reste en deuxième position dans la phrase principale ; dans une subordonnée introduite par \all{dass} ou \all{weil}, il va à la fin.

\begin{exoresolu}[Questions de compréhension sur le texte support]
\textbf{Frage 1 :} Warum vertrauen viele Menschen im Dorf den Heilpflanzen?\\
\textbf{Frage 2 :} Was macht der traditionelle Heiler, wenn ein Patient zu ihm kommt?\\
\textbf{Frage 3 :} Welche Probleme gab es während der Gesundheitskrise?\\
\textbf{Frage 4 :} Wie kann man sich vor Krankheiten schützen?\\
\textbf{Frage 5 :} Was denkt der Autor über die Zusammenarbeit von moderner und traditioneller Medizin?
\tcblower
\textbf{Antwort 1 :} Viele Menschen im Dorf vertrauen den Heilpflanzen, weil diese Medizin billig ist und weil sie sie von ihren Großeltern kennen. « Beaucoup de gens au village font confiance aux plantes médicinales parce que cette médecine est bon marché et parce qu'ils la connaissent de leurs grands-parents. »\\
\textbf{Antwort 2 :} Wenn ein Patient zu ihm kommt, untersucht ihn der Heiler und bereitet ihm ein Heilmittel aus Blättern und Wurzeln zu. « Quand un patient vient le voir, le guérisseur l'examine et lui prépare un remède à base de feuilles et de racines. »\\
\textbf{Antwort 3 :} Während der Gesundheitskrise gab es zu wenige Ärzte, und viele Kranke konnten sich nicht ins Krankenhaus begeben. « Pendant la crise sanitaire, il y avait trop peu de médecins et beaucoup de malades ne pouvaient pas se rendre à l'hôpital. »\\
\textbf{Antwort 4 :} Man kann sich schützen, wenn man sich impfen lässt, sich oft die Hände wäscht und sich gesund ernährt. « On peut se protéger en se faisant vacciner, en se lavant souvent les mains et en se nourrissant sainement. »\\
\textbf{Antwort 5 :} Der Autor denkt, dass sich die moderne und die traditionelle Medizin ergänzen sollten. « L'auteur pense que la médecine moderne et la médecine traditionnelle devraient se compléter. »
\end{exoresolu}

\begin{attention}[La place du verbe dans la réponse]
Erreur fréquente des francophones : après \all{dass} et \all{weil}, le verbe conjugué va \textbf{en fin de proposition}, contrairement au français. On écrit : \all{Ich denke, dass man sich impfen lassen sollte.} « Je pense qu'on devrait se faire vacciner. » — et jamais \emph{dass man sollte sich impfen lassen}. Attention aussi au pronom réfléchi : \all{sich impfen lassen}, \all{sich erkälten}, \all{sich erholen} ne perdent jamais leur \all{sich}.
\end{attention}

\courssub{Résumé du texte avec connecteurs temporels}

Un résumé (\all{die Zusammenfassung}) reprend les idées essentielles du texte en 4 à 5 phrases, avec ses propres mots, au présent, et reliées par des \cle{connecteurs}. Voici la frise des connecteurs temporels à mobiliser :

\begin{popfigure}
\begin{center}
\begin{tikzpicture}[scale=1, every node/.style={font=\small}]
\draw[-{Stealth[length=3mm]}, line width=1.2pt, popblue] (0,0) -- (11.5,0);
\foreach \x/\mot/\trad in {0.8/{zuerst}/{d'abord}, 3.9/{dann}/{puis}, 7.0/{danach}/{ensuite}, 10.3/{schließlich}/{finalement}}{
  \fill[poporange] (\x,0) circle (2.5pt);
  \node[at={(\x,0.6)}, align=center, text=popdark] {\textbf{\mot}\\ \emph{\trad}};
}
\end{tikzpicture}
\end{center}
\legende{Frise des connecteurs temporels pour structurer un résumé}
\end{popfigure}

\begin{center}
\begin{tabularx}{\linewidth}{|l|X|X|}
\hline
\rowcolor{popblueL} \textbf{Connecteur} & \textbf{Emploi} & \textbf{Exemple} \\
\hline
\all{zuerst} & première étape & \all{Zuerst beschreibt der Text das Dorf.} \\
\hline
\all{dann} & étape suivante & \all{Dann spricht er über die Krise.} \\
\hline
\all{danach} & après cela & \all{Danach erklärt er die Rolle der Heiler.} \\
\hline
\all{außerdem} & ajout d'une idée & \all{Außerdem nennt er die Impfung.} \\
\hline
\all{schließlich} & conclusion & \all{Schließlich gibt er seine Meinung.} \\
\hline
\end{tabularx}
\end{center}

\begin{exemplebox}[Résumé modèle du texte support]
\all{Zuerst stellt der Text das Dorf und seine Heilpflanzen vor. Dann beschreibt er die Gesundheitskrise: Es gab zu wenige Ärzte und Medikamente. Außerdem erklärt er, wie die traditionellen Heiler den Kranken geholfen haben. Schließlich meint der Autor, dass sich die moderne Medizin und die traditionelle Medizin ergänzen sollten.}\\
« D'abord, le texte présente le village et ses plantes médicinales. Puis il décrit la crise sanitaire : il y avait trop peu de médecins et de médicaments. De plus, il explique comment les guérisseurs traditionnels ont aidé les malades. Finalement, l'auteur estime que la médecine moderne et la médecine traditionnelle devraient se compléter. »
\end{exemplebox}

\courssub{Production guidée : \all{Gesundheit in meinem Dorf}}

\begin{methode}[Rédiger un paragraphe guidé en quatre étapes]
\begin{enumerate}
\item \textbf{Noter 5 mots-clés du thème} : par exemple \all{die Heilpflanze, der Heiler, sich erkälten, die Impfung, das Heilmittel}.
\item \textbf{Écrire une phrase par idée} : une phrase simple et correcte vaut mieux qu'une phrase longue et fausse. Suis la grille : début (\all{In meinem Dorf...}), milieu (\all{Viele Menschen...} / \all{Die Heiler...}), conseils (\all{Man sollte...} / \all{Es wäre besser, wenn...}), conclusion (\all{Ich denke, dass...}).
\item \textbf{Vérifier la place du verbe et les pronoms réfléchis} : verbe en 2\up{e} position dans la principale, en fin après \all{dass}, \all{weil}, \all{wenn} ; \all{sich} présent avec chaque verbe pronominal.
\item \textbf{Relire en vérifiant articles et pluriels} : \all{die Heilpflanze, die Heilpflanzen} ; \all{das Heilmittel, die Heilmittel} ; majuscule à tous les noms.
\end{enumerate}
\end{methode}

\begin{exemplebox}[Phrases utiles traduites]
\all{In meinem Dorf benutzen viele Menschen Heilpflanzen.} « Dans mon village, beaucoup de gens utilisent des plantes médicinales. »\\
\all{Man sollte die traditionelle Medizin respektieren.} « On devrait respecter la médecine traditionnelle. »\\
\all{Es wäre besser, wenn die Heiler mit den Ärzten zusammenarbeiten würden.} « Il serait mieux que les guérisseurs travaillent ensemble avec les médecins. »\\
\all{Wenn ich mich erkälte, bereitet mir meine Großmutter einen Tee aus Blättern zu.} « Quand je m'enrhume, ma grand-mère me prépare une tisane de feuilles. »
\end{exemplebox}

\begin{textebox}[Production modèle : „Gesundheit in meinem Dorf“]
In meinem Dorf in der Nähe von Yaoundé benutzen viele Menschen Heilpflanzen, wenn sie krank sind. Wenn sich jemand erkältet oder sich schlecht fühlt, geht er zuerst zum traditionellen Heiler. Der Heiler untersucht den Patienten und bereitet ihm ein Heilmittel aus Blättern, Rinden und Wurzeln zu. Viele alte Menschen erinnern sich gern an die Zeit, als es noch kein Krankenhaus im Dorf gab, und sie vertrauen den Pflanzen mehr als den Tabletten. Während der Gesundheitskrise haben sich die Heiler um viele Kranke gekümmert, weil die Krankenhäuser voll waren. Aber Vorsicht: Manche Heilmittel sind gefährlich, wenn man die Dosis nicht kennt. Man sollte die traditionelle Medizin respektieren, aber man sollte sich auch impfen lassen und zum Arzt gehen, wenn die Krankheit ernst ist. Es wäre besser, wenn die Heiler und die Ärzte zusammenarbeiten würden. Ich denke, dass sich die beiden Medizinen sehr gut ergänzen können: Die Pflanzen meiner Großmutter und die Impfung des Arztes sind keine Feinde, sondern Partner für unsere Gesundheit.

\textbf{Glossaire :} die Heilpflanze, -n (plante médicinale) ; sich erkälten (s'enrhumer) ; der Heiler, - (guérisseur) ; das Heilmittel, - (remède) ; die Rinde, -n (écorce) ; die Wurzel, -n (racine) ; sich erinnern an + Akk. (se souvenir de) ; sich kümmern um + Akk. (s'occuper de) ; die Dosis, die Dosen (dose) ; sich impfen lassen (se faire vacciner) ; sich ergänzen (se compléter).
\end{textebox}

Comptons les exigences dans ce modèle : \textbf{cinq} verbes pronominaux (\all{sich erkältet}, \all{erinnern sich}, \all{haben sich ... gekümmert}, \all{sich impfen lassen} et \all{sich ergänzen}) et deux conseils au subjonctif II (\all{Man sollte...}, \all{Es wäre besser, wenn...}). C'est exactement le contrat de l'exercice.

\begin{exercice}[À ton tour]
Rédige un paragraphe de 8 à 10 lignes intitulé \all{Gesundheit in meinem Dorf} en suivant la grille : début (\all{In meinem Dorf...}), milieu (\all{Viele Menschen...} / \all{Die Heiler...}), deux conseils (\all{Man sollte...} / \all{Es wäre besser, wenn...}), conclusion (\all{Ich denke, dass...}). Utilise au moins 3 verbes pronominaux différents et souligne-les.
\end{exercice}

\begin{corrige}[Exercice : éléments de correction]
Vérifier : 1) la présence d'au moins 3 verbes pronominaux avec \all{sich} correctement placé (\all{...haben sich die Heiler um die Kranken gekümmert}) ; 2) deux conseils au Konjunktiv II (\all{sollte}, \all{wäre}, \all{würden}) ; 3) le verbe en fin de subordonnée après \all{dass}, \all{weil}, \all{wenn} ; 4) articles et pluriels corrects ; 5) majuscules des noms. Une production type peut reprendre la structure du modèle ci-dessus avec des détails personnels (le nom du village, une plante locale comme \all{die Rinde des Prunus africanus} ou simplement \all{die Blätter}).
\end{corrige}

\courssub{Expression orale : débat en binômes}

\begin{experience}[Débat : „Moderne Medizin oder Heilpflanzen ?“]
Par binômes, l'élève A défend la médecine moderne, l'élève B la médecine traditionnelle. Chacun parle une minute, puis répond aux arguments de son partenaire.\\
\textbf{Prendre la parole :} \all{Ich bin der Meinung, dass...} / \all{Meiner Meinung nach...} « À mon avis... »\\
\textbf{Argumenter :} \all{Ein Argument dafür ist, dass...} / \all{Außerdem...}\\
\textbf{Contredire poliment :} \all{Das stimmt, aber...} / \all{Ich sehe das anders, weil...} « C'est vrai, mais... / Je vois cela autrement, parce que... »\\
\textbf{Conclure :} \all{Zum Schluss denke ich, dass sich beide Medizinen ergänzen sollten.}\\
Le professeur circule et note les verbes pronominaux et les conseils au Konjunktiv II effectivement employés.
\end{experience}

\courssub{Correction collective et auto-évaluation}

La correction collective se fait au tableau : deux productions d'élèves sont recopiées, la classe repère ensemble les réussites et les erreurs (place du verbe, \all{sich} oublié, article faux). Puis chaque élève s'auto-évalue avec la grille suivante :

\begin{center}
\begin{tabularx}{\linewidth}{|X|c|c|c|}
\hline
\rowcolor{popblueL} \textbf{Critère} & \textbf{Oui} & \textbf{Un peu} & \textbf{Non} \\
\hline
J'ai utilisé au moins 3 verbes pronominaux corrects. & & & \\
\hline
J'ai formulé 2 conseils au Konjunktiv II (\all{sollte}, \all{wäre}). & & & \\
\hline
Le verbe est en fin de phrase après \all{dass} / \all{weil}. & & & \\
\hline
Mes articles et mes pluriels sont corrects. & & & \\
\hline
J'ai utilisé des connecteurs (\all{zuerst}, \all{dann}, \all{schließlich}). & & & \\
\hline
J'ai mis la majuscule à tous les noms. & & & \\
\hline
\end{tabularx}
\end{center}

\begin{attention}[Dernière vérification avant de rendre]
Relis toujours ta production à voix basse : chaque verbe pronominal a-t-il son \all{sich} ? Chaque subordonnée a-t-elle son verbe à la fin ? Chaque nom a-t-il sa majuscule et son article ? Ces trois questions éliminent 80 \% des erreurs des copies francophones.
\end{attention}

\begin{savaistu}[L'OMS et la médecine traditionnelle]
\all{Die Weltgesundheitsorganisation} (WHO), l'Organisation mondiale de la santé, reconnaît officiellement la médecine traditionnelle et encourage les États à la réglementer et à étudier scientifiquement les plantes médicinales. Plusieurs pays africains, dont le Cameroun, ont créé des structures pour valoriser les guérisseurs tout en protégeant les patients. C'est un argument précieux et vérifiable pour ton essai au Bac : la médecine traditionnelle n'est pas l'ennemie de la science, elle peut devenir sa partenaire.
\end{savaistu}

\begin{aretenir}[Ce qu'il faut retenir de cette section]
\begin{itemize}
\item Répondre aux questions en phrases complètes, verbe en fin après \all{dass} et \all{weil}.
\item Résumer en 4-5 phrases avec \all{zuerst, dann, außerdem, schließlich}.
\item Rédiger un paragraphe guidé : 5 mots-clés, une phrase par idée, vérification du verbe et des pronoms réfléchis, relecture des articles et pluriels.
\item Débattre à l'oral avec des formules de prise de parole et des conseils au Konjunktiv II.
\end{itemize}
\end{aretenir}

\newpage

\coursec{Activité d'intégration}

\courssub{Objectifs de l'activité}

Cette activité d'intégration clôt la leçon sur les maladies, la crise sanitaire et la médecine traditionnelle. Elle vise à faire réinvestir, dans une tâche de communication authentique, l'ensemble des acquis de la leçon :

\begin{itemize}
\item \cle{Compréhension} : comprendre des arguments oraux et écrits sur la santé et réagir de manière pertinente ;
\item \cle{Lexique} : employer correctement le vocabulaire thématique (\all{die Krankheit}, \all{der Arzt}, \all{das Heilmittel}, \all{die Heilpflanze}, \all{die Impfung}, \all{die Pandemie}) avec les bons articles et pluriels ;
\item \cle{Grammaire 1} : conjuguer et placer correctement les verbes pronominaux (\all{sich erkälten}, \all{sich kümmern um}, \all{sich erholen von}, \all{sich impfen lassen}, \all{sich behandeln}) ;
\item \cle{Grammaire 2} : formuler des conseils au subjonctif II (\all{Du solltest \ldots}, \all{Man könnte \ldots}, \all{Es wäre besser, wenn \ldots}, \all{Ich würde \ldots}) ;
\item \cle{Production} : rédiger un court article de presse structuré (environ 10 lignes) donnant un avis personnel argumenté.
\end{itemize}

\courssub{Situation}

Le lycée de Nkol-Eton à Yaoundé organise, dans le cadre de la semaine de la santé, un débat public animé par le club d'allemand. Le thème retenu cette année concerne directement la vie des élèves et de leurs familles : face à une maladie, faut-il faire confiance à la médecine moderne, aux plantes médicinales, ou aux deux ? Le journal du lycée publiera le meilleur article rédigé à l'issue du débat. Voici le document qui lance la réflexion, affiché au tableau :

\begin{textebox}[Affiche du club d'allemand — Lycée de Nkol-Eton, Yaoundé]
\textbf{DEBATTE: Impfung oder Heilpflanzen?}\par
Nach der Pandemie fragen sich viele Kameruner: Sollen wir uns nur auf die moderne Medizin verlassen? Oder sind die Heilpflanzen unserer Großeltern genauso wichtig?\par
Der Arzt sagt: „Die Impfung schützt vor gefährlichen Krankheiten. Man sollte sich impfen lassen.“\par
Der Heiler antwortet: „Die Natur bietet viele Heilmittel. Man könnte zuerst die traditionelle Medizin probieren.“\par
Und Sie? Was denken Sie? Kommen Sie zur Debatte am Freitag um 14 Uhr!
\end{textebox}

\begin{definition}[Glossaire de l'affiche]
\all{sich verlassen auf + Akk.} : compter sur, se fier à ; \all{der Heiler, -} : le guérisseur ; \all{schützen vor + Dat.} : protéger contre ; \all{probieren} : essayer ; \all{die Debatte, -n} : le débat ; \all{sich fragen} : se demander.
\end{definition}

\courssub{Méthode : Rédiger un article de presse}

\begin{methode}[Structure d'un article d'opinion court]
\begin{enumerate}
\item \textbf{Titre} : accrocheur, résume le sujet (ex. \all{Gesundheit in Kamerun: Ein Weg zwischen zwei Medizinen}).
\item \textbf{Introduction (Lead)} : répond aux questions \all{Wann? Wo? Was? Wer?} (Quand ? Où ? Quoi ? Qui ?). Pose le contexte.
\item \textbf{Corps de l'article} : présente les deux positions (\all{Einerseits \ldots / Andererseits \ldots}) avec du discours rapporté direct ou indirect. Utilise le vocabulaire thématique.
\item \textbf{Prise de position personnelle} : introduite par \all{Meiner Meinung nach} / \all{Ich bin der Ansicht, dass} + verbe en 2e position. Apporte un argument nouveau (synthèse, proposition).
\item \textbf{Conclusion} : phrase de chute, ouverture ou conseil au \cle{Subjonctif II} (\all{Man könnte \ldots}, \all{Es wäre besser, \ldots}).
\item \textbf{Relecture} : vérifier l'ordre des mots (V2), les majuscules des noms, la conjugaison (pronominaux, modaux, KII), la ponctuation des guillemets „…“.
\end{enumerate}
\end{methode}

\courssub{Consignes}

\begin{enumerate}
\item \textbf{Lecture de l'affiche (5 min).} Lis l'affiche en silence, puis réponds oralement : \all{Worum geht es in der Debatte? Wer spricht? Was sagen der Arzt und der Heiler?} Souligne dans le texte deux verbes pronominaux (\all{fragen sich}, \all{sich verlassen}, \all{sich impfen lassen}) et deux conseils au subjonctif II (\all{sollte \ldots lassen}, \all{könnte \ldots probieren}).

\item \textbf{Formation des groupes (5 min).} La classe se divise en trois groupes : les \all{Ärzte} (médecins modernes), les \all{Heiler} (guérisseurs traditionnels) et les \all{Journalisten} (journalistes). Choisissez un porte-parole par groupe.

\item \textbf{Préparation écrite des arguments (15 min).} Chaque groupe rédige au moins \textbf{cinq phrases} :
\begin{itemize}
\item les \all{Ärzte} défendent la vaccination et la médecine scientifique ;
\item les \all{Heiler} défendent les plantes médicinales et le savoir traditionnel ;
\item les \all{Journalisten} préparent six questions avec des mots interrogatifs (\all{Warum? Wie? Wann? Welche? Wer? Was?}).
\end{itemize}
\textbf{Obligation de langue} : chaque texte doit contenir au moins \textbf{deux verbes pronominaux} correctement conjugués (place du pronom, accord) et \textbf{deux conseils au subjonctif II} (forme correcte du modal/auxiliaire + infinitif à la fin).

\item \textbf{Le débat (20 min).} Les journalistes posent leurs questions. Chaque camp répond et argumente. Le professeur veille à la prise de parole équitable et note les erreurs de grammaire (ordre des mots, cas après prépositions, KII) sans interrompre.

\item \textbf{Production écrite individuelle (20 min).} Chaque élève rédige seul un court article de \textbf{10 lignes} intitulé \all{„Gesundheit in Kamerun: Ein Weg zwischen zwei Medizinen“}. L'article doit : résumer les deux positions du débat, contenir au moins \textbf{deux verbes pronominaux} et \textbf{un conseil au subjonctif II}, et donner ton avis personnel avec \all{Meiner Meinung nach \ldots} (verbe immédiatement après).

\item \textbf{Relecture croisée (10 min).} Échange ta copie avec un voisin : vérifie la place du pronom réfléchi (après le sujet), la conjugaison du subjonctif II (\all{sollte, könnte, wäre, würde}) et les majuscules des noms. Corrige au crayon.
\end{enumerate}

\courssub{Ressources à mobiliser}

\begin{aretenir}[Boîte à outils linguistique]
\textbf{Lexique thématique (articles \& pluriels)} : \all{die Krankheit, -en}, \all{der Arzt, die Ärzte}, \all{das Heilmittel, -}, \all{die Heilpflanze, -n}, \all{die Impfung, -en}, \all{die Pandemie, -n}, \all{das Fieber}, \all{die Nebenwirkung, -en}, \all{die Medizin, -en}, \all{der Heiler, -}, \all{die Wissenschaft}.\par
\textbf{Verbes pronominaux (contexte santé)} :
\begin{itemize}
\item \all{sich erkälten} (s'enrhumer) — \all{Ich erkälte mich oft.}
\item \all{sich kümmern um + Akk.} (s'occuper de) — \all{Kümmert euch um eure Gesundheit!}
\item \all{sich erholen von + Dat.} (se remettre de) — \all{Er erholt sich von der Grippe.}
\item \all{sich fürchten vor + Dat.} (avoir peur de) — \all{Viele fürchten sich vor der Spritze.}
\item \all{sich impfen lassen} (se faire vacciner) — \all{Man sollte sich impfen lassen.}
\item \all{sich behandeln} (se soigner / se traiter) — \all{Er behandelt sich mit Heilpflanzen.}
\end{itemize}
\textbf{Conseils au Subjonctif II (KII)} :
\begin{itemize}
\item \all{Du solltest mehr Wasser trinken.} (\all{sollen} $\rightarrow$ \all{sollte})
\item \all{Man könnte die Heilpflanzen erforschen.} (\all{können} $\rightarrow$ \all{könnte})
\item \all{Es wäre besser, beide Medizinen zu verbinden.} (\all{sein} $\rightarrow$ \all{wäre} + infinitif)
\item \all{Ich würde zum Arzt gehen.} (\all{werden} $\rightarrow$ \all{würde} + infinitif)
\end{itemize}
\textbf{Prise de position \& articulation} : \all{Meiner Meinung nach \ldots} (verbe 2e) / \all{Ich bin der Ansicht, dass \ldots} (verbe à la fin) / \all{Einerseits \ldots, andererseits \ldots} / \all{Zudem / Außerdem / Allerdings}.
\end{aretenir}

\begin{attention}[Pièges à éviter pour les francophones]
\begin{itemize}
\item \textbf{Place du pronom réfléchi} : en phrase principale, il suit \textbf{immédiatement le sujet} (verbe en 2e position). \all{Ich erkälte mich.} / \all{Man sollte sich impfen lassen.} — Pas \all{*Ich mich erkälte} ni \all{*Man sich sollte impfen lassen}.
\item \textbf{Subjonctif II (KII) \& ordre des mots} : le verbe conjugué (\all{sollte, könnte, wäre, würde}) reste en 2e position ; l'infinitif (ou le participe passé au Perfekt) va à la \textbf{fin de la phrase}. \all{Man könnte die Pflanzen erforschen.}
\item \textbf{Gestion des cas avec verbes pronominaux} : \all{sich kümmern um + Akkusativ}, \all{sich erholen von + Dativ}, \all{sich fürchten vor + Dativ}, \all{schützen vor + Dativ}. Apprends le couple verbe + préposition + cas.
\item \textbf{Genre \& Majuscules} : \all{die Impfung} (f.), \all{das Heilmittel} (n.), \all{der Arzt} (m.). \textbf{Tout nom prend une majuscule} en allemand.
\item \textbf{Après \all{Meiner Meinung nach}} : le verbe conjugué vient \textbf{immédiatement} (place 2). \all{Meiner Meinung nach sollten wir \ldots} — Pas \all{*Meiner Meinung nach wir sollten}.
\item \textbf{Faux ami} : \all{die Medizin} = la médecine / le médicament (science) ; \all{das Medikament} = le médicament (produit). \all{die Gift} = le poison (pas le cadeau !).
\end{itemize}
\end{attention}

\courssub{Proposition de production et corrigé commenté}

\begin{exoresolu}[Article modèle : \all{„Gesundheit in Kamerun: Ein Weg zwischen zwei Medizinen“}]
Rédige un article de 10 lignes résumant le débat et donnant ton avis, avec au moins deux verbes pronominaux et un conseil au subjonctif II.
\tcblower
\textbf{Production modèle :}\par
\all{Am Freitag fand an unserer Schule eine Debatte über die Gesundheit in Kamerun statt. Die Ärzte verteidigten die moderne Medizin: „Die Impfung schützt vor gefährlichen Krankheiten, und man sollte sich impfen lassen“, sagte ihre Sprecherin. Die Heiler antworteten, dass die Natur viele Heilmittel biete. „Wenn man sich erkältet, kann man sich zuerst mit Heilpflanzen behandeln“, erklärte ein Heiler. Die Journalisten stellten viele Fragen, zum Beispiel: „Warum fürchten sich manche Menschen vor der Impfung?“ Meiner Meinung nach sollten wir beide Medizinen verbinden. Es wäre falsch, die traditionelle Medizin zu vergessen, aber bei einer Pandemie brauchen wir die Wissenschaft. Man könnte zum Beispiel die Heilpflanzen wissenschaftlich erforschen. So könnte sich Kamerun um die Gesundheit aller kümmern.}\par
\textbf{Commentaire des choix de langue :}
\begin{itemize}
\item \textbf{Verbes pronominaux} : \all{sich impfen lassen} (construction avec \all{lassen}, réfléchi \all{sich} après le sujet \all{man}) ; \all{sich erkältet} (Présent, accord avec \all{man}) ; \all{sich behandeln} (infinitif après modal \all{kann}, réfléchi \all{sich} après sujet \all{man}) ; \all{fürchten sich} (réfléchi après sujet \all{manche Menschen} dans la question) ; \all{könnte sich \ldots kümmern} (réfléchi \all{sich} après sujet \all{Kamerun}, verbe modal \all{könnte} en position 2, infinitif \all{kümmern} à la fin).
\item \textbf{Conseils au KII} : \all{man sollte sich impfen lassen} (\all{sollen} $\rightarrow$ \all{sollte}) ; \all{Es wäre falsch \ldots} (\all{sein} $\rightarrow$ \all{wäre} + infinitif) ; \all{Man könnte \ldots erforschen} (\all{können} $\rightarrow$ \all{könnte}) ; \all{sollten wir \ldots verbinden} (\all{sollen} $\rightarrow$ \all{sollten} après \all{Meiner Meinung nach}). Quatre formes variées.
\item \textbf{Discours rapporté} : utilisation du \textbf{Subjonctif I} (\all{biete} au lieu de \all{bietet}) pour marquer la distance dans \all{Die Heiler antworteten, dass die Natur viele Heilmittel biete} (niveau excellence).
\item \textbf{Lexique thématique} : \all{die Impfung}, \all{die Heilpflanzen}, \all{die Heilmittel}, \all{die Pandemie}, \all{die Wissenschaft}, \all{sich erkälten}, \all{sich behandeln}, \all{sich kümmern um} réinvestis avec articles corrects.
\item \textbf{Structure d'article} : introduction factuelle (\all{Am Freitag \ldots statt}), présentation équilibrée des deux camps (discours direct/indirect), avis personnel argumenté (\all{Meiner Meinung nach} + V2), proposition concrète (\all{Man könnte \ldots}), conclusion ouverte (\all{So könnte \ldots}).
\end{itemize}
\end{exoresolu}

\courssub{Bilan de l'activité}

À l'issue de ce débat-délibéré, l'élève a mobilisé simultanément les quatre compétences : il a lu et compris un document authentique (l'affiche), écouté et interrogé ses camarades (compréhension et expression orales), puis rédigé un article de presse structuré (expression écrite). Les deux points de grammaire de la leçon — les verbes pronominaux (place, cas, conjugaison) et le conseil au subjonctif II (formes \all{sollte, könnte, wäre, würde} + syntaxe) — ont été réemployés dans un contexte de communication réel et non dans de simples exercices mécaniques. Le débat a en outre permis d'aborder une question de société qui touche directement les élèves camerounais : la complémentarité entre médecine moderne et médecine traditionnelle, sujet culturel et citoyen au cœur du chapitre « Environnement, santé et bien-être ». Le professeur pourra conserver les meilleurs articles pour le journal du lycée, donnant ainsi une finalité concrète et motivante à la production écrite.

\newpage

\coursec{Méthodes et exercices résolus}

\coursec{Les maladies : la crise sanitaire et la médecine traditionnelle}

\courssub{Texte support : lecture et première compréhension}

\begin{textebox}[Texte authentique : médecine traditionnelle et crise sanitaire au Cameroun]
\all{In Kamerun vertrauen viele Menschen der traditionellen Medizin. Heiler benutzen Heilpflanzen gegen Krankheiten wie Malaria oder Magenschmerzen. Während der Gesundheitskrise haben sich viele Familien auch an die moderne Medizin gewandt: Sie haben sich impfen lassen und die Hygieneregeln beachtet. Experten sagen: Die beste Lösung ist die Zusammenarbeit zwischen traditioneller und moderner Medizin.}
\end{textebox}

\begin{exemplebox}[Premier contact avec le texte]
\begin{enumerate}
\item \textbf{Titre et source} : ce texte (article de presse simulé) traite de la coexistence de deux systèmes de santé au Cameroun.
\item \textbf{Mots transparents} : \all{die Pandemie}, \all{das Medikament}, \all{die Tradition}, \all{modern}, \all{Experten}, \all{die Lösung}.
\item \textbf{Idée générale} : \all{Der Text handelt von der traditionellen und modernen Medizin in Kamerun während der Gesundheitskrise.}
\end{enumerate}
\end{exemplebox}

\courssub{Lexique thématique : la santé et la médecine}

\begin{definition}[Vocabulaire de base : noms avec article et pluriel]
\begin{center}
\begin{tabularx}{\linewidth}{|X|X|X|}
\hline
\rowcolor{popblueL} \textbf{Allemand (article + pluriel)} & \textbf{Français} & \textbf{Exemple contextuel} \\
\hline
\all{die Krankheit, -en} & la maladie & \all{Maria ist eine gefährliche Krankheit.} \\
\hline
\all{der Arzt, die Ärzte} & le médecin & \all{Der Arzt untersucht den Patienten.} \\
\hline
\all{das Heilmittel, -} & le remède & \all{Es gibt kein einfaches Heilmittel gegen Malaria.} \\
\hline
\all{die Heilpflanze, -n} & la plante médicinale & \all{Der Heiler kennt viele Heilpflanzen.} \\
\hline
\all{die Impfung, -en} & la vaccination & \all{Die Impfung schützt vor schweren Krankheiten.} \\
\hline
\all{die Pandemie, -n} & la pandémie & \all{Die Pandemie hat das Gesundheitssystem belastet.} \\
\hline
\all{die Gesundheitskrise} & la crise sanitaire & \all{Während der Gesundheitskrise...} \\
\hline
\all{der Heiler, -} & le guérisseur & \all{Der Heiler behandelt mit Naturmitteln.} \\
\hline
\all{die Hygieneregeln (Pl.)} & les règles d'hygiène & \all{Man muss die Hygieneregeln beachten.} \\
\hline
\end{tabularx}
\end{center}
\end{definition}

\begin{propriete}[Familles de mots et verbes à constructions fixes]
\textbf{Familles lexicales} (dérivation) :
\begin{itemize}
\item \all{heilen} (guérir) $\rightarrow$ \all{der Heiler} (guérisseur), \all{das Heilmittel} (remède), \all{die Heilpflanze} (plante médicinale), \all{heilsam} (curatif).
\item \all{krank} (malade) $\rightarrow$ \all{die Krankheit} (maladie), \all{das Krankenhaus} (hôpital), \all{der Krankenpfleger} (infirmier), \all{krankheitsbedingt} (lié à la maladie).
\end{itemize}

\textbf{Verbes thématiques avec leur régime (préposition + cas)} :
\begin{center}
\begin{tabular}{|l|l|l|}
\hline
\rowcolor{popblueL} \textbf{Verbe} & \textbf{Construction} & \textbf{Exemple} \\
\hline
\all{sich erkälten} & (pronominal, pas de préposition) & \all{Ich habe mich erkältet.} \\
\hline
\all{sich erholen} & \all{von} + \textbf{Datif} & \all{Er erholt sich von der Grippe.} \\
\hline
\all{sich impfen lassen} & (causatif, pronominal) & \all{Wir lassen uns gegen Corona impfen.} \\
\hline
\all{leiden} & \all{an} + \textbf{Datif} & \all{Er leidet an Malaria.} \\
\hline
\all{behandeln} & + \textbf{Akkusativ} (transitif) & \all{Der Arzt behandelt den Patienten.} \\
\hline
\all{sich interessieren} & \all{für} + \textbf{Akkusativ} & \all{Ich interessiere mich für Medizin.} \\
\hline
\all{sich kümmern} & \all{um} + \textbf{Akkusativ} & \all{Sie kümmert sich um die Kranken.} \\
\hline
\end{tabular}
\end{center}
\end{propriete}

\begin{methode}[Maîtriser le lexique : 5 étapes]
\begin{enumerate}
\item \textbf{Apprendre chaque nom avec son article et son pluriel} (voir tableau ci-dessus).
\item \textbf{Regrouper par familles de mots} (ex. \all{heilen} $\rightarrow$ \all{Heiler, Heilmittel, Heilpflanze}).
\item \textbf{Mémoriser les verbes avec leur construction} (préposition + cas) — \emph{le cas est indissociable du verbe}.
\item \textbf{Utiliser les expressions figées de la crise sanitaire} : \all{die Hygieneregeln beachten}, \all{sich die Hände waschen}, \all{eine Maske tragen}, \all{Abstand halten}.
\item \textbf{Réemployer} chaque mot dans une phrase personnelle liée au contexte camerounais (ex. \all{In Yaoundé gibt es viele Heiler.}).
\end{enumerate}
\end{methode}

\begin{exoresolu}[Lexique : compléter un texte à trous]
\textbf{Énoncé} : Complétez avec les mots du tableau : \all{Heilpflanzen — Impfung — Krankheit — Heilmittel — Arzt}.

\all{Malaria ist eine gefährliche (1) \dots\ in Afrika. Viele Patienten gehen zum (2) \dots, aber andere benutzen traditionelle (3) \dots\ aus der Natur. Gegen Malaria gibt es kein einfaches (4) \dots. Die (5) \dots\ gegen andere Krankheiten schützt die Bevölkerung.}
\tcblower
\textbf{Raisonnement} : on analyse le sens de chaque phrase et la nature grammaticale attendue (genre, nombre, cas).

\textbf{Solutions justifiées} :
\begin{enumerate}
\item \all{Krankheit} : « le paludisme est une \emph{maladie} dangereuse » ; \all{eine gefährliche} exige un nom féminin singulier à l'accusatif.
\item \all{Arzt} : « beaucoup de patients vont chez le \emph{médecin} » ; \all{zum} = \all{zu dem} (Datif), donc \all{dem Arzt}.
\item \all{Heilpflanzen} : « d'autres utilisent des \emph{plantes médicinales} traditionnelles » ; accusatif pluriel sans article $\rightarrow$ \all{traditionelle Heilpflanzen}.
\item \all{Heilmittel} : « il n'existe pas de \emph{remède} simple » ; \all{kein einfaches} exige un nom neutre singulier à l'accusatif $\rightarrow$ \all{Heilmittel}.
\item \all{Impfung} : « la \emph{vaccination} protège la population » ; sujet nominatif féminin singulier $\rightarrow$ \all{die Impfung}.
\end{enumerate}

\textbf{Texte corrigé complet} :
\all{Malaria ist eine gefährliche Krankheit in Afrika. Viele Patienten gehen zum Arzt, aber andere benutzen traditionelle Heilpflanzen aus der Natur. Gegen Malaria gibt es kein einfaches Heilmittel. Die Impfung gegen andere Krankheiten schützt die Bevölkerung.}
\end{exoresolu}

\courssub{Grammaire 1 : les verbes pronominaux (Reflexivverben)}

\begin{propriete}[Formation, cas du pronom et place]
\begin{enumerate}
\item \textbf{Deux types} :
\begin{itemize}
\item \emph{Essentiels} : le pronom fait partie du sens du verbe, pas d'équivalent non pronominal (\all{sich erholen}, \all{sich beeilen}, \all{sich erkälten}).
\item \emph{Accidentels (ou "corps propre")} : l'action porte sur une partie du corps, le pronom est au \textbf{Datif} car le corps est COD (\all{sich die Hände waschen}, \all{sich die Zähne putzen}).
\end{itemize}
\item \textbf{Cas du pronom réfléchi} :
\begin{itemize}
\item \textbf{Akkusatif} si pas d'autre COD : \all{Ich wasche mich.} (Je me lave = je lave mon corps entier).
\item \textbf{Datif} s'il y a déjà un COD (partie du corps) : \all{Ich wasche mir die Hände.} (\all{die Hände} = COD Akk, \all{mir} = DAT).
\end{itemize}
\item \textbf{Place du pronom} :
\begin{itemize}
\item Proposition principale : \textbf{juste après le verbe conjugué} (\all{Er erholt sich schnell.} / \all{Er wäscht sich die Hände.}).
\item Subordonnée (\all{weil, dass, wenn...}) : \textbf{après le sujet} (\all{..., weil er sich schnell erholt.}).
\item Impératif : pronom après le verbe (\all{Wasch dir die Hände!}).
\end{itemize}
\item \textbf{Perfekt} : toujours avec \all{haben} + Participe II (\all{Ich habe mich erkältet.} / \all{Ich habe mir die Hände gewaschen.}). Jamais \all{sein}.
\end{enumerate}

\textbf{Tableau des pronoms réfléchis} :
\begin{center}
\begin{tabular}{|l|l|l|}
\hline
\rowcolor{popblueL} \textbf{Personne} & \textbf{Akkusativ} & \textbf{Dativ} \\
\hline
ich & mich & mir \\
\hline
du & dich & dir \\
\hline
er/sie/es & sich & sich \\
\hline
wir & uns & uns \\
\hline
ihr & euch & euch \\
\hline
sie/Sie & sich & sich \\
\hline
\end{tabular}
\end{center}
\end{propriete}

\begin{exemplebox}[Verbes pronominaux fréquents au programme (avec régime)]
\begin{itemize}
\item \all{sich interessieren für} + Akk (s'intéresser à)
\item \all{sich erinnern an} + Akk (se souvenir de)
\item \all{sich freuen über} + Akk (se réjouir de [présent/passé])
\item \all{sich freuen auf} + Akk (se réjouir de [futur])
\item \all{sich kümmern um} + Akk (s'occuper de)
\item \all{sich erholen von} + Dat (se remettre de)
\item \all{sich wenden an} + Akk (s'adresser à / se tourner vers)
\item \all{sich impfen lassen} (se faire vacciner — causatif)
\end{itemize}
\end{exemplebox}

\begin{exoresolu}[Thème : traduction avec verbes pronominaux]
\textbf{Énoncé} : Traduisez en allemand.
\begin{enumerate}
\item Pendant la crise sanitaire, beaucoup de Camerounais se sont fait vacciner.
\item Ma sœur s'est enrhumée, mais elle se remet vite.
\item Nous nous intéressons à la médecine traditionnelle.
\item Lave-toi les mains avant le repas !
\end{enumerate}
\tcblower
\textbf{Raisonnement et solutions} :
\begin{enumerate}
\item \all{Während der Gesundheitskrise haben sich viele Kameruner impfen lassen.}
      \textit{Analyse} : \all{während} + Génitif (\all{der Gesundheitskrise}) ; \all{sich impfen lassen} (causatif) au Perfekt $\rightarrow$ \all{haben} + pronom \all{sich} + \all{impfen lassen} (double infinitif). Verbe en 2\textsuperscript{e} position.
\item \all{Meine Schwester hat sich erkältet, aber sie erholt sich schnell.}
      \textit{Analyse} : \all{sich erkälten} (Perfekt \all{hat sich erkältet}) ; \all{sich erholen} (Présens \all{erholt sich}), ici sans \all{von} (absolu = "aller mieux").
\item \all{Wir interessieren uns für die traditionelle Medizin.}
      \textit{Analyse} : \all{sich interessieren für} + Akkusativ. Faux ami : ne jamais dire \all{interessieren} seul.
\item \all{Wasch dir die Hände vor dem Essen!}
      \textit{Analyse} : Impératif 2\textsuperscript{e} pers. sg \all{Wasch} ; COD \all{die Hände} $\rightarrow$ pronom réfléchi au \textbf{Datif} \all{dir}. \all{vor} + Datif \all{dem Essen}.
\end{enumerate}
\textbf{Conclusion} : trois réflexes — bon pronom (personne + cas), bonne place (après verbe conjugué), auxiliaire \all{haben} au Perfekt.
\end{exoresolu}

\courssub{Grammaire 2 : le conseil avec le subjonctif II (Konjunktiv II)}

\begin{propriete}[Formation et emploi du Konjunktiv II pour le conseil]
\textbf{1. Formation} :
\begin{itemize}
\item \textbf{Verbes forts} : base Präteritum + Umlaut (+ \all{-e} si possible) : \all{haben $\rightarrow$ hätte}, \all{sein $\rightarrow$ wäre}, \all{kommen $\rightarrow$ käme}, \all{gehen $\rightarrow$ ginge}, \all{wissen $\rightarrow$ wüsste}, \all{können $\rightarrow$ könnte}, \all{müssen $\rightarrow$ müsste}, \all{sollen $\rightarrow$ sollte}, \all{dürfen $\rightarrow$ dürfte}.
\item \textbf{Verbes faibles / mixtes / modaux (souvent)} : \all{würde} + \textbf{infinitif} à la fin (\all{ich würde gehen, ich würde machen}). C'est la forme standard pour presque tous les verbes sauf \all{haben, sein, modaux}.
\item \textbf{Modaux au Konjunktif II (formes irrégulières à connaître par cœur)} :
\end{itemize}
\begin{center}
\begin{tabular}{|l|l|l|}
\hline
\rowcolor{popblueL} \textbf{Infinitif} & \textbf{Konjunktiv II (ich/er/sie/es)} & \textbf{Sens conseil} \\
\hline
\all{sollen} & \all{sollte} & devrait (obligation morale/conseil) \\
\hline
\all{können} & \all{könnte} & pourrait (possibilité/suggestion) \\
\hline
\all{müssen} & \all{müsste} & devrait (nécessité) \\
\hline
\all{dürfen} & \all{dürfte} & pourrait (permission) \\
\hline
\all{wollen} & \all{wollte} & voudrait (sans Umlaut !) \\
\hline
\all{mögen} & \all{möchte} & aimerait (politesse) \\
\hline
\end{tabular}
\end{center}

\textbf{2. Structures types du conseil (à mémoriser)} :
\begin{itemize}
\item \all{Du solltest zum Arzt gehen.} (Conseil direct, verbe à la fin).
\item \all{An deiner Stelle würde ich mich impfen lassen.} (Conseil empathique, "à ta place").
\item \all{Es wäre besser, wenn du dich ausruhen würdest.} (\all{wenn} + subordonnée $\rightarrow$ verbe conjugué \all{würdest} \textbf{à la fin}).
\item \all{Ich würde dir raten, Heilpflanzen zu benutzen.} (\all{zu}-infinitif après virgule).
\end{itemize}

\textbf{3. Nuance} : Le Konjunktiv II rend le conseil \emph{politesse, hypothétique, non imposé}. L'impératif (\all{Geh zum Arzt!}) est direct/autoritaire.
\end{propriete}

\begin{exemplebox}[Attention : pièges classiques]
\begin{itemize}
\item \all{wollen} $\rightarrow$ \all{wollte} (PAS \all{wöllte}).
\item \all{An deiner Stelle} (Datif !) — jamais \all{*Auf deiner Stelle} (calque de "sur ta place").
\item Dans la subordonnée \all{wenn/dass} : le verbe conjugué (\all{würde, könnte, sollte} ou \all{würde} + infinitif) va \textbf{à la toute fin}.
\end{itemize}
\end{exemplebox}

\begin{exoresolu}[Expression guidée : donner des conseils à un ami]
\textbf{Énoncé} : Votre ami Paul est souvent malade et ne veut pas aller chez le médecin. Donnez-lui quatre conseils en utilisant : \all{sollen} (Konj. II), \all{würde} + infinitif, \all{Es wäre besser, wenn...}, \all{können} (Konj. II).
\tcblower
\textbf{Raisonnement} : chaque structure impose sa forme verbale et son ordre des mots.

\textbf{Production corrigée} :
\begin{enumerate}
\item \all{Du solltest endlich zum Arzt gehen, Paul.} (\all{solltest} 2\textsuperscript{e} pos, \all{gehen} fin).
\item \all{An deiner Stelle würde ich mich regelmäßig untersuchen lassen.} (\all{würde} 2\textsuperscript{e} pos, \all{mich} après sujet, \all{untersuchen lassen} fin).
\item \all{Es wäre besser, wenn du dich gesund ernähren würdest.} (\all{wenn} $\rightarrow$ \all{würdest} tout à la fin après \all{ernähren}).
\item \all{Du könntest auch traditionelle Heilpflanzen probieren, aber sprich zuerst mit einem Arzt.} (\all{könntest} 2\textsuperscript{e} pos, \all{probieren} fin ; \all{aber} coordination $\rightarrow$ \all{sprich} 2\textsuperscript{e} pos, \all{mit} + Dat \all{einem Arzt}).
\end{enumerate}
\textbf{Conclusion} : un bon conseil au Bac = Konjunktiv II correct + infinitif(s) en fin de phrase + vocabulaire du thème sanitaire.
\end{exoresolu}

\courssub{Méthode du commentaire de texte et commentaire modèle}

\begin{methode}[Comprendre un texte sur la santé (lecture globale puis détaillée)]
\begin{enumerate}
\item \textbf{Lire le titre et la source} : ils annoncent le thème (crise sanitaire, médecine traditionnelle, vaccination) et le point de vue de l'auteur.
\item \textbf{Première lecture rapide} (2 minutes) : repérer les \cle{mots transparents} (\all{die Pandemie}, \all{das Medikament}, \all{die Tradition}) et les chiffres. Ne pas s'arrêter sur les mots inconnus.
\item \textbf{Dégager l'idée générale} en une phrase : \all{Der Text handelt von...} (le texte traite de...).
\item \textbf{Deuxième lecture, paragraphe par paragraphe} : souligner les mots du champ lexical de la santé (\all{die Krankheit}, \all{der Arzt}, \all{die Heilpflanze}, \all{die Impfung}) et noter la structure : problème $\rightarrow$ causes $\rightarrow$ solutions.
\item \textbf{Répondre aux questions} en reformulant avec ses propres mots ; citer le texte entre guillemets allemands „…“ pour justifier.
\item \textbf{Vérifier} : la réponse doit être en allemand correct (verbe en 2e position, majuscules aux noms, Umlaute).
\end{enumerate}
\textbf{Réflexe} : une réponse de compréhension = phrase complète + justification tirée du texte.
\end{methode}

\begin{exoresolu}[Compréhension de texte : la médecine traditionnelle au Cameroun]
\textbf{Texte} (repris de la section 1) :
\all{In Kamerun vertrauen viele Menschen der traditionellen Medizin. Heiler benutzen Heilpflanzen gegen Krankheiten wie Malaria oder Magenschmerzen. Während der Gesundheitskrise haben sich viele Familien auch an die moderne Medizin gewandt: Sie haben sich impfen lassen und die Hygieneregeln beachtet. Experten sagen: Die beste Lösung ist die Zusammenarbeit zwischen traditioneller und moderner Medizin.}

\textbf{Questions} :
1. \all{Worum geht es in diesem Text?}
2. \all{Wogegen benutzen die Heiler Heilpflanzen?}
3. \all{Was haben die Familien während der Krise gemacht?}
4. \all{Was ist nach Meinung der Experten die beste Lösung?}
\tcblower
\textbf{Raisonnement} : on repère d'abord le thème (médecine traditionnelle et crise sanitaire au Cameroun), puis on localise chaque information dans le texte avant de reformuler.

\textbf{Réponses rédigées} :
\begin{enumerate}
\item \all{Der Text handelt von der traditionellen Medizin in Kamerun und von der Gesundheitskrise.} (Le texte traite de la médecine traditionnelle au Cameroun et de la crise sanitaire.) — Formule type : \all{Der Text handelt von...} + Datif.
\item \all{Die Heiler benutzen Heilpflanzen gegen Krankheiten wie Malaria oder Magenschmerzen.} — Reprise de \all{gegen} + Akkusativ.
\item \all{Während der Krise haben sich die Familien an die moderne Medizin gewandt: Sie haben sich impfen lassen und die Hygieneregeln beachtet.} — \all{sich wenden an} + Akk ; Perfekt \all{haben sich gewandt} / \all{haben sich impfen lassen}.
\item \all{Nach Meinung der Experten ist die beste Lösung die Zusammenarbeit zwischen traditioneller und moderner Medizin.} — \all{nach Meinung der Experten} (selon les experts) ; \all{zwischen} + Datif (\all{zwischen traditioneller und moderner Medizin}).
\end{enumerate}
\textbf{Conclusion} : chaque réponse est une phrase complète, verbe conjugué en 2e position, informations justifiées par le texte.
\end{exoresolu}

\courssub{Production guidée et réinvestissement}

\begin{methode}[Résumer un texte et donner son avis (Expression écrite)]
\begin{enumerate}
\item \textbf{Relever les idées essentielles} (qui ? quoi ? où ? pourquoi ?) et éliminer les exemples secondaires.
\item \textbf{Rédiger le résumé} en 3 à 5 phrases, avec ses propres mots, au présent ; début type : \all{Der Text informiert darüber, dass...} ou \all{Im Text geht es um...}
\item \textbf{Utiliser des connecteurs logiques} : \all{zuerst} (d'abord), \all{dann} (ensuite), \all{außerdem} (de plus), \all{deshalb} (c'est pourquoi), \all{schließlich} (finalement).
\item \textbf{Donner son avis} avec marqueurs explicites : \all{Meiner Meinung nach...}, \all{Ich bin der Ansicht, dass...} (verbe à la fin !), \all{Ich finde es wichtig, dass...}.
\item \textbf{Conclure} par une ouverture : \all{Ich hoffe, dass traditionelle und moderne Medizin in Zukunft zusammenarbeiten.}
\item \textbf{Relire} : verbe 2e position, majuscules noms, Umlaute, accords, guillemets „…“.
\end{enumerate}
\end{methode}

\begin{exoresolu}[Production écrite : résumé et opinion personnelle]
\textbf{Énoncé} : Résumez le texte suivant en trois phrases, puis donnez votre avis en deux phrases (environ 60 mots au total).

\all{In vielen Dörfern Kameruns spielen die traditionellen Heiler eine wichtige Rolle. Sie kennen die Heilpflanzen des Waldes und behandeln Krankheiten wie Malaria oder Hautprobleme. Während der letzten Gesundheitskrise haben aber viele Menschen gemerkt, dass die moderne Medizin mit ihren Impfungen unverzichtbar ist. Heute arbeiten in einigen Regionen Heiler und Ärzte zusammen.}
\tcblower
\textbf{Raisonnement} : idées essentielles = (1) rôle des guérisseurs traditionnels, (2) importance de la médecine moderne pendant la crise, (3) coopération actuelle. On reformule sans copier-coller.

\textbf{Résumé rédigé (3 phrases)} :
\all{Der Text handelt von der Medizin in Kamerun. Zuerst erklärt er die wichtige Rolle der traditionellen Heiler und ihrer Heilpflanzen. Dann zeigt er, dass die moderne Medizin während der Gesundheitskrise unverzichtbar war, und dass Heiler und Ärzte heute manchmal zusammenarbeiten.}

\textbf{Avis personnel (2 phrases)} :
\all{Meiner Meinung nach ist diese Zusammenarbeit sehr positiv, weil sie das Beste der beiden Welten verbindet. Ich finde es wichtig, dass der Staat die traditionelle Medizin kontrolliert, damit die Patienten sicher behandelt werden.}

\textbf{Justifications grammaticales} :
\begin{itemize}
\item \all{weil sie ... verbindet} (verbe à la fin).
\item \all{dass der Staat ... kontrolliert} (verbe à la fin).
\item \all{damit die Patienten sicher behandelt werden} (subordonnée finale \all{damit} + passif \all{werden behandelt}).
\end{itemize}
\textbf{Conclusion} : résumé objectif au présent + avis marqué par \all{meiner Meinung nach} = consigne respectée.
\end{exoresolu}

\begin{exercice}[Entraînement personnel : réinvestissement]
\begin{enumerate}
\item \textbf{Grammaire (Verbes pronominaux)} : Mettez le verbe entre parenthèses à la forme correcte (Présent ou Perfekt).
      \begin{enumerate}
      \item Gestern \_\_\_\_\_\_\_\_\_ ich \_\_\_\_\_\_\_\_\_ (sich erkälten).
      \item Meine Mutter \_\_\_\_\_\_\_\_\_ sich immer \_\_\_\_\_\_\_\_\_ (sich kümmern um) die Enkelkinder.
      \item Wir \_\_\_\_\_\_\_\_\_ uns \_\_\_\_\_\_\_\_\_ (sich freuen auf) die Ferien.
      \item \_\_\_\_\_\_\_\_\_ du \_\_\_\_\_\_\_\_\_ (sich waschen) die Hände vor dem Essen! (Impératif)
      \end{enumerate}
\item \textbf{Grammaire (Konjunktiv II - Conseil)} : Transformez les phrases à l'impératif en conseils polis (Konjunktiv II).
      \begin{enumerate}
      \item Geh zum Arzt! $\rightarrow$ \all{Du \_\_\_\_\_\_\_\_\_ zum Arzt gehen.} (\all{sollen})
      \item Nimm traditionelle Heilmittel! $\rightarrow$ \all{An deiner Stelle \_\_\_\_\_\_\_\_\_ ich traditionelle Heilmittel \_\_\_\_\_\_\_\_\_.} (\all{würde} + infinitif)
      \item Iss mehr Gemüse! $\rightarrow$ \all{Es wäre besser, wenn du mehr Gemüse \_\_\_\_\_\_\_\_\_.} (\all{essen})
      \end{enumerate}
\item \textbf{Vocabulaire} : Complétez les phrases avec le mot convenable (article + nom correct).
      \begin{enumerate}
      \item Der \_\_\_\_\_\_\_\_\_ (guérisseur) kennt viele \_\_\_\_\_\_\_\_\_ (plantes médicinales).
      \item Die \_\_\_\_\_\_\_\_\_ (vaccination) gegen Gelbfieber ist in Kamerun Pflicht.
      \item Während der \_\_\_\_\_\_\_\_\_ (crise sanitaire) haben wir \_\_\_\_\_\_\_\_\_ (règles d'hygiène) beachtet.
      \end{enumerate}
\item \textbf{Expression écrite} : Écrivez un court paragraphe (5-6 lignes) sur le thème : „Traditionelle und moderne Medizin in Kamerun“. Utilisez au moins 5 mots du vocabulaire de la leçon, un verbe pronominal et un conseil au Konjunktiv II.
\end{enumerate}
\end{exercice}


\begin{corrige}[Exercice : Entraînement personnel]
\textbf{1. Verbes pronominaux} :
\begin{enumerate}
\item \all{hat sich erkältet} (Perfekt, Akk).
\item \all{kümmert ... um} (Présens, \all{sich kümmern um} + Akk).
\item \all{freuen ... auf} (Présens, \all{sich freuen auf} + Akk).
\item \all{Wasch ... dir} (Impératif, Datif car COD \all{die Hände} sous-entendu ou explicite : \all{Wasch dir die Hände!}).
\end{enumerate}

\textbf{2. Konjunktiv II - Conseil} :
\begin{enumerate}
\item \all{solltest} (\all{Du solltest zum Arzt gehen.}).
\item \all{würde ... nehmen} (\all{An deiner Stelle würde ich traditionelle Heilmittel nehmen.}).
\item \all{isst} / \all{essen würdest} (Deux solutions : \all{Es wäre besser, wenn du mehr Gemüse isst} [indicatif réel possible mais moins poli] OU \all{Es wäre besser, wenn du mehr Gemüse essen würdest} [Konj II standard]). La forme attendue est \all{essen würdest}.
\end{enumerate}

\textbf{3. Vocabulaire} :
\begin{enumerate}
\item \all{Heiler} (masc, sing), \all{Heilpflanzen} (fém, plur).
\item \all{Impfung} (fém, sing, sujet/acc).
\item \all{Gesundheitskrise} (fém, sing, après \all{während} + Génitif $\rightarrow$ \all{der Gesundheitskrise}), \all{Hygieneregeln} (fém, plur, Akk).
\end{enumerate}

\textbf{4. Expression écrite (proposition de corrigé)} :
\all{In Kamerun gibt es sowohl traditionelle Heiler als auch moderne Ärzte. Die Heiler benutzen Heilpflanzen aus dem Wald gegen Krankheiten wie Malaria. Während der Gesundheitskrise haben sich viele Menschen impfen lassen. Meiner Meinung nach ist die Zusammenarbeit sehr wichtig. Du solltest immer einen Arzt fragen, bevor du Heilpflanzen nimmst.}
\end{corrige}

\begin{attention}[Les 5 erreurs qui coûtent des points au Bac]
\begin{enumerate}
\item \textbf{Oublier le pronom réfléchi ou le mettre au mauvais cas} : on écrit \all{Ich habe mich erkältet} (pas \all{*ich habe erkältet}) ; avec un COD (partie du corps), le pronom passe au \textbf{Datif} : \all{Ich wasche mir die Hände} (pas \all{*mich die Hände}).
\item \textbf{Mauvaise place du verbe en subordonnée} : avec \all{weil, dass, wenn, damit}, le verbe conjugué va \textbf{à la fin} : \all{..., weil er sich nicht impfen lassen wollte.} Jamais \all{*weil er wollte sich nicht impfen lassen}.
\item \textbf{Calquer le français pour les conseils} : « tu devrais » = \all{du solltest} (Konj. II de \all{sollen}), PAS futur \all{*wirst sollen}. « À ta place » = \all{an deiner Stelle} (Dativ), PAS \all{*auf deiner Stelle}.
\item \textbf{Faux amis et mauvaises constructions prépositionnelles} :
      \begin{itemize}
      \item « s'intéresser à » = \all{sich interessieren für} + \textbf{Akkusativ}.
      \item « souffrir de » = \all{leiden an} + \textbf{Dativ}.
      \item « se remettre de » = \all{sich erholen von} + \textbf{Dativ}.
      \item « se tourner vers » = \all{sich wenden an} + \textbf{Akkusativ}.
      \end{itemize}
      Vérifier toujours la préposition et le cas dans le dictionnaire.
\item \textbf{Négliger l'orthographe allemande} : \textbf{majuscule à tous les noms} (\all{die Krankheit}, \all{das Heilmittel}, \all{die Impfung}), \textbf{Umlaute obligatoires} (\all{Ärzte}, \all{würde}, \all{könnte}, \all{müsste}, \all{Hautprobleme}), \textbf{ß après voyelle longue/diphtongue} (\all{gesundheit} $\rightarrow$ \all{Gesundheit}, \all{Maßnahme}), et \textbf{guillemets allemands „…“} pour citer le texte.
\end{enumerate}
\end{attention}

\newpage

\coursec{Exercices}

\courssub{Je vérifie mes connaissances}

\begin{exercice}[QCM : vocabulaire de la santé]
Entoure la bonne réponse.
\begin{enumerate}
\item \all{die Impfung} signifie :
\begin{enumerate}
\item la maladie \item la vaccination \item l'opération
\end{enumerate}
\item \all{die Heilpflanze} signifie :
\begin{enumerate}
\item la plante médicinale \item l'hôpital \item la pharmacie
\end{enumerate}
\item \all{sich erkälten} signifie :
\begin{enumerate}
\item se reposer \item s'enrhumer \item se guérir
\end{enumerate}
\item \all{das Heilmittel} signifie :
\begin{enumerate}
\item le patient \item le remède \item le médecin
\end{enumerate}
\item \all{die Ansteckung} signifie :
\begin{enumerate}
\item la contagion \item la consultation \item la guérison
\end{enumerate}
\end{enumerate}
\end{exercice}

\begin{exercice}[Vrai ou faux ? Justifie ta réponse]
Réponds par \all{richtig} ou \all{falsch} et justifie en une phrase en français.
\begin{enumerate}
\item \all{Die Pandemie} est une maladie limitée à un seul village.
\item \all{Der Patient} est la personne qui soigne les malades.
\item Le verbe \all{sich beeilen} se construit avec l'auxiliaire \all{haben} au Perfekt.
\item \all{Du solltest zum Arzt gehen} est une phrase qui exprime un conseil.
\item \all{Die traditionelle Medizin} utilise souvent des \all{Heilpflanzen}.
\end{enumerate}
\end{exercice}

\begin{exercice}[Lexique : articles et pluriels]
Complète le tableau avec l'article et le pluriel de chaque nom.
\begin{center}
\begin{tabular}{|l|l|l|}
\hline
\rowcolor{popblueL} Nom & Article & Pluriel \\
\hline
Krankheit & \ldots & \ldots \\
\hline
Arzt & \ldots & \ldots \\
\hline
Heilmittel & \ldots & \ldots \\
\hline
Impfung & \ldots & \ldots \\
\hline
Pandemie & \ldots & \ldots \\
\hline
Patient & \ldots & \ldots \\
\hline
\end{tabular}
\end{center}
\end{exercice}

\begin{exercice}[Conjugaison à compléter]
Complète avec la forme correcte du verbe pronominal au Präsens.
\begin{enumerate}
\item Ich \ldots\ (sich erkälten) oft im Winter.
\item \ldots\ du \ldots\ (sich fühlen) heute besser ?
\item Er \ldots\ (sich interessieren) für die traditionelle Medizin.
\item Wir \ldots\ (sich beeilen), denn der Arzt wartet.
\item Ihr \ldots\ (sich erinnern) an die Impfung im letzten Jahr.
\end{enumerate}
\end{exercice}

\courssub{Je m'entraîne}

\begin{exercice}[Le bon pronom réfléchi]
Complète les huit phrases avec le pronom réfléchi correct : \all{mich}, \all{mir}, \all{dich}, \all{dir}, \all{sich}, \all{uns}, \all{euch}.
\begin{enumerate}
\item Ich habe \ldots\ beim Fußballspielen verletzt.
\item Fühlst du \ldots\ heute besser, mein Kind ?
\item Die Patientin erinnert \ldots\ an ihre erste Impfung.
\item Du solltest \ldots\ mehr ausruhen, du bist müde.
\item Der Junge hat \ldots\ die Hand gebrochen.
\item Wir interessieren \ldots\ für die Heilpflanzen unserer Großeltern.
\item Erkältest du \ldots\ oft im Regen ?
\item Die Schüler beeilen \ldots\, um rechtzeitig im Krankenhaus anzukommen.
\end{enumerate}
\end{exercice}

\begin{exercice}[Appariements : mots et traductions]
Relie chaque mot allemand à sa traduction française.
\begin{center}
\begin{tabular}{|l|l|}
\hline
\rowcolor{popblueL} Deutsch & Français \\
\hline
1. die Impfung & a. le patient \\
\hline
2. das Heilmittel & b. la contagion \\
\hline
3. die Heilpflanze & c. la vaccination \\
\hline
4. der Patient & d. le remède \\
\hline
5. die Ansteckung & e. la plante médicinale \\
\hline
\end{tabular}
\end{center}
Puis complète : \all{\ldots\ Impfung (Pl. \ldots)}, \all{\ldots\ Heilmittel (Pl. \ldots)}, \all{\ldots\ Heilpflanze (Pl. \ldots)}.
\end{exercice}

\begin{textebox}[In der Klinik]
In Kamerun gehen viele \ldots\ zuerst zum traditionellen Heiler, bevor sie einen \ldots\ besuchen. Gegen manche \ldots\ benutzt der Heiler \ldots\ aus dem Wald. Er bereitet daraus ein \ldots\ zu. Aber gegen gefährliche Krankheiten wie Polio hilft nur die \ldots\ im Gesundheitszentrum.
\end{textebox}

\begin{exercice}[Texte à trous]
Complète le texte « In der Klinik » ci-dessus avec les mots de la liste : \all{Heilpflanzen -- Impfung -- Arzt -- Krankheit -- Heilmittel -- Patienten}. Attention : un seul mot de la liste est au pluriel, et un mot devra être mis au pluriel.
\end{exercice}

\begin{exercice}[Transformation : du Präsens au Perfekt]
Mets les phrases suivantes au Perfekt. Attention au choix de l'auxiliaire et à la place du pronom réfléchi.
\begin{enumerate}
\item Ich erkälte mich im Regen.
\item Fühlst du dich besser ?
\item Er interessiert sich für die Schulmedizin.
\item Wir beeilen uns am Morgen.
\item Die Großmutter erinnert sich an den alten Heiler des Dorfes.
\end{enumerate}
\end{exercice}

\courssub{Je me prépare au Bac}

\begin{textebox}[Impfen in Deutschland]
In Deutschland diskutiert man oft über die Impfung. Viele Eltern fragen sich : Sollen wir unsere Kinder impfen lassen ? Die meisten Ärzte sagen : „Ja, unbedingt !“ Denn die Impfung schützt nicht nur das eigene Kind, sondern die ganze Gesellschaft. Wenn sich viele Menschen impfen lassen, kann sich eine gefährliche Krankheit nicht so schnell ausbreiten. Das hat man während der letzten Pandemie deutlich gesehen.

Trotzdem gibt es Menschen, die sich vor der Impfung fürchten. Sie haben Angst vor Nebenwirkungen oder sie vertrauen der Schulmedizin nicht. Manche erinnern sich an die Heilpflanzen ihrer Großeltern und glauben, dass die Natur alles heilen kann. Die Ärzte antworten : „Die traditionelle Medizin kann bei kleinen Beschwerden helfen, aber gegen schwere Krankheiten braucht man die moderne Medizin.“

Der Staat informiert die Bürger mit Plakaten, im Radio und im Internet. Auch in den Schulen sprechen die Lehrer mit den Schülern über Gesundheit und Vorsorge. Denn wer sich gut informiert, kann sich besser entscheiden.
\end{textebox}

\begin{exercice}[Compréhension de texte type Bac]
Lis le texte « Impfen in Deutschland » ci-dessus, puis réponds aux questions en allemand, par des phrases complètes.

\textbf{Fragen zum Text :}
\begin{enumerate}
\item Welche Frage stellen sich viele Eltern in Deutschland ?
\item Warum empfehlen die meisten Ärzte die Impfung ? Nenne zwei Gründe.
\item Warum fürchten sich manche Menschen vor der Impfung ?
\item Was antworten die Ärzte auf das Argument der Heilpflanzen ?
\item Wie informiert der Staat die Bürger ? Nenne drei Möglichkeiten.
\end{enumerate}

\textbf{Fragen zur Sprache :}
\begin{enumerate}
\item Suche im Text zwei Verben pronominaux und konjugiere sie im Präsens an allen Personen.
\item Suche im Text eine Beratungsstruktur (Modalverb oder Konjunktiv II) und erkläre ihre Funktion.
\end{enumerate}
\end{exercice}

\begin{exercice}[Thème et version]
\textbf{A. Thème.} Traduis en allemand les cinq phrases suivantes. Utilise le Konjunktiv II ou un verbe modal pour exprimer le conseil.
\begin{enumerate}
\item Tu devrais te reposer, tu es très fatigué.
\item Il vaudrait mieux que tu ailles chez le médecin.
\item On pourrait utiliser des plantes médicinales contre cette maladie.
\item Vous devriez vous faire vacciner avant le voyage.
\item Si j'étais toi, je me dépêcherais d'aller à l'hôpital.
\end{enumerate}

\textbf{B. Version.} Traduis en français le texte « Der Heiler aus dem Dorf » ci-dessous.
\end{exercice}

\begin{textebox}[Der Heiler aus dem Dorf]
In unserem Dorf bei Bafoussam lebt ein alter Heiler. Jeden Morgen beeilt er sich, um seine Heilpflanzen im Wald zu sammeln. Die Kranken fühlen sich bei ihm wohl, denn er hört ihnen lange zu. Er erinnert sich noch an seinen Großvater, der ihm alles über die Pflanzen erklärt hat. Viele Patienten interessieren sich heute wieder für diese traditionelle Medizin.
\end{textebox}

\begin{exercice}[Production écrite type Bac]
\textbf{Sujet :} \all{Die traditionelle Medizin in meinem Land}

Rédige un texte de 10 à 12 lignes en allemand. Tu présenteras le rôle de la médecine traditionnelle au Cameroun, ses avantages et ses limites, et tu donneras ton avis personnel.

\textbf{Contraintes obligatoires :}
\begin{itemize}
\item utilise au moins \textbf{trois verbes pronominaux} différents (par exemple \all{sich erinnern}, \all{sich fühlen}, \all{sich interessieren}, \all{sich erkälten}) ;
\item formule au moins \textbf{deux conseils au Konjunktiv II} (par exemple avec \all{sollte}, \all{könnte}, \all{würde}) ;
\item emploie au moins cinq mots du vocabulaire du chapitre (\all{die Krankheit}, \all{das Heilmittel}, \all{die Heilpflanze}, \all{der Patient}, \all{die Impfung}, \ldots) ;
\item soigne l'ordre des mots et la place du verbe.
\end{itemize}
\end{exercice}

\begin{methode}[Réussir la production écrite]
\begin{enumerate}
\item \textbf{Analyse le sujet} : trois parties sont attendues (le rôle, les avantages et limites, ton avis). Fais un plan rapide au brouillon.
\item \textbf{Liste ton vocabulaire} avant d'écrire : \all{der Heiler}, \all{die Heilpflanze}, \all{das Heilmittel}, \all{die Krankheit}, \all{der Patient}, \all{die Impfung}, \all{das Krankenhaus}.
\item \textbf{Place correctement les verbes pronominaux} : verbe conjugué en 2\up{e} position, pronom réfléchi juste après : \all{Viele Patienten fühlen sich beim Heiler wohl.}
\item \textbf{Formule tes conseils au Konjunktiv II} : \all{Man sollte die Heilpflanzen kontrollieren.} / \all{Der Staat könnte die Heiler ausbilden.}
\item \textbf{Relis-toi} : majuscules des noms, Umlaute, terminaisons des verbes, place du verbe dans les subordonnées (\all{weil}, \all{dass}, \all{wenn} : verbe à la fin).
\end{enumerate}
\end{methode}

\begin{exemplebox}[Modèle de production écrite]
\all{In Kamerun spielt die traditionelle Medizin eine wichtige Rolle. Viele Patienten gehen zuerst zum Heiler, weil er in der Nähe wohnt und nicht viel Geld kostet. Die Kranken fühlen sich bei ihm wohl, denn er hört ihnen lange zu. Die Heiler benutzen Heilpflanzen aus dem Wald und bereiten daraus Heilmittel zu. Meine Großmutter erinnert sich noch an die Pflanzen ihrer Kindheit.}

\all{Aber die traditionelle Medizin hat auch Grenzen : Gegen schwere Krankheiten wie Polio hilft nur die Impfung im Gesundheitszentrum. Manche Heilmittel sind auch gefährlich, weil niemand die Dosis kontrolliert.}

\all{Meiner Meinung nach sollte man die beiden Medizinen verbinden. Der Staat könnte die Heiler ausbilden und ihre Heilpflanzen untersuchen. Wenn ich krank wäre, würde ich zuerst zum Arzt gehen und dann die traditionelle Medizin als Ergänzung benutzen.}

\emph{Traduction :} Au Cameroun, la médecine traditionnelle joue un rôle important. Beaucoup de patients vont d'abord chez le guérisseur, parce qu'il habite près de chez eux et ne coûte pas cher. Les malades se sentent bien chez lui, car il les écoute longuement. Les guérisseurs utilisent des plantes médicinales de la forêt et en préparent des remèdes. Ma grand-mère se souvient encore des plantes de son enfance. Mais la médecine traditionnelle a aussi des limites : contre les maladies graves comme la polio, seule la vaccination au centre de santé est efficace. Certains remèdes sont aussi dangereux, parce que personne ne contrôle la dose. À mon avis, on devrait combiner les deux médecines. L'État pourrait former les guérisseurs et examiner leurs plantes médicinales. Si j'étais malade, j'irais d'abord chez le médecin, puis j'utiliserais la médecine traditionnelle en complément.
\end{exemplebox}

\begin{attention}[Erreurs fréquentes à éviter]
\begin{itemize}
\item Ne dis pas \all{Ich habe mich erkältet worden} : les verbes pronominaux se conjuguent simplement avec \all{haben} + participe : \all{Ich habe mich erkältet.}
\item N'oublie pas le pronom réfléchi : on écrit \all{Er interessiert sich für Medizin}, jamais \all{Er interessiert für Medizin}.
\item Attention au cas du pronom : \all{sich die Hand brechen} exige le datif (\all{Ich habe mir die Hand gebrochen}), car la partie du corps porte déjà l'accusatif.
\item Au Konjunktiv II, n'oublie pas l'Umlaut : \all{könnte}, \all{müsste}, \all{würde}, \all{sollte} (sans Umlaut pour \all{sollen} et \all{wollen} : \all{sollte}, \all{wollte}).
\item Faux ami : \all{die Medizin} signifie « la médecine » (la discipline), le remède se dit \all{das Medikament} ou \all{das Heilmittel}.
\end{itemize}
\end{attention}

\newpage

\coursec{Corrigés des exercices}

\begin{corrige}[Exercice 1 — QCM : vocabulaire de la santé]
\begin{enumerate}
\item b -- \all{die Impfung} = la vaccination. (Famille : \all{impfen} = vacciner ; \all{sich impfen lassen} = se faire vacciner.)
\item a -- \all{die Heilpflanze} = la plante médicinale. (Composé de \all{heilen} « guérir » + \all{die Pflanze} « la plante ».)
\item b -- \all{sich erkälten} = s'enrhumer. (Nom associé : \all{die Erkältung} = le rhume.)
\item b -- \all{das Heilmittel} = le remède. (\all{das Mittel} = le moyen ; \all{heil-} = guérison.)
\item a -- \all{die Ansteckung} = la contagion. (Verbe : \all{sich anstecken} = se contaminer ; adjectif : \all{ansteckend} = contagieux.)
\end{enumerate}
\end{corrige}

\begin{corrige}[Exercice 2 — Vrai ou faux ?]
\begin{enumerate}
\item \all{Falsch} : une pandémie (\all{die Pandemie}) s'étend à plusieurs pays, voire au monde entier ; une maladie limitée à une région est une épidémie (\all{die Epidemie}).
\item \all{Falsch} : \all{der Patient} est la personne malade ; la personne qui soigne est le médecin (\all{der Arzt}, féminin \all{die Ärztin}).
\item \all{Richtig} : \all{sich beeilen} est un verbe pronominal, et tous les verbes pronominaux forment le Perfekt avec \all{haben} : \all{Ich habe mich beeilt.} (« Je me suis dépêché. »)
\item \all{Richtig} : \all{solltest} est le Konjunktiv II de \all{sollen} ; il exprime le conseil atténué : « tu devrais ».
\item \all{Richtig} : la médecine traditionnelle (\all{die traditionelle Medizin}) utilise fréquemment des plantes médicinales (\all{die Heilpflanzen}).
\end{enumerate}
\end{corrige}

\begin{corrige}[Exercice 3 — Lexique : articles et pluriels]
\begin{center}
\begin{tabular}{|l|l|l|}
\hline
\rowcolor{popblueL} Nom & Article & Pluriel \\
\hline
Krankheit & die & die Krankheiten \\
\hline
Arzt & der & die Ärzte \\
\hline
Heilmittel & das & die Heilmittel \\
\hline
Impfung & die & die Impfungen \\
\hline
Pandemie & die & die Pandemien \\
\hline
Patient & der & die Patienten \\
\hline
\end{tabular}
\end{center}
Points de vigilance :
\begin{itemize}
\item \all{der Arzt} prend l'Umlaut au pluriel : \all{die Ärzte}.
\item Les noms en \all{-ung} et \all{-heit} sont toujours féminins et prennent \all{-en} au pluriel : \all{die Impfung -- die Impfungen}, \all{die Krankheit -- die Krankheiten}.
\item \all{der Patient} est un nom faible : il prend \all{-en} à tous les cas sauf au nominatif singulier : \all{Ich besuche den Patienten.} (« Je rends visite au patient. »)
\end{itemize}
\end{corrige}

\begin{corrige}[Exercice 4 — Conjugaison à compléter]
\begin{enumerate}
\item Ich \all{erkälte mich} oft im Winter. (« Je m'enrhume souvent en hiver. »)
\item \all{Fühlst} du \all{dich} heute besser ? (« Te sens-tu mieux aujourd'hui ? »)
\item Er \all{interessiert sich} für die traditionelle Medizin. (« Il s'intéresse à la médecine traditionnelle. »)
\item Wir \all{beeilen uns}, denn der Arzt wartet. (« Nous nous dépêchons, car le médecin attend. »)
\item Ihr \all{erinnert euch} an die Impfung im letzten Jahr. (« Vous vous souvenez de la vaccination de l'année dernière. »)
\end{enumerate}
Rappel grammatical : le pronom réfléchi se place immédiatement après le verbe conjugué ; à la 2\up{e} personne du pluriel, on utilise \all{euch} (accusatif et datif identiques). Attention à la construction : \all{sich interessieren für} + accusatif ; \all{sich erinnern an} + accusatif.
\end{corrige}

\begin{corrige}[Exercice 5 — Le bon pronom réfléchi]
\begin{enumerate}
\item Ich habe \all{mich} beim Fußballspielen verletzt. (accusatif : \all{sich verletzen} = se blesser)
\item Fühlst du \all{dich} heute besser, mein Kind ? (accusatif : \all{sich fühlen} = se sentir)
\item Die Patientin erinnert \all{sich} an ihre erste Impfung. (accusatif : \all{sich erinnern an} + A)
\item Du solltest \all{dich} mehr ausruhen, du bist müde. (accusatif : \all{sich ausruhen} = se reposer)
\item Der Junge hat \all{sich} die Hand gebrochen. (datif : la partie du corps, \all{die Hand}, porte déjà l'accusatif ; le pronom passe donc au datif : \all{sich (D) die Hand (A) brechen})
\item Wir interessieren \all{uns} für die Heilpflanzen unserer Großeltern. (accusatif : \all{sich interessieren für} + A)
\item Erkältest du \all{dich} oft im Regen ? (accusatif : \all{sich erkälten})
\item Die Schüler beeilen \all{sich}, um rechtzeitig im Krankenhaus anzukommen. (accusatif : \all{sich beeilen} ; à la 3\up{e} personne du pluriel, le pronom est \all{sich})
\end{enumerate}
Point de vigilance : c'est la phrase 5 qui piège le plus souvent les francophones — quand le verbe pronominal a déjà un complément d'objet direct (une partie du corps), le pronom réfléchi est au datif : \all{Ich habe mir die Hand gebrochen.}
\end{corrige}

\begin{corrige}[Exercice 6 — Appariements : mots et traductions]
1 -- c (\all{die Impfung} = la vaccination) ; 2 -- d (\all{das Heilmittel} = le remède) ; 3 -- e (\all{die Heilpflanze} = la plante médicinale) ; 4 -- a (\all{der Patient} = le patient) ; 5 -- b (\all{die Ansteckung} = la contagion).

Articles et pluriels : \all{die Impfung (Pl. die Impfungen)}, \all{das Heilmittel (Pl. die Heilmittel)}, \all{die Heilpflanze (Pl. die Heilpflanzen)}.

Remarque : les noms en \all{-ung} et en \all{-e} (féminins) prennent \all{-en} ou \all{-n} au pluriel ; \all{das Heilmittel} reste identique au pluriel, comme beaucoup de noms neutres en \all{-el}.
\end{corrige}

\begin{corrige}[Exercice 7 — Texte à trous]
In Kamerun gehen viele \all{Patienten} zuerst zum traditionellen Heiler, bevor sie einen \all{Arzt} besuchen. Gegen manche \all{Krankheit} benutzt der Heiler \all{Heilpflanzen} aus dem Wald. Er bereitet daraus ein \all{Heilmittel} zu. Aber gegen gefährliche Krankheiten wie Polio hilft nur die \all{Impfung} im Gesundheitszentrum.

Traduction : « Au Cameroun, beaucoup de patients vont d'abord chez le guérisseur traditionnel avant de consulter un médecin. Contre certaines maladies, le guérisseur utilise des plantes médicinales de la forêt. Il en prépare un remède. Mais contre les maladies dangereuses comme la polio, seule la vaccination au centre de santé est efficace. »

Justifications :
\begin{itemize}
\item \all{viele} + pluriel impose \all{Patienten}.
\item \all{einen} + accusatif masculin impose \all{Arzt} (singulier).
\item \all{manche} + singulier impose \all{Krankheit} au singulier.
\item \all{Heilpflanzen} est le seul mot déjà au pluriel dans la liste : le guérisseur utilise plusieurs plantes.
\item \all{ein} + neutre impose \all{Heilmittel} ; \all{die} + féminin impose \all{Impfung}.
\end{itemize}
\end{corrige}

\begin{corrige}[Exercice 8 — Transformation : du Präsens au Perfekt]
\begin{enumerate}
\item Ich \all{habe mich} im Regen \all{erkältet}. (« Je me suis enrhumé sous la pluie. »)
\item \all{Hast} du \all{dich} besser \all{gefühlt} ? (« T'es-tu senti mieux ? »)
\item Er \all{hat sich} für die Schulmedizin \all{interessiert}. (« Il s'est intéressé à la médecine moderne. »)
\item Wir \all{haben uns} am Morgen \all{beeilt}. (« Nous nous sommes dépêchés le matin. »)
\item Die Großmutter \all{hat sich} an den alten Heiler des Dorfes \all{erinnert}. (« La grand-mère s'est souvenue du vieux guérisseur du village. »)
\end{enumerate}
Règle appliquée : tous les verbes pronominaux forment le Perfekt avec l'auxiliaire \all{haben} + participe passé, sans exception ; le pronom réfléchi se place juste après le verbe conjugué. Participes : \all{erkältet} (verbe faible en \all{-t-} : \all{ge-...-et}), \all{gefühlt}, \all{interessiert} (verbe en \all{-ieren} : pas de \all{ge-}), \all{beeilt}, \all{erinnert}.
\end{corrige}

\begin{corrige}[Exercice 9 — Compréhension de texte type Bac]
\textbf{Fragen zum Text :}
\begin{enumerate}
\item Viele Eltern fragen sich, ob sie ihre Kinder impfen lassen sollen. (« Beaucoup de parents se demandent s'ils doivent faire vacciner leurs enfants. »)
\item Die Ärzte empfehlen die Impfung, weil sie das eigene Kind schützt und weil sich eine gefährliche Krankheit nicht so schnell ausbreiten kann, wenn sich viele Menschen impfen lassen.
\item Manche Menschen fürchten sich vor der Impfung, weil sie Angst vor Nebenwirkungen haben oder weil sie der Schulmedizin nicht vertrauen.
\item Die Ärzte antworten, dass die traditionelle Medizin bei kleinen Beschwerden helfen kann, dass man aber gegen schwere Krankheiten die moderne Medizin braucht.
\item Der Staat informiert die Bürger mit Plakaten, im Radio und im Internet ; auch die Lehrer sprechen in den Schulen über Gesundheit und Vorsorge.
\end{enumerate}

\textbf{Fragen zur Sprache :}
\begin{enumerate}
\item Verbes pronominaux du texte (au choix) : \all{sich fragen}, \all{sich fürchten}, \all{sich erinnern}, \all{sich entscheiden}, \all{sich ausbreiten}, \all{sich impfen lassen}.

Conjugaison de \all{sich fürchten} au Präsens : ich fürchte mich, du fürchtest dich, er/sie/es fürchtet sich, wir fürchten uns, ihr fürchtet euch, sie/Sie fürchten sich.

Conjugaison de \all{sich erinnern} au Präsens : ich erinnere mich, du erinnerst dich, er/sie/es erinnert sich, wir erinnern uns, ihr erinnert euch, sie/Sie erinnern sich.
\item Structure de conseil : \all{Sollen wir unsere Kinder impfen lassen ?} — le modal \all{sollen} sert ici à demander un avis. Au Konjunktiv II, \all{sollten} exprime le conseil atténué : \all{Wir sollten unsere Kinder impfen lassen.} (« Nous devrions faire vacciner nos enfants. »)
\end{enumerate}

Barème indicatif (sur 10) : questions de compréhension 5 × 1 pt (contenu 0,5 + correction linguistique 0,5) ; questions de langue : 2 verbes pronominaux correctement conjugués 2 pts ; structure de conseil identifiée et expliquée 2 pts ; qualité globale de l'expression 1 pt.

Points de vigilance : répondre par des phrases complètes avec le verbe en 2\up{e} position ; dans les subordonnées en \all{weil} et \all{dass}, le verbe conjugué va à la fin ; ne pas oublier le pronom réfléchi (\all{sich fürchten vor} + datif).
\end{corrige}

\begin{corrige}[Exercice 10 — Thème et version]
\textbf{A. Thème.}
\begin{enumerate}
\item \all{Du solltest dich ausruhen, du bist sehr müde.} (Konjunktiv II de \all{sollen} ; verbe pronominal \all{sich ausruhen}.)
\item \all{Es wäre besser, wenn du zum Arzt gehen würdest.} (Variante acceptée : \all{Du solltest lieber zum Arzt gehen.})
\item \all{Man könnte Heilpflanzen gegen diese Krankheit benutzen.} (Konjunktiv II de \all{können} avec Umlaut.)
\item \all{Sie sollten sich vor der Reise impfen lassen.} (\all{sich impfen lassen} = se faire vacciner ; \all{vor} + datif.)
\item \all{Wenn ich du wäre, würde ich mich beeilen, ins Krankenhaus zu gehen.} (Variante : \all{An deiner Stelle würde ich mich beeilen, ins Krankenhaus zu gehen.})
\end{enumerate}

\textbf{B. Version.}

Dans notre village près de Bafoussam vit un vieux guérisseur. Chaque matin, il se dépêche pour aller cueillir ses plantes médicinales dans la forêt. Les malades se sentent bien chez lui, car il les écoute longuement. Il se souvient encore de son grand-père, qui lui a tout expliqué sur les plantes. Aujourd'hui, beaucoup de patients s'intéressent de nouveau à cette médecine traditionnelle.

Barème indicatif (sur 10) : thème 5 × 1 pt (Konjunktiv II correct 0,5 + pronom réfléchi et ordre des mots 0,5) ; version 5 pts (fidélité du sens 3 pts, qualité du français 2 pts).

Points de vigilance : ne pas oublier l'Umlaut au Konjunktiv II (\all{könnte}, \all{würde}, mais \all{sollte} sans Umlaut) ; \all{sich beeilen, ... zu + Infinitiv} : le pronom reste avec le verbe conjugué ; en version, rendre \all{jemandem zuhören} par « écouter quelqu'un » et \all{sich erinnern an} par « se souvenir de ».
\end{corrige}

\begin{corrige}[Exercice 11 — Production écrite type Bac]
\textbf{Proposition de rédaction modèle :}

\all{In Kamerun spielt die traditionelle Medizin eine wichtige Rolle. Viele Patienten gehen zuerst zum Heiler, weil er in der Nähe wohnt und nicht viel Geld kostet. Die Kranken fühlen sich bei ihm wohl, denn er hört ihnen lange zu. Die Heiler benutzen Heilpflanzen aus dem Wald und bereiten daraus Heilmittel zu. Meine Großmutter erinnert sich noch an die Pflanzen ihrer Kindheit.}

\all{Aber die traditionelle Medizin hat auch Grenzen : Gegen schwere Krankheiten wie Polio hilft nur die Impfung im Gesundheitszentrum. Manche Heilmittel sind auch gefährlich, weil niemand die Dosis kontrolliert.}

\all{Meiner Meinung nach sollte man die beiden Medizinen verbinden. Der Staat könnte die Heiler ausbilden und ihre Heilpflanzen untersuchen. Wenn ich krank wäre, würde ich zuerst zum Arzt gehen und dann die traditionelle Medizin als Ergänzung benutzen.}

\textbf{Vérification des contraintes :}
\begin{itemize}
\item Trois verbes pronominaux différents : \all{sich fühlen}, \all{sich erinnern}, (+ \all{sich erkälten} possible) — présents.
\item Deux conseils au Konjunktiv II : \all{man sollte verbinden}, \all{der Staat könnte ausbilden}, \all{ich würde gehen} — présents.
\item Cinq mots du chapitre : \all{die Medizin}, \all{der Patient}, \all{die Heilpflanze}, \all{das Heilmittel}, \all{die Impfung}, \all{die Krankheit} — présents.
\item Ordre des mots : verbe en 2\up{e} position dans les principales, verbe à la fin après \all{weil} et \all{wenn}.
\end{itemize}

Barème indicatif (sur 20) : respect du sujet et du plan en trois parties 4 pts ; contraintes grammaticales (3 verbes pronominaux + 2 Konjunktiv II) 6 pts ; vocabulaire du chapitre 4 pts ; correction linguistique (ordre des mots, cas, orthographe, majuscules) 4 pts ; cohérence et avis personnel argumenté 2 pts.

Points de vigilance : majuscule à tous les noms ; ne pas écrire \all{Er interessiert für} sans \all{sich} ; \all{Meiner Meinung nach} est suivi du verbe en 2\up{e} position ; éviter le calque du français « la médecine est bonne » — préférer \all{die traditionelle Medizin kann helfen}.
\end{corrige}

\newpage

\coursec{Fiche bilan}

\begin{popfigure}
\begin{center}
\begin{tikzpicture}[mindmap, grow cyclic, every node/.style={concept, font=\small}, concept color=popblue!60, text=popdark, level 1/.append style={level distance=3.4cm, sibling angle=60, font=\footnotesize}, level 2/.append style={level distance=2.2cm, font=\scriptsize}]
\node[concept color=popblue, text=white, font=\bfseries\small, align=center] {Les maladies\\ Krankheiten}
  child[concept color=poporange!70] { node[align=center] {Lexique santé\\ \all{die Krankheit}}
    child { node[align=center] {\all{der Arzt}\\ \all{die Impfung}} }
    child { node[align=center] {\all{das Heilmittel}\\ \all{die Heilpflanze}} }
  }
  child[concept color=popgreen!70] { node[align=center] {Médecine\\ traditionnelle}
    child { node[align=center] {\all{der Heiler}\\ \all{die Erfahrung}} }
  }
  child[concept color=poppink!70] { node[align=center] {Crise sanitaire\\ \all{die Pandemie}}
    child { node[align=center] {\all{sich schützen}\\ \all{die Hygiene}} }
  }
  child[concept color=popgold!80] { node[align=center] {Verbes\\ pronominaux}
    child { node[align=center] {\all{sich erkälten}\\ \all{sich fühlen}} }
    child { node[align=center] {acc. / dat.\\ place du pronom} }
  }
  child[concept color=poppurple!70] { node[align=center] {Konjunktiv II\\ le conseil}
    child { node[align=center] {\all{du solltest}\\ \all{du könntest}} }
    child { node[align=center] {\all{an deiner Stelle}} }
  }
  child[concept color=popblue!50] { node[align=center] {Méthode\\ commentaire}
    child { node[align=center] {lire -- repérer\\ analyser -- rédiger} }
  };
\end{tikzpicture}
\end{center}
\legende{Carte mentale de la leçon : maladies, crise sanitaire et médecine traditionnelle}
\end{popfigure}

\begin{bilan}
\textbf{1. Lexique clé à connaître par cœur}
\begin{itemize}
  \item \all{die Krankheit, -en} : la maladie ; \all{die Pandemie, -n} : la pandémie ; \all{die Epidemie, -n} : l'épidémie
  \item \all{der Arzt, Ärzte} / \all{die Ärztin, -nen} : le médecin / la médecin ; \all{der Patient, -en} : le patient
  \item \all{das Heilmittel, -} : le remède ; \all{die Heilpflanze, -n} : la plante médicinale ; \all{der Heiler, -} : le guérisseur
  \item \all{die Impfung, -en} : la vaccination ; \all{impfen} : vacciner ; \all{das Medikament, -e} : le médicament
  \item \all{die Hygiene} : l'hygiène ; \all{die Behandlung, -en} : le traitement ; \all{das Fieber} : la fièvre
  \item \all{die Gesundheit} : la santé ; \all{gesund / krank} : en bonne santé / malade ; \all{sich erholen} : se rétablir
\end{itemize}

\textbf{2. Grammaire à maîtriser}
\begin{center}
\begin{tabularx}{\linewidth}{|X|X|}
\hline
\rowcolor{popblueL} \textbf{Règle} & \textbf{Exemple} \\
\hline
Verbe pronominal : pronom réfléchi à l'accusatif (\all{mich, dich, sich, uns, euch, sich}), placé juste après le verbe conjugué & \all{Ich erkälte mich oft.} « Je m'enrhume souvent. » \\
\hline
Quelques verbes pronominaux au datif : \all{sich (Dat.) etwas waschen, sich etwas kaufen} & \all{Ich wasche mir die Hände.} « Je me lave les mains. » \\
\hline
Verbes pronominaux du thème : \all{sich fühlen, sich erkälten, sich erholen, sich schützen, sich impfen lassen, sich Sorgen machen} & \all{Er lässt sich impfen.} « Il se fait vacciner. » \\
\hline
Conseil au Konjunktiv II : \all{sollen $\to$ solltest}, \all{können $\to$ könntest}, \all{müssen $\to$ müsstest} + infinitif en fin de phrase & \all{Du solltest mehr Sport treiben.} « Tu devrais faire plus de sport. » \\
\hline
Formule de conseil : \all{An deiner Stelle würde ich ...} / \all{Es wäre besser, wenn ...} & \all{An deiner Stelle würde ich zum Arzt gehen.} « À ta place, j'irais chez le médecin. » \\
\hline
\end{tabularx}
\end{center}

\textbf{3. Méthodes express}
\begin{itemize}
  \item \textbf{Compréhension de texte} : 1) lire le titre et la source, 2) première lecture globale (qui ? quoi ? où ?), 3) lecture détaillée avec soulignage des mots-clés, 4) répondre en phrases complètes en allemand.
  \item \textbf{Commentaire} : introduction (présenter le texte et le thème), développement (idées principales + vocabulaire du champ lexical de la santé), conclusion (avis personnel bref).
  \item \textbf{Conseil oral} : identifier le problème $\to$ choisir le modal adapté $\to$ Konjunktiv II + infinitif final.
\end{itemize}
\end{bilan}

\begin{aretenir}[Checklist avant le Bac]
\begin{itemize}
  \item $\square$ Je connais l'article et le pluriel des noms clés : \all{die Krankheit, -en ; der Arzt, Ärzte ; das Heilmittel, - ; die Impfung, -en}.
  \item $\square$ Je sais conjuguer un verbe pronominal au présent : \all{ich fühle mich, du fühlst dich, er fühlt sich ...}.
  \item $\square$ Je place correctement le pronom réfléchi juste après le verbe conjugué.
  \item $\square$ Je distingue pronom réfléchi à l'accusatif (\all{mich / sich}) et au datif (\all{mir die Hände waschen}).
  \item $\square$ Je connais les formes de Konjunktiv II des modaux : \all{solltest, könntest, müsstest, dürftest, wolltest}.
  \item $\square$ Je sais donner un conseil : \all{Du solltest zum Arzt gehen.} et \all{An deiner Stelle würde ich ...}.
  \item $\square$ Je n'oublie pas l'infinitif à la fin de la phrase avec un modal.
  \item $\square$ Je mets une majuscule à tous les noms allemands (\all{die Pandemie, das Fieber}).
  \item $\square$ Je sais parler de la médecine traditionnelle : \all{die Heilpflanze, der Heiler, die traditionelle Medizin}.
  \item $\square$ Je sais résumer un texte en 3 ou 4 phrases avec mes propres mots.
  \item $\square$ Je relis ma production : place du verbe en position 2, accords, orthographe.
\end{itemize}
\end{aretenir}

\findecours

\end{document}
